\ifx\danslelivre\undefined
\RequirePackage{luatex85}
\documentclass[10pt]{../fichiers-configuration-latex/smfart}
\usepackage{polyglossia}
%\usepackage[francais]{babel}
\usepackage{fontspec}
\usepackage[math-style=ISO,bold-style=upright,sans-style=upright]{unicode-math}
%\usepackage[mathletters]{ucs}
%\usepackage[utf8]{inputenc}
%\usepackage[T1]{fontenc}
\usepackage[unicode,colorlinks=true,linkcolor=blue,citecolor=Sepia,urlcolor=magenta]{hyperref}
\setmainlanguage{french}
\setmainfont[Ligatures=TeX]{Linux Libertine O}
% variante possible : TeXGyrePagella ?
\setmathfont{XITS Math}

% maths
\usepackage{amsmath}
%\usepackage{amssymb}
%\usepackage{mathrsfs}
%\usepackage{stmaryrd}
%\usepackage{euscript}

% couleurs
\usepackage[usenames,dvipsnames]{xcolor}

% graphiques
\usepackage{graphicx}
\usepackage[all]{xy}
% Xy-pic 3.8.9's PDF driver assumes that \fontdimen8\textfont3 is a
% legacy TeX rule width.  Under LuaTeX with unicode-math/XITS Math it is
% instead 6.62pt, which turns every arrow shaft into a thick black bar.
% Override only xypdf's private cached width; do not alter the math font.
\makeatletter
\AtBeginDocument{%
  \@ifundefined{xP@preclw}{}{%
    \edef\xP@preclw{0.4pt}%
    \edef\xP@lw{\xP@dim\xP@preclw}%
    \PackageInfo{sga4-xypdf-compat}{Xy-pic PDF line width set to \xP@preclw}%
  }%
}
\makeatother
%\usepackage{tikz}
%\usetikzlibrary{matrix,arrows,calc}
%\tikzset{column sep=3em,row sep=2em,text height=1.5ex, text depth=0.25ex,>=angle 90}
%\tikzset{description/.style={fill=white,inner sep=1pt}}

% hyperliens, 'etiquettes
\usepackage{xr-hyper}
\usepackage{../fichiers-configuration-latex/smfthm}
%\ifx\cacheretiquettes\undefined
%\usepackage[inline]{showlabels} \else \fi % pour rendre visibles les 'etiquettes
%\usepackage{../fichiers-configuration-latex/makeidx} % pour faire un index
%\usepackage[pagebackref,unicode]{hyperref}  % pour voir o`u sont cit'ees les r'ef'erences
\hypersetup{unicode,colorlinks=true,linkcolor=blue,citecolor=Sepia,urlcolor=magenta,
  hypertexnames=false}
%\hypersetup{breaklinks=true} %permet le retour `a la ligne dans les liens trop longs. Incompatible avec pdflatex
%\renewcommand*{\backref}[1]{{\small [cit\'e page(s)~#1.]}}

\usepackage{xspace}
\usepackage{pifont} % pour expose XVII

\renewcommand{\thefootnote}{\roman{footnote}} % pour mettre les notes en bas de page sous la forme (ii) et non (2), qui ressemble trop (?) a un exposant.

\newcounter{pageoriginal}
\newcommand{\pageoriginale}{\refstepcounter{pageoriginal}\marginpar{\footnotesize\textbf{\thepageoriginal}}\xspace}
\let\pageoriginaled\pageoriginale

\usepackage{ifthen}
\newboolean{original}
\def\sisi{\ifthenelse{\boolean{original}}}
\setboolean{original}{false}

\newcounter{cnote}
\def\nde#1{\ifthenelse{\boolean{original}}{}{\stepcounter{cnote}\def\thefootnote{\arabic{cnote}}\protect\footnote[\value{cnote}]{N.D.E.: #1}}}
\def\ndemark{\ifthenelse{\boolean{original}}{}{\stepcounter{cnote}\def\thefootnote{\arabic{cnote}}\protect\footnotemark[\value{cnote}]}}
\def\ndetext#1{\ifthenelse{\boolean{original}}{}{\def\thefootnote{\arabic{cnote}}\protect\footnotetext[\value{cnote}]{N.D.E.: #1}}}

\setmainlanguage{french}
\setotherlanguage{english}

% unicode-math deliberately disables the legacy \Arrowvert glyph.  The
% source uses it once as a rotated equality marker; \Vert has the intended
% appearance and keeps the historical diagram buildable with LuaLaTeX.
\renewcommand{\Arrowvert}{\Vert}


%%
%% Pour afficher le numéro de version Git (exige que LuaTeX soit appelé avec --shell-escape)
%%

\catcode`\%=12
\directlua{
function command_output(cmd)
  local output_pipe = io.popen(cmd)
  local line_iter = output_pipe and output_pipe:lines()
  local line = line_iter and line_iter()
  if output_pipe then output_pipe:close() end
  return line
end
local git_command1 = [[git log --pretty=format:"%h %ai" -1]]
local git_command2 = [[if git diff --quiet HEAD ; then echo clean ; else echo dirty ; fi]]
function git_status()
  local status = command_output(git_command1)
  if status then
    tex.tprint({luatexbase.catcodetables.CatcodeTableLaTeX, string.char(0x5c).."texttt{"}, {-2,status}, {luatexbase.catcodetables.CatcodeTableLaTeX, "}"})
    if not (command_output(git_command2) == "clean") then tex.sprint(-2, " (modifications non validées)") end
  else
    tex.sprint(-2, "?")
  end
end
}
\catcode`\%=14
\def\showgitstatus{\directlua{git_status()}}


% macros pour LaTeX
\renewcommand{\mathsf}{}

% Police différente pour U+2113 SCRIPT SMALL L
\setmathfont[range="2113]{latinmodern-math.otf}

% Give every explicit SGA label its own destination and preserve the printed
% number supplied as the optional argument.  The former definition inherited
% whichever counter anchor happened to precede it and ignored #1 entirely.
\makeatletter
\def\etiquette[#1]#2{%
  \phantomsection
  \protected@edef\@currentlabel{#1}%
  \label{#2}%
}
\makeatother
\def\eqhyperref[#1]#2{(\hyperref[#1]{#2})}
\def\pagehyperref[#1]#2{(page \hyperref[#1]{#2})}
\def\hyperrefext[#1][#2]#3{\hyperref[#2]{#1 #3}}
\def\eqhyperrefext[#1][#2]#3{\hyperref[#2]{#1 #3}}     %cf. exposes 17
\def\neweqhyperrefext[#1][#2]#3{\hyperref[#2]{#1 #3}}  %cf. exposes 17

\def\shat#1{\protect\widehat{#1}}
\def\stilde#1{\protect\widetilde{#1}}
%%%%%%%%%%%%%%%%%%


\def\etcn{etc.\xspace}
\def\etco{etc...\xspace}
\def\etc{\sisi{\etco}{\etcn}}
\def\ldotsn{\sisi{...}{\etcn}}

\def\cf{{cf.}\kern.3em}
\def\Cf{{Cf.}\kern.3em}
\def\ie{{i.e.}\kern.3em}
\def\iev{{i.e.},\xspace}
\def\loccito{loc.\kern2pt cit.\xspace}
\def\loccitn{\textit{loc.\kern2pt cit.}\xspace}
\def\loccit{\sisi{\loccito}{\loccitn}}
\def\EGA{\textup{EGA}\kern.2em}
\def\SGA{\textup{SGA}\kern.3em}

\def\quoi{}

\def\quav{\quad,\quad}
\def\quaet{\quad \text{et}\quad}
\def\fibreproduct#1#2#3{#1 ×_{#3} #2}
\let\fprod\fibreproduct
\def\ccheck#1#2{\mbox{${\check{\check{\rm #1}}}^{\lower2pt\hbox{$\scriptstyle #2$}}$}}
\def\sheaf#1{\protect\mathcal{#1}}
\def\vieuxsheaf#1{{#1}} % cf. problème avec sheaf{a\Phi_i} par exemple...
\DeclareMathOperator{\Hom}{Hom}
\DeclareMathOperator{\can}{can}
\DeclareMathOperator{\Ind}{Ind}
\DeclareMathOperator{\ind}{ind}
%\DeclareMathOperator{\Hhomtop}{\mathcal{H}\!\mathit{omtop}}
\DeclareMathOperator{\Hhomtop}{\mathbf{Homtop}}
\DeclareMathOperator{\Point}{Point}
\DeclareMathOperator{\Pro}{Pro}
\DeclareMathOperator{\Top}{Top}
\DeclareMathOperator{\Mod}{Mod}
\DeclareMathOperator{\Cart}{Cart}
%\DeclareMathOperator{\Mmorsite}{\mathcal{M}\!\mathit{orsite}}
\DeclareMathOperator{\Mmorsite}{\mathbf{Morsite}}
\DeclareMathOperator{\loc}{loc}
\DeclareMathOperator{\Loc}{Loc}
\DeclareMathOperator{\Pic}{Pic}
\DeclareMathOperator{\Et}{Et}
\DeclareMathOperator{\card}{card}
\DeclareMathOperator{\Ens}{Ens}
\DeclareMathOperator{\h}{\underline{H}}
\DeclareMathOperator{\Fonct}{Fonct}
\DeclareMathOperator{\Ker}{Ker}
\DeclareMathOperator{\but}{but}
\DeclareMathOperator{\Ob}{Ob}
\DeclareMathOperator{\ob}{ob}
\DeclareMathOperator{\Lib}{Lib}
\DeclareMathOperator{\esj}{esj}
\DeclareMathOperator{\id}{id}
\DeclareMathOperator{\Id}{Id}
\DeclareMathOperator{\ad}{ad}
\DeclareMathOperator{\Iim}{Im}
\DeclareMathOperator{\pr}{pr}
\DeclareMathOperator{\Cat}{Cat}
\DeclareMathOperator{\car}{car}
\DeclareMathOperator{\SysInd}{SysInd}
\DeclareMathOperator{\Sup}{Sup}
\DeclareMathOperator{\Coker}{Coker}
\DeclareMathOperator{\Filt}{Filt}
\DeclareMathOperator{\acc}{acc}
\DeclareMathOperator{\sur}{sur}
\DeclareMathOperator{\cat}{cat}
\DeclareMathOperator{\Ord}{Ord}
\DeclareMathOperator{\Gr}{Gr}
\DeclareMathOperator{\Cov}{Cov}
\DeclareMathOperator{\Caf}{Caf}
\DeclareMathOperator{\coker}{coker}
\DeclareMathOperator{\coim}{coim}
\DeclareMathOperator{\im}{im}
\DeclareMathOperator{\Ouv}{Ouv}
\DeclareMathOperator{\diag}{diag}
\DeclareMathOperator{\Esp}{Esp}
\DeclareMathOperator{\TOP}{TOP}
\DeclareMathOperator{\Unif}{Unif}
\DeclareMathOperator{\Iind}{\mathbf{Ind}}
\DeclareMathOperator{\Hhomcart}{\mathbf{Homcart}}
\DeclareMathOperator{\Hhomtopess}{\mathbf{Homtopess}}
\DeclareMathOperator{\Ppoints}{\mathbf{Points}}
\DeclareMathOperator{\Ppoint}{\mathbf{Point}}
\DeclareMathOperator{\Ttors}{\mathbf{Tors}}
\DeclareMathOperator{\Ssex}{\mathbf{Sex}}
\DeclareMathOperator{\Ffib}{\mathbf{Fib}}
\DeclareMathOperator{\Vvois}{\mathbf{Vois}}
\DeclareMathOperator{\Ppointess}{\mathbf{Pointess}}
\DeclareMathOperator{\scrHom}{\mathbf{Hom}}
\DeclareMathOperator{\sheafPic}{\mathbf{Pic}}
\DeclareMathOperator{\sheafAut}{\mathbf{Aut}}
\DeclareMathOperator{\Hhomtoplocan}{\mathbf{Homtoplocan}}
\DeclareMathOperator{\Hhomtopan}{\mathbf{Homtopan}}
\DeclareMathOperator{\Pptgeom}{\mathbf{Ptgéom}}
\DeclareMathOperator{\Morsite}{Morsite}
\DeclareMathOperator{\Fib}{Fib}
\DeclareMathOperator{\Fais}{Fais}
\DeclareMathOperator{\kar}{kar}
\DeclareMathOperator{\Groupes}{Groupes}
\DeclareMathOperator{\Restr}{Restr}
\DeclareMathOperator{\Isom}{Isom}
\DeclareMathOperator{\Prol}{Prol}
\DeclareMathOperator{\Res}{Res}
\DeclareMathOperator{\Tors}{Tors}
\DeclareMathOperator{\esp}{esp}
\DeclareMathOperator{\es}{es}
\DeclareMathOperator{\od}{od}
\DeclareMathOperator{\gr}{gr}
\DeclareMathOperator{\sob}{sob}
\DeclareMathOperator{\dfn}{dfn}
\DeclareMathOperator{\Sch}{Sch}
\DeclareMathOperator{\ouv}{ouv}
\DeclareMathOperator{\cons}{cons}
\DeclareMathOperator{\Nerf}{Nerf}
\DeclareMathOperator{\Points}{Points}
\DeclareMathOperator{\Espsob}{Espsob}
\DeclareMathOperator{\Homtop}{Homtop}
\DeclareMathOperator{\Pt}{Pt}
\DeclareMathOperator{\Quot}{Quot}
\DeclareMathOperator{\glob}{glob}
\DeclareMathOperator{\Homfib}{Homfib}
\DeclareMathOperator{\Qouv}{Qouv}
\DeclareMathOperator{\Quotouv}{Quotouv}
\DeclareMathOperator{\ens}{ens}
\DeclareMathOperator{\ab}{ab}
\DeclareMathOperator{\an}{an}
\DeclareMathOperator{\fs}{fs}
\DeclareMathOperator{\red}{red}
\DeclareMathOperator{\Ab}{Ab}
\DeclareMathOperator{\cosq}{cosq}
\DeclareMathOperator{\Ann}{Ann}
\DeclareMathOperator{\Zar}{Zar}


%%%%%%%%%%%%
%%%%% SGA4-2
\DeclareMathOperator{\Sex}{Sex}
\DeclareMathOperator{\vois}{vois}
\DeclareMathOperator{\Vois}{Vois}
\DeclareMathOperator{\Ext}{Ext}
\DeclareMathOperator{\ext}{ext}
\DeclareMathOperator{\Fer}{Fer}
\DeclareMathOperator{\Add}{Add}
\DeclareMathOperator{\Cocart}{Cocart}
\DeclareMathOperator{\cosk}{cosk}
\DeclareMathOperator{\plat}{plat}

\DeclareMathOperator{\faisc}{faisc}
\DeclareMathOperator{\sk}{sk}
\DeclareMathOperator{\ch}{ch}

\def\HH{\mathbf{H}}
\def\qc{\mathrm{qc}}
\def\et{\mathrm{et}}
\def\coh{\mathrm{coh}}

\def\Hvr#1{%
\begin{array}[b]{l}
{\kern2pt\scriptstyle r}\\[-8pt]
{\kern1pt\scriptstyle \vee}\\[-5pt]
H^{#1}
\end{array}}

\def\Hvv#1{%
\begin{array}[b]{l}
\kern2pt\rotatebox{90}{${\scriptscriptstyle\ll}$}\\[-6pt]
H^{#1}
\end{array}}

\def\Hvi#1{%
\begin{array}[b]{l}
{\kern0.2pt\scriptscriptstyle \infty}\\[-8pt]
{\kern1pt\scriptstyle \vee}\\[-5pt]
H^{#1}
\end{array}}

\def\Hv#1{%
\begin{array}[b]{l}
{\kern1pt\scriptstyle \vee}\\[-5pt]
H^{#1}
\end{array}}

\DeclareMathOperator{\cart}{cart}
\DeclareMathOperator{\Ffaisc}{\mathbf{Faisc}}


\def\Ev#1{%
\begin{array}[b]{l}
{\kern1pt\scriptstyle \vee}\\[-5pt]
E'_{#1}
\end{array}}

\DeclareMathOperator{\Cont}{Cont}
\DeclareMathOperator{\Ccont}{\mathscr{C}\!\mathit{ont}}
\DeclareMathOperator{\Site}{Site}
\newcommand\MM{\mathbf{M}}
\def\limtop#1{{\displaystyle{\varprojlim_{#1}\!\mbox{\rm top }}}}
\def\Limtop{\mathop{\varprojlim\!\mbox{\rm top}}}
\DeclareMathOperator{\spec}{spec}
\def\ss{\scriptstyle{\mathscr{S}}}
\DeclareMathOperator{\Gal}{Gal}
\def\smallm{\scriptstyle{\mathscr{M}}}
\DeclareMathOperator{\End}{End}
\DeclareMathOperator{\ass}{ass}
\DeclareMathOperator{\zar}{zar}
\DeclareMathOperator{\Ogg}{Ogg}
\DeclareMathOperator{\res}{res}
\DeclareMathOperator{\tr}{tr}
\DeclareMathOperator{\Supp}{Supp}
\DeclareMathOperator{\codim}{codim}
\DeclareMathOperator{\rk}{rk}
\DeclareMathOperator{\cur}{cur}
\DeclareMathOperator{\cdqc}{cdqc}
\newcommand\DD{\mathbf{D}}
\DeclareMathOperator{\cl}{cl}
\def\JJ{\mathbf{J}}
\def\FF{\mathbf{F}}
\def\BB{\mathbf{B}}
\DeclareMathOperator{\Alget}{Alget}
\DeclareMathOperator{\Ass}{Ass}
\DeclareMathOperator{\supp}{supp}
\DeclareMathOperator{\rad}{rad}
\DeclareMathOperator{\prof}{prof}
\DeclareMathOperator{\cint}{int}
\DeclareMathOperator{\sic}{sic}
\DeclareMathOperator{\Tor}{Tor}
\DeclareMathOperator{\tors}{tors}
\DeclareMathOperator{\Sym}{Sym}
\DeclareMathOperator{\parf}{parf}
\DeclareMathOperator{\tor}{tor}
\DeclareMathOperator{\defn}{defn}
\DeclareMathOperator{\Div}{Div}
\DeclareMathOperator{\Pol}{Pol}
\DeclareMathOperator{\Fex}{Fex}
\DeclareMathOperator{\Lie}{Lie}
\DeclareMathOperator{\tot}{tot}
\DeclareMathOperator{\Aut}{Aut}
\DeclareMathOperator{\HOM}{HOM}
\DeclareMathOperator{\pch}{pch}
\DeclareMathOperator{\Ch}{Ch}
\DeclareMathOperator{\EXT}{EXT}
\DeclareMathOperator{\PIC}{PIC}
\DeclareMathOperator{\Rf}{Rf}
\DeclareMathOperator{\inj}{inj}
\newcommand\VV{\mathbf{V}}

\def\Xv{%
\begin{array}[b]{l}
{\kern1pt\scriptstyle \vee}\\[-5pt]
X
\end{array}}

%%%%%%%
\newcommand\TT{\mathbf{T}}
\newcommand\XX{\mathbf{X}}
\newcommand\EE{\mathbf{E}}
%%%%%%

\def\qedsymbol{\textup{C.Q.F.D.}}
\let\CQFD\qed
\let\cqfd\qed

%%%%%\bigdot
\font\bigsymb=cmsy10 at 4pt
\def\bigdot{{\kern1.2pt\raise 1.5pt\hbox{\bigsymb\char15}}}
%%%%%

\def\Lright{\mathop{\vcenter{\hbox{\hbox{$\longrightarrow$}\lower8pt \hbox{\kern-16pt\hbox{$\longrightarrow$}}}}}}

\def\Aa{{\mathcal A}}
\def\Bb{{\mathcal B}}
\def\Cc{{\mathcal C}}
\def\Dd{{\mathcal D}}
\def\Ee{{\mathcal E}}
\def\Ff{{\mathcal F}}
\def\Gg{{\mathcal G}}
\def\Hh{{\mathcal H}}
\def\Ii{{\mathcal I}}
\def\Jj{{\mathcal J}}
\def\Kk{{\mathcal K}}
\def\Ll{{\mathcal L}}
\def\Mm{{\mathcal M}}
\def\Nn{{\mathcal N}}
\def\Oo{{\mathcal O}}
\def\Pp{{\mathcal P}}
\def\Qq{{\mathcal Q}}
\def\Rr{{\mathcal R}}
\def\Ss{{\mathcal S}}
\def\Tt{{\mathcal T}}
\def\Uu{{\mathcal U}}
\def\Vv{{\mathcal V}}
\def\Ww{{\mathcal W}}
\def\Xx{{\mathcal X}}
\def\Yy{{\mathcal Y}}
\def\Zz{{\mathcal Z}}

\newcommand{\CC}{\mathbf{C}}
\newcommand{\GG}{\mathbf{G}}
\newcommand{\KK}{\mathbf{K}}
\newcommand{\LL}{\mathbf{L}}
\newcommand{\NN}{\mathbf{N}}
\newcommand{\PP}{\mathbf{P}}
\newcommand\QQ{\mathbf{Q}}
\newcommand\RR{\mathbf{R}}
\newcommand\ZZ{\mathbf{Z}}

\def\to{\mathchoice{\longrightarrow}{\rightarrow}{\rightarrow}{\rightarrow}}
\def\from{\mathchoice{\longleftarrow}{\leftarrow}{\leftarrow}{\leftarrow}}
\def\mto{\mathchoice{\longmapsto}{\mapsto}{\mapsto}{\mapsto}}
\def\mfrom{\mathrel{\from\joinrel\mapstochar\,}}
\def\hto{\mathrel{\lhook\joinrel\to}}
\def\hfrom{\mathrel{\from\joinrel\rhook}}

\def\implies{\mathchoice{\Longrightarrow}{\Rightarrow}{\Rightarrow}{\Rightarrow}}
\def\implied{\mathchoice{\Longleftarrow}{\Leftarrow}{\Leftarrow}{\Leftarrow}}
\def\iff{\mathchoice{\Longleftrightarrow}{\Leftrightarrow}{\Leftrightarrow}{\Leftrightarrow}}

\def\To#1{\mathchoice{\xrightarrow{\labelstyle~~#1~}}{\stackrel{#1}{\longrightarrow}}{}{}}
\def\Mto#1{\mathrel{\mapstochar\To{#1}}}
\def\From#1{\mathchoice{\xleftarrow{\labelstyle~~#1~}}{\stackrel{#1}{\longleftarrow}}{}{}}

\newcommand{\isomto}{\overset{\sim}{\longrightarrow}}
\newcommand{\isomfrom}{\overset{\sim}{\longleftarrow}}

\newcommand{\IFF}{\unskip~$\iff$~\ignorespaces}
\newcommand{\THEN}{\unskip~$\implies$~\ignorespaces}

\let\leq\leqslant
\let\geq\geqslant
\let\le\leqslant
\let\ge\geqslant

\def\cf{{cf.}\kern.3em}
\def\Cf{{Cf.}\kern.3em}
\def\ie{{i.e.}\kern.3em}
\def\iev{{i.e.},\xspace}
%\def\loccito{loc.\kern2pt cit.\xspace}
%\def\loccitn{\textit{loc.\kern2pt cit.}\xspace}
%\def\loccit{\sisi{\loccito}{\loccitn}}
\def\EGA{\textup{EGA}\kern.2em}
\def\SGA{\textup{SGA}\kern.3em}

\let\oldS\S
\def\S{\oldS\kern .15em}
\def\ptbl{.\kern .2em}
\def\resp{resp.\kern.3em}
\def\Exp{Exp.\kern.3em}

\let\bar\overline
\let\hat\widehat
\let\tilde\widetilde
\let\dpl\displaystyle
\let\ldots\dots
\let\emptyset\varnothing
\let\skpt\skippointrait
\let\epsilon\varepsilon

\def\sfrac#1#2{{#1}/{#2}}

\DeclareMathOperator{\cd}{cd}
%\DeclareMathOperator{\Hhom}{\mathcal{H}\!\mathit{om}}
\DeclareMathOperator{\Hhom}{\mathbf{Hom}}
%\DeclareMathOperator{\Eext}{\mathcal{E}\!\mathit{xt}}
\DeclareMathOperator{\Eext}{\mathbf{Ext}}
\DeclareMathOperator{\Spec}{Spec}
\DeclareMathOperator{\Tr}{Tr}

\let\emptyset=\varnothing

%%%%%%%%%%%%%%%
%%% expose XVII
%%%%%%%%%%%%%%%

\makeatletter
\let\old@optexpose\optexpose
\newcommand{\empty@optexpose}[1]{}
\newcommand{\paren@optexpose}[1]{\old@optexpose{#1}(}
\newcommand{\refnoexpose}[1]{\let\optexpose\empty@optexpose\ref{#1}\let\optexpose\old@optexpose}%\refnoexpose{X.cor5.2} engendre "5.2" au lieu de "X 5.2"
\newcommand{\neweqref}[1]{\let\optexpose\paren@optexpose\textup{\ref{#1})}\let\optexpose\old@optexpose}%\neweqref{XII.eq4.1} engendre "XII (4.1)" au lieu de "(XII 4.1)"
\makeatother
\newcommand{\refvb}[2][VI${}_{\text{B}}$ ]{\sisi{#1\refnoexpose{#2}}{\ref{#2}}}

\newcommand{\leftexp}[2]{{\vphantom{#2}}^{#1}{#2}}
\newcommand{\AAA}{\mathbf{A}}
\newcommand{\Sf}{\mathfrak{S}}
\DeclareMathOperator*{\quotprojlim}{\text{``$\varprojlim$''}}
\DeclareMathOperator*{\quotinjlim}{\text{``$\varinjlim$''}}
\DeclareMathOperator{\Fl}{F\sisi{\ell}{l}}
\DeclareMathOperator{\Pl}{P\sisi{\ell}{l}}
%\DeclareMathOperator{\cl}{c\sisi{\ell}{l}}
\DeclareMathOperator{\tdf}{tdf}
\DeclareMathOperator{\rg}{rg}
\newcommand{\iInd}{\sisi{I}{i}nd}
\newcommand{\iind}{ind}
\newcommand{\ppro}{pro}
\newcommand{\tTor}{\sisi{T}{t}or}
\newcommand{\ttor}{tor}
\newcommand{\chap}{Chap.}
\newcommand{\Ffex}{\Ff\mathit{ex}}
\newcommand{\sheafk}{\underline{k}}
\newcommand{\sheafn}{\underline{n}}
\newcommand{\sheafGamma}{\underline{\Gamma}}

\sisi{
\DeclareMathOperator{\simarrow}{\rule{.15mm}{1.8mm}\kern -3.4pt\rightsquigarrow}
\newcommand{\sheafell}{\underline{\ell}}
\newcommand{\sheafZZ}{\underline{\ZZ}}
}{
\newcommand{\nderef}[1]{la N.D.E.~(\ref{#1}) page~\pageref{#1}}
}

%%%%%%%%%%%%%%%%%
%%%%% expose XVII
%%%%%%%%%%%%%%%%%



\externaldocument{../01/I.tex}
\externaldocument{../02/II.tex}
\externaldocument{../03/III.tex}
\externaldocument{../04/IV.tex}
%\externaldocument{../05/V.tex}
\externaldocument{../05bis/Vbis.tex}
\externaldocument{../06/VI.tex}
\externaldocument{../07/VII.tex}
\externaldocument{../08/VIII.tex}
\externaldocument{../09/IX.tex}
\externaldocument{../10/X.tex}
\externaldocument{../11/XI.tex}
\externaldocument{../12/XII.tex}
\externaldocument{../13/XIII.tex}
\externaldocument{../14/XIV.tex}
\externaldocument{../15/XV.tex}
\externaldocument{../16/XVI.tex}
\externaldocument{../17/XVII.tex}
\externaldocument{../18/XVIII.tex}
\externaldocument{../19/XIX.tex}

\setboolean{original}{false}

%% À faire :
%% ajouter étiquettes
%% uniformiser i) versus (i) dans énoncés et démonstrations
%% uniformiser X., X_∙, etc.

\begin{document}

\begin{center}
V. Cohomologie dans les topos\\
J.-L. Verdier
\end{center}

\selectlanguage{english}version : \showgitstatus\selectlanguage{french}
\tableofcontents

\else
\chapter{Cohomologie dans les topos}\etiquette[V]{V}
%\addtocontents{toc}{J.-L. Verdier}
\begin{center}
J.-L. Verdier
\end{center}
\fi

\setcounter{pageoriginal}{1}

\section*{Introduction}

On\pageoriginale présente dans cet exposé les invariants
cohomologiques commutatifs et élémentaires des topos. Dans le ${\rm
n^{\circ}}$ 1, on étudie les modules plats et les morphismes plats
de topos annelés. Les démonstrations sont faites en utilisant
l'hypothèse, le plus souvent vérifiée dans la pratique, que les
topos ont suffisamment de points (IV 6). Ces démonstrations sont
reprises dans le cas général dans l'appendice nᵒ 8
où Deligne, à l'aide de la technique des limites inductives
locales, généralise en outre au cas des topos, le théorème de D.
Lazard sur la structure des modules plats. Les théorèmes de cet
exposé\sisi{,}{} sont des théorèmes d'existence de suites spectrales reliant
les différents invariants cohomologiques (nᵒ 3, 5, 6). On
sait que, même pour les espaces topologiques, la cohomologie de Čech
ne coïncide pas en général avec la cohomologie des faisceaux
\cite{SGA4-05:11}. On introduit dans l'appendice nᵒ 7, un
calcul de Čech modifié permettant d'obtenir, à l'aide de
recouvrement, la cohomologie des faisceaux dans un topos quelconque.
On est amené dans cet appendice à utiliser des recouvrements
simpliciaux (hyper-recouvrements) dont les invariants homotopiques
ont été étudiés dans \cite{SGA4-05:1} (\cf aussi \cite{SGA4-05:17}).

Les invariants cohomologiques introduits sont élémentaires en ce
sens que nous n'utilisons pas les catégories dérivées \cite{SGA4-05:12}. Le
lecteur familier avec ce langage fera immédiatement la traduction
des différents énoncés de cet exposé et pourra alors les
généraliser aux complexes et à l'hypercohomologie. Ce langage des
catégories dérivées est d'ailleurs utilisé dans la suite de ce
séminaire.

On se limite ici à la cohomologie commutative. Pour le $H^1$ non
commutatif, utilisé dans ce séminaire, et pour le $H^2$ non
commutatif, nous renvoyons à \cite{SGA4-05:9}. Les foncteurs qu'on dérive sont
additifs; on reste muet sur les structures multiplicatives (\cf \cite{SGA4-05:7}).

\setcounter{section}{-1}
\section{Généralités sur les catégories abéliennes}\etiquette[0]{V.sec0}

Dans \pageoriginale ce numéro nous rappelons quelques lemmes dont la
plupart se trouvent dans \cite{SGA4-05:11}.

\begin{enonce*}{Proposition 0.1}\etiquette[0.1]{V.prop0.1}
Soit $\mathfrak{A}$ une catégorie abélienne possédant un
générateur. Les conditions suivantes sont équivalentes:
\begin{enumerate}
\renewcommand{\theenumi}{\roman{enumi}}
\renewcommand{\labelenumi}{\theenumi)}
\item La catégorie $\mathfrak{A}$ vérifie l'axiome $AB$ $5$): Les
  petites sommes directes sont représentables et si $(X_i)$, $i \in
  I$, est une petite famille filtrante croissante de sous-objets d'un
  objet $X$ de $\mathfrak{A}$ et $Y$ est un sous-objet de $X$, on a
$$
(\displaystyle{\sup_i}\ X_i) \cap Y = \displaystyle{\sup_i}\ (X_i \cap Y) {\quoi}.
$$
\item Les petites limites inductives pseudo-filtrantes \emph{(I.2.7)} %% étiquette \XXX
sont représentables et commutent aux limites projectives finies.
\end{enumerate}

De plus, si les conditions ci-dessus sont remplies, les petites
limites inductives filtrantes sont universelles \emph{(I 2.6)}.
\end{enonce*}

\begin{enonce*}[remark]{Preuve}
Il est clair que (ii) $\Rightarrow$ (i). Pour montrer que
\sisi{i)}{(i)} $\Rightarrow$ (ii), et pour prouver l'assertion
supplémentaire, il suffit d'utiliser que $\mathfrak{A}$ est une
sous-catégorie pleine d'une catégorie de modules $\Mm$
sur un anneau convenable, telle que le foncteur d'inclusion $u :
\mathfrak{A} \rightarrow \Mm$ admette un adjoint à gauche
$v$ exact \cite{SGA4-05:5}. La vérification est alors triviale.
\end{enonce*}

\setcounter{subsection}{1}
\subsubsection{}\etiquette[0.1.1]{V.subsubsec0.1.1}

On sait \cite{SGA4-05:11} qu'une catégorie abélienne $\mathfrak{A}$
possédant un générateur et vérifiant l'axiome $AB$ 5) possède
suffisamment d'injectifs \ie tout objet se plonge dans un objet
injectif. De plus, d'après le résultat déjà cité \cite{SGA4-05:5}, les
petits produits sont représentables dans $\mathfrak{A}$ (axiome $AB$
3) $^{\ast}$).

\begin{enonce*}{Proposition 0.2}\etiquette[0.2]{V.prop0.2}
Soient $\mathfrak{A}$ et $\mathfrak{B}$ deux
  catégories abéliennes et
$\xymatrix{\mathfrak{A} \ar@<-2pt>[r]_{u}& \mathfrak{B} \ar@<-2pt>[l]_{v}
}
$
deux foncteurs adjoints ($u$ est adjoint à gauche de
$v$). Considérons les deux propriétés:
\begin{enumerate}
\renewcommand{\theenumi}{\roman{enumi}}
\renewcommand{\labelenumi}{\theenumi)}
\item Le foncteur $u$ est exact.

\item Le\pageoriginale foncteur $v$ transforme les objets injectifs de $\mathfrak{B}$
  en objets injectifs de $\mathfrak{A}$.
\begin{enumerate}
\renewcommand{\theenumii}{\arabic {enumii}}
\renewcommand{\labelenumii}{\theenumii)}
\item On a toujours l'implication \emph{(i) $\Rightarrow$ (ii)}.

Si, de plus, tout objet non nul de $\mathfrak{B}$ est source d'un morphisme non
  nul dans un objet injectif, alors \emph{(ii) $\Leftrightarrow$ (i)}.
\item Supposons que:
\begin{enumerate}
\renewcommand{\theenumiii}{\alph {enumiii}}
\renewcommand{\labelenumiii}{\theenumiii)}
\item la catégorie $\mathfrak{B}$ possède suffisamment
  d'injectifs.
\item l'une des deux conditions équivalentes {\rm(i)} et {\rm(ii)} ci-dessus
  soit remplie.
\item le foncteur $u$ soit fidèle.
\end{enumerate}
\end{enumerate}
\end{enumerate}
Alors, la catégorie $\mathfrak{A}$ possède suffisamment d'injectifs.
\end{enonce*}

\begin{enonce*}[remark]{Preuve}
La preuve est laissée au lecteur.
\end{enonce*}

\begin{enonce*}[remark]{Remarque 0.2.1}\etiquette[0.2.1]{V.rem0.2.1}
La catégorie des groupes commutatifs possède suffisamment
d'injectifs. Appliquant le lemme \hyperref[V.prop0.2]{0.2} on en déduit que toute
catégorie de modules unitaires sur un anneau à élément unité
possède suffisamment d'injectifs. Appliquant alors le résultat de
\cite{SGA4-05:5} (utilisé dans la preuve de \hyperref[V.prop0.1]{0.1}) et \hyperref[V.prop0.2]{0.2}, on en déduit que toute
catégorie abélienne possédant un générateur et des petites
limites inductives filtrantes exactes possède suffisamment
d'injectifs; ce qui fournit une nouvelle démonstration de ce fait.
\end{enonce*}

\begin{enonce*}{Proposition 0.3}\etiquette[0.3]{V.prop0.3}
Soient $\mathfrak{A}$, $\mathfrak{B}$, $\mathfrak{C}$ trois
catégories abéliennes et $u : \mathfrak{A} \rightarrow
\mathfrak{B}$, $v : \mathfrak{B} \rightarrow \mathfrak{C}$ deux
foncteurs additifs exacts à gauche. Supposons que $\mathfrak{A}$ et
$\mathfrak{B}$ possèdent suffisamment d'objets injectifs. Les deux
propriétés suivantes sont équivalentes:
\begin{enumerate}
\renewcommand{\theenumi}{\roman{enumi}}
\renewcommand{\labelenumi}{\theenumi)}
\item Il existe un foncteur spectral dont le terme $E^{p, q}_2$ est:
$$
R^p vR^qu
$$
et qui aboutit à $R^{p+q} vu$ (convenablement filtré).
\item Le foncteur $u$ transforme les objets injectifs de
  $\mathfrak{A}$ en objets acycliques pour le foncteur $v$.
\end{enumerate}
\end{enonce*}

\begin{enonce*}[remark]{Preuve}
(i) $\Rightarrow$ (ii) est trivial car il suffit d'appliquer le foncteur
  spectral à un\pageoriginale objet injectif. L'implication (ii) $\Rightarrow$ (i)
  est démontrée dans \cite{SGA4-05:11}.
\end{enonce*}

\begin{enonce*}{Proposition 0.4}\etiquette[0.4]{V.prop0.4}
Soient $\mathfrak{A}$ et $\mathfrak{B}$ deux catégories abéliennes
et $u : \mathfrak{A} \rightarrow \mathfrak{B}$ un foncteur %% césure fonc-teur correcte ?
additif exact à gauche. Soit $M$ un sous-ensemble de l'ensemble des objets
de $\mathfrak{A}$ possédant les propriétés suivantes:
\begin{enumerate}
\renewcommand{\labelenumi}{\theenumi)}
\item Tout objet de $\mathfrak{A}$ se plonge dans un élément de
  $M$.

\item Si $X \oplus Y$ appartient à $M$, l'objet $X$ appartient à
  $M$.

\item Si
$$
0 \longrightarrow X' \longrightarrow X \longrightarrow X''
\longrightarrow 0
$$
est une suite exacte et si $X'$ et $X$ appartiennent à $M$, alors
$X''$ appartient à $M$ et la suite
$$
0 \longrightarrow u(X') \longrightarrow u(X) \longrightarrow u(X'')
\longrightarrow 0
$$
est exacte. Les objets nuls appartiennent à $M$.
\end{enumerate}

Alors tout injectif appartient à $M$, et les objets de $M$ sont
acycliques pour $u$, \ie pour tout $q \neq 0$ et tout objet $X$ de
$M$, on a $R^q u(X) = 0$. (En particulier les résolutions par des
objets de $M$ permettent de calculer les foncteurs dérivés de $u$.)
\end{enonce*}

Pour la preuve\sisi{}{,} voir \cite{SGA4-05:11} 3.3.1.

\begin{enonce*}{Proposition 0.5}\etiquette[0.5]{V.prop0.5}
Soient $𝒰 \subset 𝒱$ deux univers, $\mathfrak{A}$
(\resp $\mathfrak{B}$) une $𝒰$-catégorie (\resp $𝒱$-catégorie)
abélienne vérifiant l'axiome $AB$ 5) relativement à $𝒰$
(\resp à $𝒱$) et possédant une famille génératrice $𝒰$-petite
(\resp $𝒱$-petite). Soit $\epsilon : \mathfrak{A} \to \mathfrak{B}$
un foncteur exact et pleinement fidèle. Les conditions suivantes
sont équivalentes:
\begin{itemize}
\item[1)] Il existe une famille génératrice $(X_i)_{i\in I}$ de
  $\mathfrak{A}$ telle que la famille $(\epsilon(X_i))_{i\in I}$
  soit génératrice dans $\mathfrak{B}$.

\item[1′)] Tout objet de $\mathfrak{B}$ est isomorphe à un quotient
  d'un objet du type $\displaystyle{⨁_{\alpha \in A}} \epsilon
  (Y_\alpha)$ où $A \in 𝒱$.

   Sous ces conditions, on a la propriété suivante:

\item [2)] $\epsilon$ transforme les produits $𝒰$-petits en
  produits (donc commute aux limites projectives $𝒰$-petites).
\end{itemize}

De plus, sous les conditions équivalentes 1) ou 1′), les conditions
suivantes sont équivalentes\pageoriginale:
\begin{itemize}
\item[a)] Pour tout objet $Y$ de $\mathfrak{A}$, tout sous-objet de
  $\epsilon(Y)$ est isomorphe à l'image par $\epsilon$ d'un
  sous-objet de $Y$.

\item[a′)] Il existe une famille génératrice $(X_i)_{i \in I}$ de
  $\mathfrak{A}$ telle que la famille $(\epsilon(X_i))_{i \in I}$
  soit génératrice dans $\mathfrak{B}$ et telle que pour tout $i
  \in I$, tout sous-objet de $\epsilon(X_i)$ soit isomorphe à
  l'image par $\epsilon$ d'un sous-objet de $X_i$.

\item[b)] Tout objet de $\mathfrak{B}$ est isomorphe à un sous-objet d'un
  objet du type $\prod_{\alpha \in A} \epsilon(Y_\alpha)$ où
  $A \in 𝒱$.

\item[c)] $\epsilon$ commute aux sommes directes $𝒰$-petites (donc
  commute aux limites inductives $𝒰$-petites.

De plus, sous les conditions 1) et a′) on a :

\item[d)] $\epsilon$ transforme les objets injectifs en objets injectifs.
\end{itemize}
\end{enonce*}

\begin{enonce*}[remark]{Remarque 0.5.1}\etiquette[0.5.1]{V.rem0.5.1}
Lorsque dans \hyperref[V.prop0.5]{0.5} on a $𝒰 = 𝒱$, les conditions 1) et
a′) entraînent
que $\epsilon$ est une équivalence de catégories (car on a
alors b), 2) et a)).
\end{enonce*}

\setcounter{subsection}{5}
\subsubsection{}\etiquette[0.5.2]{V.subsubsec0.5.2}
Nous nous bornerons à donner des indications sur la démonstration.
Les implications 1) $\Leftrightarrow$ $1'$), 1) $\Rightarrow$ 2),
sont laissées au lecteur. Il est clair que a) $\Rightarrow$
a′). Montrons que a′) entraîne d). Soit $M$ un injectif
dans $\mathfrak{A}$. Pour tout $i \in I$ et tout sous-objet $Y
\hookrightarrow X_i$ de $X_i$, l'homomorphisme $\Hom(X_i, M)
\rightarrow \Hom(Y, M)$ est surjectif. Donc, en vertu de a′) et de
la pleine fidélité de $\epsilon$, pour tout sous-objet $U
\hookrightarrow \epsilon(X_i)$, l'homomorphisme $\Hom(\epsilon
(X_i), \epsilon(M)) \to \Hom(U, \epsilon(M))$ est surjectif. Comme
les $\epsilon(X_i)$ forment une famille génératrice de
$\mathfrak{B}$, $\epsilon(M)$ est injectif \cite{SGA4-05:11}.

Montrons que a′) $\Rightarrow$ b). Quitte à augmenter la famille des
$X_i$, on peut supposer que pour tout $i \in I$, tout quotient de
$X_i$ est isomorphe à un $X_j$ pour un $j$ convenable. Soit alors,
pour tout $i \in I$, $X_i \hookrightarrow M_i$ un monomorphisme dans
un objet injectif. La famille $(\epsilon (X_i))_{i\in I}$ est stable
par quotient et pour $i \in I$, le morphisme $\epsilon(X_i)
\hookrightarrow \epsilon(M_i)$ est, d'après d), un monomorphisme
dans un objet injectif. On vérifie alors immédiatement que, la
famille $\epsilon(X_i)$ étant génératrice, la famille
$(\epsilon(M_i))_{i \in I}$ est cogénératrice, d'où b).

Montrons que b) $\Rightarrow$ c). Soit $(Z_\alpha)_{\alpha \in A}$ une famille
$𝒰$-petite d'objets de $\mathfrak{A}$ et montrons que le morphisme
canonique
$$
⨁_{\alpha\in A} \epsilon(Z_\alpha) \to \epsilon(⨁_{\alpha \in A} Z_\alpha)
$$
est\pageoriginale un isomorphisme. Comme la famille des objets
$\epsilon(Y)$ est cogénératrice, il suffit montrer que pour
tout objet $Y$ de $\mathfrak{A}$, l'homomorphisme
$$
\Hom(\epsilon(⨁_\alpha Z_\alpha), \epsilon(Y)) \to
\Hom(⨁_\alpha \epsilon(Z_\alpha), \epsilon(Y))
$$
est un isomorphisme ce qui résulte de la pleine fidélité de $\epsilon$.

Il reste à montrer que c) $\Rightarrow$ a). Soit $Z$ un
sous-objet de $\epsilon(Y)$. Il existe, en vertu de $1'$), une
famille $𝒰$-petite $Y_\alpha$ d'objets de $\mathfrak{A}$ et
un épimorphisme de $\bigoplus_{\alpha} \epsilon(Y_\alpha)$ sur $Z$.
Comme $\epsilon$ commute aux sommes directes $𝒰$-petites, il existe
donc un épimorphisme $\epsilon(Y') \to Z$. Notons $u : \epsilon(Y')
\to Z \to \epsilon(Y)$ le morphisme composé. On a $Z =
\text{Im}(u)$. Comme \sisi{}{$\epsilon$} est pleinement fidèle, on a $u \sisi{:}{=} \epsilon(v)$
et par suite $Z = \text{Im}(\epsilon(\sisi{V}{v})) = \epsilon(\text{Im}(v))$.


\begin{enonce*}[remark]{Exercice 0.5.2}\etiquette[0.5.2]{V.exe0.5.2}
Soit $\Sex_{𝒱} (\mathfrak{A})$ la catégorie des foncteurs
contravariants
de $\mathfrak{A}$ dans la catégorie des $𝒱$-groupes commutatifs qui
commutent aux limites inductives $𝒰$-petites. Montrer que sous les
conditions 1) et a′) le foncteur canonique de $\mathfrak{B}$ dans
$\Sex_{𝒱}(\mathfrak{A})$ est une équivalence.
\end{enonce*}

\section{Modules plats}\etiquette[1]{V.sec1}

\begin{enonce*}{Définition 1.1}\etiquette[1.1]{V.def1.1}
Soit $(E, A)$ un topos annelé \emph{(IV 11.1.1)}. %% étiquette
Un $A$-Module à droite
(\resp à gauche) $M$ est dit plat si le foncteur $M
\otimes_A ⋅$ (\resp $⋅\otimes_A M$) de la catégorie des $A$-Modules
à gauche (\resp à droite) dans la catégorie des faisceaux
abéliens de $E$\sisi{,}{} est exact.
\end{enonce*}

\begin{enonce*}{Proposition 1.2}\etiquette[1.2]{V.prop1.2}
Soit $M$ un $B$-$A$ bi-Module.
\begin{enumerate}
\renewcommand{\labelenumi}{\theenumi)}
\item Les propriétés suivantes sont équivalentes:
\renewcommand{\theenumii}{\roman {enumii}}
\renewcommand{\labelenumii}{\theenumii)}
\begin{enumerate}
\item Le module $M$ est $A$-plat à gauche\sisi{}{.}
\item Pour tout $B$-Module injectif $I$, le  $A$-Module à gauche
  $\Hhom_{B}(M, I)$ est injectif.
\end{enumerate}
\item Un Module $M$, limite inductive pseudo-filtrante \emph{(I 2.7.1)} %% étiquette
de Modules plats, est plat.

\item Enfin, si $M_∙ = \cdots \sisi{}{\to} M_{i + 1} \to M_i \sisi{}{\to} \cdots$ est un complexe
  acyclique de Modules plats $(M_i = 0$ pour $i < i_0)$, alors
  pour tout Module $F$, le complexe:
$$
M_∙ \otimes_A F = \cdots \sisi{}{\to} M_{i + 1} \otimes_A F \to M_i \otimes_A F \sisi{}{\to} \cdots
$$
est\pageoriginale acyclique.
\end{enumerate}
\end{enonce*}

\begin{enonce*}[remark]{Preuve}
D'après (IV 12.12) %% étiquette
on a un isomorphisme d'adjonction :
$$
\Hom_B (M \otimes_A \cdot , \cdot) \xrightarrow{\sim} \Hom_A (\cdot ,
\Hhom_B(M , \cdot) ){\quoi}.
$$
Il suffit alors d'appliquer \hyperref[V.prop0.2]{0.2} pour obtenir
l'équivalence (i) $\Leftrightarrow$ (ii). Cet isomorphisme
d'adjonction montre par ailleurs que le produit tensoriel commute
aux limites inductives. Le fait que les limites inductives
pseudo-filtrantes soient exactes (\hyperref[V.prop0.1]{0.1}) entraîne la
deuxième assertion. Pour montrer que le complexe $M_∙ \otimes_A F$
est acyclique, il suffit de montrer que pour tout faisceau abélien
injectif I le complexe $\Hom^{∙}_{\ZZ} (M_∙ \otimes_A F, I)$ est
acyclique. Ce complexe est isomorphe, en vertu des formules
d'adjonction, au complexe $\Hom^{∙}_A(F, \Hhom_{\ZZ}(M_∙ , I))$. Or
d'après l'équivalence (i) $\Leftrightarrow$ (ii), le complexe de
faisceaux
$$
\Hhom_{\ZZ}(M_∙ , I)
$$
est un complexe acyclique dont les objets sont injectifs, d'où la conclusion.
\end{enonce*}

\begin{enonce*}{Proposition 1.3.1}\etiquette[1.3.1]{V.prop1.3.1}
Soient $(E, A)$ un topos annelé, $H$ un objet de $E$, $M$ un
$A_{/H}$-Module plat. Alors $j_{H!}M$ est un $A$-Module plat. En
particulier $A_H$ est plat à droite et à gauche.
\end{enonce*}

Supposons, pour fixer les idées, que $M$ soit un $A_{/H}$-Module à
droite. Pour tout $A$-Module à gauche $P$, on a un isomorphisme
canonique (IV 12) %% étiquette
$$
P \otimes_A j_{H!} M \simeq j_{H!} (P \otimes_{A_{/H}} M){\quoi}.
$$
Les foncteurs $j_{H!}$ et $j^{\ast}_H$ sont exacts (IV 11.3.1 et 11.12.2)
et par hypothèse, le foncteur $-\otimes_{A_{/H}} M$ est exact. Par
suite le foncteur $P \to P \otimes_A j_{H!}M$ est exact et $j_{H!}M$
est plat.

\begin{enonce*}{Proposition 1.3.2}\etiquette[1.3.2]{V.prop1.3.2}
(Formule de projection pour les immersions fermées):

Soient $(E,A)$ un topos annelé, $i : F \to E$ un sous-topos fermé de
$E$. Posons $i^{\ast} A = A_{/F}$. Pour tout $A_{/F}$-Module (à droite) $M$ et
tout $A$-Module (à gauche) $P$, on\pageoriginale a un isomorphisme fonctoriel
$$
i_\ast (M \otimes_{A_{/F}} i^\ast P) \simeq (i_\ast M) \otimes_A P
$$
\end{enonce*}

Soient $U$ l'ouvert complémentaire de $F$ et $j : U \to E$ le
morphisme canonique d'immersion. On a $j^{\ast} (i_\ast M \otimes_A P)
\simeq 0 \otimes_{A_{/U}} j^\ast (P)$ (IV 12); par suite $i_{\ast} M
\otimes_{A} P$ a son support dans $F$. Donc le morphisme d'adjonction
$i_{\ast} M \otimes_{A} P \to i_\ast i^\ast (i_\ast M_A P)$ est un
isomorphisme (IV 14). On a $i^{\ast} (i_\ast M \otimes_{A} P) \simeq
i^\ast i_\ast M \otimes_{A_{/F}} i^\ast P$ (IV 12) et comme $i : F \to
E$ est une immersion fermée, $i^{\ast} i_{\ast} M \simeq M$ (IV 14);
d'où l'isomorphisme annoncé.

\begin{enonce*}{Corollaire 1.3.3}\etiquette[1.3.3]{V.cor1.3.3}
Pour tout $A_{/F}$-Module plat $M$, $i_{\ast} M$ est plat.
\end{enonce*}

Il résulte de \hyperref[V.prop1.3.2]{1.3.2}, et de (IV 14) que le foncteur $P \mapsto
(i_\ast M) \otimes_A P$ est un foncteur exact.

\setcounter{subsection}{3}
\subsection{}\etiquette[1.4]{V.subsec1.4}
Soient $(E, A)$ un topos annelé, $x : P \to E$ un point de $E$ (IV
6.1), $\Vois(x)$ la catégorie des voisinages de $x$ (IV 6.8). À tout
objet $V$ de $\Vois(x)$ correspond un objet de $E$, encore noté $V$,
et un point $x_V : P \to E/V$ de $E/V$. De plus, à tout morphisme $u
: V \to W$ de $\Vois(x)$, correspond un diagramme essentiellement
commutatif de morphisme de topos (IV 6.7)
\[\vcenter{
\xymatrix@C=2.5cm@R=1.5cm{
P \ar[r]^{x_V} \ar[rd]_{x_W}& E/V \ar[d]^{j_u}\\
& E/W\quad .
}}\leqno{(1.4.1)}
\]\etiquette[1.4.1]{V.1.4.1}
Soit $N$ un $A_x$-module (\resp un ensemble). Du diagramme \eqhyperref[V.1.4.1]{1.4.1}, on
déduit un diagramme de $A_x$-modules (\resp d'ensembles):
\[\vcenter{
\xymatrix@C=2cm@R=1.5cm{
x^{\ast}_W x_{W\ast}N \ar[r]^{\ad_W} \ar[d]_{\phi(u)}& N\\
x^{\ast}_V x_{V\ast}N \ar[ru]_{\ad_{V}}&
}}\leqno{(1.4.2)}
\]\etiquette[1.4.2]{V.1.4.2}
où \pageoriginale $\ad_W$ et $\ad_V$ sont les morphismes
d'adjonction. De plus, on vérifie immédiatement que $\phi (uv) =
\phi (v) \phi{(u)}$. On a donc un foncteur de la catégorie filtrante
$\Vois(x)^{\circ}$ dans la catégorie des $A_x$-modules
(\resp ensembles) et un homomorphisme:
\begin{equation}
\Lambda (N) : \varinjlim_{V \in \Vois(x)^{\circ}}
x^{\ast}_V x_{V\ast}N \to N {\quoi}.\tag{1.4.3}
\end{equation}\etiquette[1.4.3]{V.1.4.3}

\begin{enonce*}{Proposition 1.5}\etiquette[1.5]{V.prop1.5}
Pour tout $A_x$-module\sisi{s}{} (\resp ensemble) $N$, $\Lambda(N)$ est un isomorphisme.
\end{enonce*}

Cette proposition est un cas particulier d'une proposition due à
Deligne (\hyperref[V.prop8.2.6]{8.2.6}). Donnons-en une démonstration directe. En passant
aux ensembles sous-jacents\sisi{}{,} il suffit de démontrer la proposition
lorsque $N$ est un ensemble (I 2.8). %% étiquette
La catégorie cofiltrante $\mathrm{Fl}(\Vois(x))$ des flèches de $\Vois(x)$ est fibrée sur la
catégorie $\Vois(x)$ par le foncteur $p : \mathrm{Fl}(\Vois(x)) \to
\Vois(x)$ qui, à une flèche, associe son but. Soit $D$ une
catégorie et $F : \mathrm{Fl}(\Vois(x)) \to D$ un pro-objet (I 8.10). %% étiquette
Pour tout objet $V$ de $\Vois(x)$, notons $F_V$ le pro-objet obtenu en
restreignant le foncteur $F$ à la catégorie fibre $\Vois(x)/V$. A
tout morphisme $m : U \to V$ de $\Vois(x)$, le foncteur changement de
base par $m$ associe un morphisme du pro-objet $F_U$ dans le pro-objet
$F_V$ et les morphismes canoniques de pro-objets $F \to F_V$
déterminent un morphisme de pro-objets:
\begin{equation*}
F \to \varprojlim_{V\in \Vois(x)^{\circ}} F_V\tag{1.5.1}
\end{equation*}\etiquette[1.5.1]{V.1.5.1}
dont on vérifie immédiatement que c'est un
isomorphisme. Appliquons cette remarque au pro-objet d'ensembles
pointés:
\begin{equation*}
(U, V, m) \longmapsto x^{\ast}_V (m) = F(U, V, m){\quoi},\tag{1.5.2}
\end{equation*}\etiquette[1.5.2]{V.1.5.2}
on obtient un isomorphisme de pro-objets:
\begin{equation*}
F \xrightarrow{\sim} \varprojlim_{\Vois(x)}
\varprojlim_{\Vois(x)/V} x^{\ast}_V(m){\quoi};\tag{1.5.3}
\end{equation*}\etiquette[1.5.3]{V.1.5.3}
d'où, pour tout ensemble $N$, une bijection
\begin{equation*}
\varinjlim_{\mathrm{Fl} (\Vois(x))^{\circ}} \Hom(x^{\ast}_V(m), N) \simeq
\varinjlim_{\Vois(x)^{\circ}}
\varinjlim_{(\Vois(-x)/V)^{\circ}} \Hom(x^{\ast}_V (m), N).\tag{1.5.4}
\end{equation*}\etiquette[1.5.4]{V.1.5.4}
Mais,\pageoriginale pour $V$ fixé,
$\varinjlim_{(\Vois(x)/V)^{\circ}} \Hom(x^{\ast}_V (m), N)$
s'identifie à $x^{\ast}_V x_{V\ast} N$. En effet, on a
$x^{\ast}_V x_{V\ast} N \simeq \varinjlim_{\Vois(x_V)^\circ}
\Hom_{E/V} (W, x_{V\ast}N)$ d'après IV 6.8.1 donc, par adjonction,
on a $x^{\ast}_V x_{V\ast} N \simeq \varinjlim_{\Vois(x_V)^\circ}
\Hom(x^{\ast}_V(w), N)$ et de plus la catégorie $\Vois(x_V)$ est
équivalente à la catégorie $\Vois(x)/V$ (IV 6.7.2). Enfin, les
applications canoniques de transition
$\varinjlim_{(\Vois(x)/U)^{\circ}} \Hom(x^{\ast}_U(m), N) \to
\varinjlim_{(\Vois(x)/V)^{\circ}} \Hom(x^{\ast}_V(n), N)$ dans
\eqhyperref[V.1.5.4]{1.5.4} s'identifient aux applications canoniques $x^\ast _U
x_{U\ast} N \to x^{\ast}_V x_{V\ast} N$ \eqhyperref[V.1.4.3]{1.4.3} ainsi que le
lecteur voudra bien le vérifier. On a donc une bijection
\begin{equation*}
\varinjlim_{\mathrm{Fl}(\Vois(x))^{\circ}} \Hom(x^{\ast}_V(m), N) \simeq
\varinjlim_{\Vois(x)^{\circ}} x^{\ast}_V x_{V\ast} N \tag{1.5.5}
\end{equation*}\etiquette[1.5.5]{V.1.5.5}
et l'application $\Lambda(N): \varinjlim_{\Vois(x)^{\circ}} x^{\ast}_V
x_{V\ast} N \to N$ \eqhyperref[V.1.4.3]{1.4.3} composée avec la bijection
\eqhyperref[V.1.5.5]{1.5.5} provient des applications $\Hom(x^{\ast}_V (m), N) \to N$ qui a
une application $r : x^{\ast}_V(m) \to N$ associe l'image par $r$ du
point marqué de $x^{\ast}_V (m)$. Pour démontrer la proposition,
il suffit alors de remarquer que tout objet $(U, V, m)$ de
$\mathrm{Fl}(\Vois(x))$ est minoré par $(U, U, \id_U)$ et que
$x^{\ast}_U(\id_U)$ est réduit à un élément ou, en d'autres
termes, que le morphisme canonique du pro-objet constant réduit à
un élément dans $F$ est un isomorphisme de pro-objets.

\begin{enonce*}{Proposition 1.6}\etiquette[1.6]{V.prop1.6}
Soient $(E, A)$ un topos annelé, $M$ un $A$-Module.
\begin{enumerate}
\renewcommand{\labelenumi}{\theenumi)}
\item Lorsque $M$ est plat, pour tout point $x : P \to E$ de $E$, le
  $A_x$-module $M_x$ est plat.
\item Soit $(x_i)_{i \in I}$ une famille conservatrice de \sisi{point}{points} de $E$
  telle que pour tout $i \in I$, $M_{x_i}$ soit un $A_{x_i}$-module plat.
\end{enumerate}
Alors $M$ est plat.
\end{enonce*}

\begin{enonce*}{Lemme 1.6.1}\etiquette[1.6.1]{V.lem1.6.1}
Soient $M$ un Module plat et $V$ un objet de $E$. Le $A/V$-Module
$j^{\ast}_V M$ est plat.
\end{enonce*}

Il faut montrer que le foncteur $P \mapsto P \otimes_{A/V}
j^{\ast}_V M$ est exact. Comme le foncteur prolongement par zéro
$j_{V!}$ est exact et fidèle (IV 11.3.1), il suffit de montrer que
le foncteur $P \mapsto j_{V!} (P \otimes_{A/V} j^{\ast}_V M)$ est
exact. On a un isomorphisme fonctoriel \pageoriginale $j_{V!} (P
\otimes_{A/V} j^{\ast}_V M) \simeq j_{V!} (P) \otimes_A M$ (IV
12.11) et par suite le foncteur $P \mapsto j_{V!} (P) \otimes_A M$
est exact (IV 11.3.1).

\setcounter{subsection}{6}
\setcounter{subsubsection}{1}
\subsubsection{}\etiquette[1.6.2]{V.subusbsec1.6.2}
Démontrons 1).
Supposons pour fixer les idées que $M$ soit un $A$-Module à
  gauche plat et soit $x : P \to E$ un point de $E$. Montrons que
  le foncteur $N \mapsto N\otimes_{A_x} M_x$ de la catégorie des
  $A_x$-modules à droite dans la catégorie des groupes commutatifs
  est exact. Il suffit pour cela de montrer que ce foncteur est exact
  à gauche. Avec les notations de \hyperref[V.subsec1.4]{1.4}, soient $V$ un voisinage de
  $x$ et $j_V: E/V \to E$ le morphisme de localisation. On a, pour
  tout $A_x$-module $N$, $x^\ast_V x_{V\ast} N \otimes_{A_x} M_x
  \simeq x^{\ast} _V x_{V\ast} N \otimes _{A_x}
  x^\ast_{V}j^{\ast}_V M \simeq x^\ast_V (x_{V_\ast} N \otimes_{A_x}
  j^\ast_V M)$ (IV 12.11). Donc le foncteur $N \mapsto x^\ast_V
  x_{V\ast} N \otimes_{A_x} M_x$ est exact à gauche. Par suite, le
  foncteur $N \mapsto N \otimes_{A_x} M_x$, limite inductive filtrante
  de foncteurs exacts à gauche (\hyperref[V.prop1.5]{1.5}), est exact à gauche.
Démontrons 2). Soit $0 \to P' \to P \to P'' \to 0$ une suite exacte
de $A$-Modules et montrons que la suite $0 \to P' \otimes_A M \to P
\otimes_A M \to P'' \otimes_A M \to 0$ est exacte. Il suffit pour cela
de montrer que pour tout $i \in I$, la suite obtenue en passant aux
fibres en $x_i$ est exacte (IV 6). Or la suite des fibres en $x_i$ est
la suite (IV 13.5)
$0 \to P'_{x_i} \otimes_{A_x} M_x \to P_{x_i} \otimes_{A_{x_i}} M_{x_i}
\to P''_{x_i} \otimes_{A_{x_i}} M_{x_i} \to 0$ et comme $M_{x_i}$ est
plat, cette suite est exacte.

\begin{enonce*}{Proposition 1.7}\etiquette[1.7]{V.prop1.7}
Soient $(E, A)$ un topos annelé, $u : E' \to E$ un morphisme de
topos et posons $A' = u^{\ast} A$. Soit $M$ un $A'$-Module à droite
(\resp à gauche). Les conditions suivantes sont équivalentes:
\begin{enumerate}
\renewcommand{\theenumi}{\roman{enumi}}
\renewcommand{\labelenumi}{\theenumi)}
\item Le foncteur $N \mapsto M \otimes_{A'} u^{\ast} N$ (\resp $N
  \mapsto u^{\ast} N \otimes_{A'} M$) de la catégorie des $A$-Modules
  à gauche (\resp à droite) dans la catégorie des faisceaux
  abéliens de $E'$ est exact.
\item $M$ est un $A'$-Module plat.
\end{enumerate}
\end{enonce*}

Il est clair que ii) $\Rightarrow$ i).

Montrons que i) $\Rightarrow$ ii). Nous ne ferons la démonstration
que dans le cas où $E'$ possède une famille conservatrice de
points $(x_i)_{i \in I}$. Le cas général est traité en
(\hyperref[V.prop8.2.7]{8.2.7}). Dans ce cas particulier, qui couvre la plupart
des applications, on est ramené \pageoriginale au cas où $E'$
est le topos ponctuel et où, par suite, le morphisme $u$, qu'on
notera $x$, est un point de $E$ (\hyperref[V.prop1.6]{1.6} et IV 11.3.1). Nous nous
bornerons au cas où $M$ est un $A'$-Module à droite. Le cas
où $M$ est un $A'$-Module à gauche s'y ramène en passant aux
anneaux opposés. Soient $V$ un voisinage de $x$ et $V_x$ sa fibre en
$x$. On a, avec les notations de \hyperref[V.subsec1.4]{1.4}, un diagramme
essentiellement commutatif de morphismes de topos:
\[
\xymatrix@C=2.25cm@R=1.5cm{
P \ar[r]^{j} \ar[rd]_{x_V}& P/V_{x}\ar[r]^{j_{V_x}} \ar[d]^{x/V}& P\ar[d]^{x}\\
  &  E/V \ar[r]   & E\quad \sisi{,}{}
}
\]
où $x /V$ est le morphisme déduit de $x$ par localisation sur
$V$ (IV 5.10) et $j$ est le morphisme déduit du point marqué de
$V_x$. Montrons que le foncteur $N' \mapsto M \otimes_{A'} x^{\ast}_V
N'$ est exact. On a $x^\ast_V = j^\ast(x / V)^{\ast}$ et $\id = j^*
j^*_{V_x}$ et par suite, on a un isomorphisme canonique $M
\otimes_{A'} x^*_V N' \simeq j^*(j^*_{V_x} M \otimes_{A'/V_x} (x/
V)^*N')$. Comme le foncteur $j^*$ est exact et comme le foncteur
$j_{V_x !}$ est exact et fidèle (IV 11.3.1), il suffit de montrer
que le foncteur $N'\mapsto j_{V_{x}!}$ $(j^{\ast}_{V_{x}}M\otimes
(x/V)^{\ast}N')$ est exact, et par suite (IV 12.11), il suffit de
montrer que le foncteur $N' \mapsto M \otimes_{A'} j_{V_{x}!}(x/V)^{\ast}N'$ est
exact. Mais il résulte de la définition de l'image inverse d'un topos
induit (\hyperrefext[IV][IV.eq5.10.2]{5.10.2}) que $j_{V_{x}!}(x/V)^{\ast}N'$
est égal à $x^*j_{V!}N'$ et comme le foncteur $j_{V!}$ est exact
(\hyperrefext[IV][IV.prop11.3.1]{11.3.1}), le foncteur $N' \mapsto M \otimes_{A'} x^*j_{V!}N'$ est exact par hypothèse.
Pour tout $A'$-module $Q$, on a un isomorphisme
fonctoriel (\hyperref[V.prop1.5]{1.5}):
$$
Q \simeq \varinjlim_{\Vois(x)^{\circ}} x^*_V x_{V*} Q{\quoi}.
$$
D'après ce qui précède, le foncteur $Q \mapsto M \otimes_{A'} Q$ est
limite inductive filtrante de foncteurs exacts à gauche. Il est donc
exact à gauche et par suite $M$ est plat.

\begin{enonce*}{Corollaire 1.7.1}\etiquette[1.7.1]{V.cor1.7.1}
Soient $u : (E', A') \to (E, A)$ un morphisme de topos annelés,
\sisi{}{et} $M$ un $A$-Module plat. Alors $u^*M$ est un $A'$-Module plat.
\end{enonce*}

Résulte du critère donné dans \hyperref[V.prop1.7]{1.7} et de la formule (IV 13.4.5).

\begin{enonce*}{Corollaire 1.7.2}\etiquette[1.7.2]{V.cor1.7.2}
Soit\pageoriginale $u : (E', A') \to (E, A)$ un morphisme de topos annelés. Les
conditions suivantes sont équivalentes:
\begin{enumerate}
\renewcommand{\theenumi}{\roman{enumi}}
\renewcommand{\labelenumi}{\theenumi)}
\item le $u^{-1} (A)$-Module à droite (\resp à gauche) $A'$ est
  plat.
\item le foncteur $u^{*}$ de la catégorie des $A$-Modules à gauche
  (\resp à droite) dans la catégorie des $A'$-Modules à gauche
  (\resp à droite) est exact.
\end{enumerate}
\end{enonce*}

L'équivalence résulte de \hyperref[V.prop1.7]{1.7} et de la définition de $u^\ast$ (IV
13.2.1).

\begin{enonce*}{Définition 1.8}\etiquette[1.8]{V.def1.8}
Un morphisme de topos \sisi{annelé}{annelés} qui possède les propriétés
équivalentes de \hyperref[V.cor1.7.2]{1.7.2} est appelé un morphisme de topos plat à
gauche (\resp à droite).
\end{enonce*}

\begin{enonce*}{Proposition 1.9}\etiquette[1.9]{V.prop1.9}
Soient $𝒰 \subset 𝒱$ deux univers, $C$ un $𝒰$-site, $C^{\sim}_{𝒰}$ et
$C^{\sim}_{𝒱}$ les catégories des
$𝒰$ et $𝒱$-faisceaux d'ensembles
respectivement, $\epsilon: C^{\sim}_{𝒰} \to C^{\sim}_{𝒱}$ le foncteur
d'inclusion canonique, $A$ un $𝒰$-faisceau d'anneaux sur $C$.
\begin{enumerate}
\renewcommand{\labelenumi}{\theenumi)}
\item Le foncteur $\epsilon$ commute aux limites inductives et
  projectives $𝒰$-petites de Modules. Le foncteur $\epsilon$ est
  conservatif et pleinement fidèle sur les catégories de Modules.

\item Pour tout objet $H$ de $C^{\sim}_{𝒰}$, il
  existe un isomorphisme
  canonique $\epsilon(A_H) \simeq \epsilon
  (A)_{\epsilon(H)}$.

\item Tout $A$-Module injectif à gauche ou à droite est
  transformé par $\epsilon$ en $\epsilon(A)$-Module injectif.

\item Tout $A$-Module plat à droite ou à gauche est transformé
  par $\epsilon$ en $\epsilon(A)$-Module plat.

\item Le foncteur $\epsilon$ commute à la formation du produit
  tensoriel et du faisceau des homomorphismes.
\end{enumerate}
\end{enonce*}

La formation du faisceau associé ne dépend pas de l'univers (II
3.6). Il résulte alors de la construction des limites inductives et
projectives dans les catégories de faisceaux (II 4.1 et 6.4) que le
foncteur $\epsilon$ commute aux $𝒰$-petites limites projectives
d'ensembles ou de Modules, d'où 1). Démontrons 2). En prenant une
$𝒰$-petite sous-catégorie génératrice de $
C^{\sim}_{𝒰}$, on peut se
ramener au cas où $C$ est $𝒰$-petite\pageoriginale (III
4.1). Il résulte alors
de (II 6.5) que le $A$-Module libre engendré par $H$ s'obtient en
formant le préfaisceau de $A$-modules libres engendré par $H$,
formation qui commute à l'agrandissement des univers, puis en
formant le faisceau associé, opération qui elle aussi commute à
l'agrandissement des univers (II 3.6). Pour démontrer 3) il suffit
de montrer, en vertu de \hyperref[V.prop0.5]{0.5}, que tout sous $𝒱$-faisceau d'un
$𝒰$-faisceau est un $𝒰$-faisceau ce qui est bien
clair. Démontrons
4). Soit $M$ un $A$-Module plat. Il suffit de montrer que pour tout
$\epsilon(A)$-Module injectif $J$, le faisceau abélien
$\Hhom_{\epsilon(A)} (\epsilon(M), J)$ est injectif
(\hyperref[V.prop1.2]{1.2}). En vertu de \hyperref[V.prop0.5]{0.5}, $J$ est sous-objet, donc facteur direct d'un
objet du type $\prod_{\alpha} \epsilon(I_{\alpha})$, où les
$I_{\alpha}$ sont injectifs. On peut donc supposer que $J = \epsilon(I)$
où $I$ est un objet injectif. Si l'homomorphisme $\epsilon(
\Hhom_A (M, I)) \to \Hhom_{\epsilon \sisi{A}{(A)}}
(\epsilon\sisi{M}{(M)}, \epsilon \sisi{I}{(I)})$ est un isomorphisme on en déduit que
$\Hhom_{\epsilon(A)} (\epsilon(M), \epsilon(I))$
est injectif d'après 3). Il reste donc à démontrer 5). Le fait
que la formation du produit tensoriel commute au foncteur
$\epsilon$ résulte de IV 12.6 et du fait que la formation du
faisceau associé commute à l'agrandissement de l'univers (II
3.6). Pour tout objet $X$ de $C^{\sim}_{\sheaf U}$ et tout couple de $A$-Modules
$M$ et $N$ de $C^{\sim}_{\sheaf U}$ on a donc
$$
\Hom_{C^{\sim}_{\sheaf V}} (\epsilon(X), \sisi{}{\epsilon}\sisi{}{(}\Hhom_A (M, N)\sisi{}{)}) \simeq
\Hom_{C^{\sim}_{\sheaf U}} (X, \Hhom_A (M, N)) \simeq \Hom_A (\ZZ_X
\otimes_{\ZZ} M, N)
$$
en vertu de la pleine fidélité de $\epsilon$, puis
$$
\Hom_A (\ZZ_X \otimes_{\ZZ} M, N) \simeq \Hom_{\epsilon\sisi{A}{(A)}}(\epsilon
(\ZZ_X \otimes_{\ZZ} M), \epsilon \sisi{N}{(N)}) \simeq \Hom_{\epsilon \sisi{A}{(A)}}(\epsilon
\ZZ_X \otimes_{\ZZ} \epsilon \sisi{M}{(M)}, \epsilon \sisi{N}{(N)})
$$
en vertu de la pleine fidélité de $\epsilon$ et de ce qui
précède. Utilisant alors 2) et IV 12.14, on obtient en
définitive un isomorphisme
$$
\Hom_{C^{\sim}_{\sheaf V}} (\epsilon \sisi{X}{(X)}, \epsilon
\sisi{\Hhom_A (M,N)}{(\Hhom_A (M,N))}) \simeq \Hom_{C^{\sim}_{\sheaf V}} (\epsilon \sisi{X}{(X)},
\Hhom_{\epsilon \sisi{A}{(A)}} (\epsilon \sisi{M}{(M)}, \epsilon \sisi{N}{(N)}));
$$
d'où l'isomorphisme annoncé.

\setcounter{subsection}{9}
\subsection{}\etiquette[1.10]{V.subsec1.10}
Soient $E$ un topos, $\Uu = (U_i \to X)_{i \in I}$ une petite
famille de morphismes de même but. Pour tout ensemble fini
$\Delta_n = \sisi{[0, \ldots, n]}{\{0,\ldots,n\}}$, posons
$$
S_n(\Uu) = \coprod_{f : \sisi{\Delta}{\Delta_n} \to I} U_{f(1)} \times_{X} \ldots
\times_{X} U_{f(n)}{\quoi},
$$
la somme étant prise sur l'ensemble des applications de
$\Delta_n$ dans $I$. Soient $m$ et\pageoriginale $n$ deux entiers, $g :
\Delta_{m} \to \Delta_{n}$ une application. On définit un
morphisme
\begin{equation*}
s(g) : S_n (\Uu) \to S_m(\Uu){\quoi}\sisi{,}{}\tag{1.10.1}
\end{equation*}\etiquette[1.10.1]{V.1.10.1}
de la manière suivante: pour tout $f : \Delta_n \to I$, la
restriction de $s(g)$ à \mbox{$U_{f(1)} \times_{X}\ldots \times_X
U_{f(n)}$} est le morphisme composé de l'inclusion canonique
$$
U_{f(g(1))} \times_X \ldots \times_X U_{f(g(m))} \hookrightarrow S_m(\Uu){\quoi},
$$
et de l'unique morphisme
$$
s_f(g) : U_{f(1)} \times_X \ldots \times_X U_{f(n)} \to U_{f(g(1))}
\times_{X} \ldots
\times_X U_{f(g(m))}
$$
tel que pour tout $i \in \Delta_{m}$
$$
\pr_{f(g(i))} s_f (g) = \pr_{f(g(i))}
$$
($\pr_\alpha$ désigne la $\alpha$-\sisi{ème}{ième} projection). On obtient ainsi un foncteur
contravariant $\Delta_n \mapsto S_n$ de la catégorie des
ensembles finis dans $E$; autrement dit un complexe semi-simplicial
$S_∙ (\Uu)$ d'objets de $E$. Notons que ce complexe est
canoniquement augmenté vers $X$. Tout foncteur de $E$ dans une
catégorie $C$ transforme $S_∙ (\Uu)$ en un objet simplicial
de $C$. En particulier si $A$ est un Anneau de $E$, le foncteur «
$A$-Module libre engendré » transforme $S_∙ (\Uu)$ en un
complexe simplicial de $A$-biModules augmenté vers $A_X$ et noté
$A_∙ (\Uu)$. On a
\begin{equation*}
A_n (\Uu) = \bigoplus_{f:\Delta_n\to I} A_{U_{f(1)}
\times_X \ldots \times_X U_{f(\sisi{}{n})}{\quoi}}.\tag{1.10.2}
\end{equation*}\etiquette[1.10.2]{V.1.10.2}
Ce complexe sera souvent noté
\[
\xymatrix@C=1cm@M=.5cm{
  \cdots  \ar@<12pt>[r]^(.4){s_0} \ar[r]^(.4){s_1} \ar@<-12pt>[r]^(.4){s_2}&
\displaystyle{\mathop{\bigoplus}_{i,j}} \quad A_{U_i ×_X U_j}
 \ar@<6pt>[r]^(.55){s_0}\ar@<-6pt>[r]^(.55){s_1} &
\displaystyle{\mathop{\bigoplus}_{i}}\quad  A_{U_i} \ar@<2pt>[r]& A_X \\
}\tag{1.10.3} %% problème (numérotation au mauvais endroit) \XXX
\]\etiquette[1.10.3]{V.1.10.3}
où les flèches de \eqhyperref[V.1.10.3]{1.10.3} (sauf la dernière qui est
l'augmentation) représentent les opérateurs faces du complexe
simplicial $A_∙(\Uu)$, c'est-à-dire les morphismes \sisi{correspondants}{correspondant} aux
applications injectives croissantes de $\Delta_n$ dans $\Delta_{n +
1}$ ($s_i$ évite l'entier $i$). Au complexe $A_∙(\Uu)$, on associe un
complexe différentiel augmenté vers $A_X$:
\begin{equation}
\cdots \xrightarrow{d} \bigoplus_{i, j} A_{U_i \times_X U_j}
\xrightarrow{d} \bigoplus_{i} A_{U_i} \to A_X{\quoi},\sisi{\tag*{1.10.4}}{\tag{1.10.4}}
\end{equation}\etiquette[1.10.4]{V.1.10.4}
en \pageoriginale posant
\begin{equation}
d = \sum (-1)^i s_i.\sisi{\tag*}{\tag}{1.10.5}
\end{equation}\etiquette[1.10.5]{V.1.10.5}

\begin{enonce*}{Proposition 1.11}\etiquette[1.11]{V.prop1.11}
Lorsque la famille $\Uu = (U_i \to X)_{i \in I}$ est
épimorphique, le complexe différentiel \eqhyperref[V.1.10.4]{1.10.4} est acyclique et
fournit une résolution de $A_X$.
\end{enonce*}

Notons $\ZZ$ le faisceau constant de valeurs $\ZZ$. Par
définition du foncteur « $A$-Module libre engendré », on a, pour
tout objet $H$ de $E$,
$$
A_H \simeq {\ZZ}_H \otimes_{\ZZ} A {\quoi},
$$
d'où
$$
A_∙(\Uu)\simeq \ZZ_∙(\Uu) \otimes_{\ZZ} A{\quoi}.
$$
Comme les composantes de $\ZZ_∙(\Uu)$ sont des
$\ZZ$-Module plats (\hyperref[V.prop1.3.1]{1.3.1}), il suffit de montrer la proposition
lorsque $A = \ZZ$.

Supposons tout d'abord que $E$ soit le topos des ensembles. Alors le
complexe augmenté $S_∙(\Uu)$ est somme directe de complexes
augmentés du type
\[
\xymatrix@C=.6cm@M=.18cm@R=.3cm{
\cdots \ar@<9pt>[r] \ar@<3pt>[r]  \ar@<-3pt>[r] \ar@<-9pt>[r]&
S\times S \times S  \ar@<5pt>[r]\ar[r]\ar@<-5pt>[r] & S\times S
\ar@<2pt>[r]\ar@<-2pt>[r]&S\ar[r]& \{e\} & \txt{(\{e\} ensemble à un élément).}
}
\]
Chacun de ces complexes augmentés est homotopiquement trivial. Donc
$S_∙(\Uu)$ est un complexe augmenté homotopiquement trivial
et par suite son homologie est triviale d'où la proposition dans ce
cas.

Soit maintenant $p : \text{Ens} \to E$ un point de $E$. Comme la
formation du complexe $\ZZ_∙(\Uu)$ commute aux
foncteurs image inverse par les morphismes de topos,
$p^\ast(\ZZ_∙(\Uu)) \simeq \ZZ_∙ (p^\ast
(\Uu))$ est une résolution de ${\ZZ}_{p^{\ast}X} \simeq
p^{\ast}({\ZZ}_X)$; d'où la proposition lorsque $E$ possède
suffisamment de foncteurs fibres (IV 4.6). Ceci est le cas en
particulier lorsque $E$ est un topos de préfaisceaux
$C^{\wedge}$ car
pour tout objet $X$ de $C$, $\Gamma (X, -)$ est un foncteur
fibre. Dans le cas général, $E$ est équivalent à un topos de
faisceaux sur un petit site $C$ (IV 1) et la famille épimorphique
$\Uu = (U_i \to X)_{i \in I}$ est image, par le foncteur «
faisceau associé », d'une famille épimorphique $\Uu′ = (U'_i
\to X')_{i \in I}$. Par \pageoriginale suite
${\ZZ}_∙ (\Uu)\simeq \underline{a}
    {\ZZ}_∙ (\Uu')$ est une résolution de $\underline{a}
    {\ZZ}_{X'} \simeq {\ZZ}_X$.


\section{Cohomologie de Čech. Notation cohomologique}\etiquette[2]{V.sec2}

\subsection{Notation générale}\etiquette[2.1]{V.subsec2.1}

\subsubsection{}\etiquette[2.1.1]{V.subsubsec2.1.1}
Soient $(E, A)$ un topos annelé, $M$, $N$ deux $A$-Modules (à gauche
pour fixer les idées). On note $\Ext^q_A (E; M, N)$ la valeur en $N$
du $q$-ième foncteur dérivé droit du foncteur $\Hom_A(M,⋅)$
\cite{SGA4-05:11}. Les foncteurs $\Ext_A^q(E; M, N)$, $q \in \NN$, forment un
$\delta$-foncteur en la variable $N$. C'est aussi un
$\delta$-foncteur contravariant en la variable $M$. On a, par
définition, $\Ext_A^0(E; M, N) = \Hom_A(M, N)$.

\subsubsection{}\etiquette[2.1.2]{V.subsubsec2.1.2}

Soit $X$ un objet de $E$. Lorsque $M = A_X$, le $A$-Module libre
engendré par $X$ (IV 12), on pose $\Ext_A^q(E; A_X, N) = H^q(X,
N)$. On remarquera que dans cette nouvelle notation, l'anneau $A$ ne
figure plus. Ceci ne peut prêter à confusion car nous
montrerons (\hyperref[V.cor3.5]{3.5}) que la formation des $H^q(X,⋅)$ \sisi{commutent}{commute} à la
restriction des scalaires. Le foncteur $H^q(X,⋅)$ est le $q$-ième
foncteur dérivé droit du foncteur $\Hom_A(A_X,⋅) = \Hom_E(X,⋅)$
encore noté $\Gamma(X,⋅)$. Lorsque $M = A$, on pose $\Ext^q_A(E; A,
N) = H^q(E, N)$.

\subsection{Localisation}\etiquette[2.2]{V.subsec2.2}

Soient $X$ un objet de $E$, $j : E_{/X} \to E$ le morphisme de
localisation (IV 8). Le foncteur $j^\ast_X$ pour les Modules\sisi{,}{} est exact
(\hyperref[V.prop4.11]{4.11}) et admet un foncteur adjoint à gauche $j_{X!}$ exact
(\hyperref[V.prop4.11]{4.11}). Par suite il transforme les Modules injectifs en Modules
injectifs. On a donc, pour un $A$-Module variable $N$ de $E$ et un $A
| X$-Module $M$ variable de $E_{/X}$, des isomorphismes fonctoriels
\begin{equation*}
\Ext^q_{A | X}(E_{/X} ; M, j^\ast_X N) \simeq \Ext^q_A(E; j_{X!}M,
N){\quoi}.\tag{2.2.1}
\end{equation*}\etiquette[2.2.1]{V.2.2.1}
En particulier, on a des isomorphismes canoniques
$$
H^q(E_{/X}, j^\ast_X N) \simeq H^q (X, N){\quoi}.
$$

Pour tout objet $X$ de $E$ et tout couple $M$ et $N$ de $A$-Modules,
on pose:
\begin{equation*}
\Ext^q_A(X; M, N) \simeq \Ext^q_{A|X}(E_{/X}; M | X, N|X){\quoi}.\tag{2.2.2}
\end{equation*}\etiquette[2.2.2]{V.2.2.2}
D'après\pageoriginale ce qui précède, les foncteurs $\Ext^q_A(X; M,
⋅)$ sont les foncteurs dérivés des foncteurs $N \mapsto
\Hom_{A|X}(M|X, N|X)$. \sisi{Le foncteur}{Les foncteurs} $(M, N) \mapsto \Ext^q_A(X; M,
N)$, $q > 0$, forment un $\delta$-foncteur par rapport à chacune des
variables.

\subsection{Cas des topos de préfaisceaux. Cohomologie d'un
  recouvrement}\etiquette[2.3]{V.subsec2.3}

\subsubsection{}\etiquette[2.3.1]{V.subsubsec2.3.1}
Soient $C$ une petite catégorie munie d'un
préfaisceau d'anneaux $A$, $C^{\wedge}$ le topos
des préfaisceaux sur $C$, $X$ un objet représentable sur
$C^{\wedge}$. Le foncteur qui associe à un
$A$-Module $M$ le groupe $\Gamma(X, M) = M(X)$ est exact (I 3). Par
suite $H^q(X, M)=0$ pour tout $q > 0$ et tout $A$-Module $M$ ou encore
que $A_X$ est un $A$-Module projectif. %% \XXX syntaxe perfectible…

\subsubsection{}\etiquette[2.3.2]{V.subsubsec2.3.2}
Soit $S$ un préfaisceau sur $C$. On a un
isomorphisme canonique pour tout $A$-Module $M$ (I 2):
$$
\Gamma(S, M) \simeq \varprojlim_{C/S} M|S
$$
De plus, pour tout $A$-Module injectif $M$, le $A | S$-Module $M | S$
est injectif (2.2). Par suite, les groupes $H^{q}(S, M)$ sont les
valeurs en $M | S$ des foncteurs dérivés à droite du foncteur
$\varprojlim_{C/S}$. En notant $\varprojlim_{C/S}^{q}$ ces
foncteurs dérivés, on a des isomorphismes canoniques:
$$
H^q (S, M) \simeq {\varprojlim_{C/S}}^{q} M | S{\quoi}.
$$
En particulier, on a des isomorphismes canoniques
$$
H^q (C^{\wedge}, M) \simeq {\varprojlim_{C/S}}^{q} M{\quoi}. %% \XXX S en trop dans le terme de droite
$$

\subsubsection{}\etiquette[2.3.3]{V.subsubsec2.3.3}
Soient $X$ un objet de $C$ et $\Uu = (U_i \to X)$, $i \in I$,
une famille de morphismes de $C$ telle que pour tout $i \in I$, $U_i
\to X$ soit quarrable (I 10). Notons $A_∙$ le complexe simplicial
étudié en \hyperref[V.subsec1.10]{1.10}:
\[
\xymatrix@M=.25cm@C=1cm{
A_∙ = \cdots \ar@<7pt>[r] \ar@<1pt>[r]  \ar@<-5pt>[r] &
\displaystyle{\mathop{\coprod}_{(i,j)\in I \times I}}\quad
A_{U_i \times_X U_j} \ar@<3pt>[r]\ar@<-3pt>[r]&
\displaystyle{\mathop{\coprod}_{i\in I}} \quad A_{U_i}.
}
\]
% \XXX pourquoi ∐ et pas ⨁ comme plus haut ?
Pour tout $A$-Module $M$, on pose $C^{∙}(\Uu, M) = \Hom_A (A_∙,
M)$:
\[\vcenter{
\xymatrix@C=.75cm@M=.25cm{
C^{∙}(\mathcal{U},M): \displaystyle{\mathop{\prod}_{i\in I}}
M(U_i) \ar@<6pt>[r]\ar@<-2pt>[r]& \displaystyle{\mathop{\prod}_{(i,j)\in
    I\times I}} M(U_i ×_X U_j) \ar@<10pt>[r] \ar@<3pt>[r]
\ar@<-5pt>[r]& \cdots .
}}\leqno{(2.3.3.1)}
\]\etiquette[2.3.3.1]{V.2.3.3.1}
On pose $H^q(\Uu, M) = H^q(C^{∙} (\Uu, M))$.


\begin{enonce*}{Proposition 2.3.4}\etiquette[2.3.4]{V.prop2.3.4}
\mbox{ }
\begin{enumerate}
\renewcommand{\labelenumi}{\theenumi)}
\item Avec\pageoriginale les notations de \hyperref[V.subsubsec2.3.3]{2.3.3}, soit $R
  \hookrightarrow X$ le
  crible engendré par la famille $(U_i \to X)$, $i \in I$. On a un
  isomorphisme canonique
$$
H^{q}(\Uu, M) \simeq H^q(R, M){\quoi}.
$$

\item Les foncteurs $H^q(\Uu,⋅)$ commutent aux restrictions
  des scalaires.
\end{enumerate}
\end{enonce*}

Comme $R$ est un sous-objet de $X$ dans $C^{\wedge}$,
les produits fibrés $U_{i_1} \times_R U_{i_2} \sisi{\times}{\times_R} \cdots
\times_{R} U_{i_n}$ et $U_{i_1} \times_X U_{i_2} \sisi{\times}{\times_X} \cdots
\times_{X} U_{i_n}$ sont canoniquement isomorphes. Il résulte de
\hyperref[V.prop1.11]{1.11} que le complexe $A_{𝒰 ∙}$ est une résolution de $A_R$ et
de \hyperref[V.subsubsec2.3.1]{2.3.1} que les composants de $A_{𝒰 ∙}$ sont des Modules
projectifs. Les groupes de \sisi{cohomologies}{cohomologie} de $C^{∙}(\mathcal{U}, M)$ sont
donc canoniquement isomorphes à
$\Ext^{q}_A(C^{\wedge}; A_R, M)$, d'où
l'isomorphisme. L'assertion 2) résulte immédiatement de la
description de $C^{∙}(\mathcal{U}, M)$ \eqhyperref[V.2.3.3.1]{2.3.3.1}.

\begin{enonce*}{Corollaire 2.3.5}\etiquette[2.3.5]{V.cor2.3.5}
Soient $\mathcal{U} = (U_i \to X, i \in I)$ et
$\mathcal{U}' = (U'_j \to X, j \in J)$ deux familles de
morphismes de but $X$; Soient $\Phi = (\phi : I \to J, m_i : U_i \to
U'_{\phi(i)})$ et $\phi' = (\phi' : I \to J, m'_i : U_i \to
H'_{\phi(i)})$ deux morphismes (au-dessus de $X$) de $\mathcal{U}$ dans
  $\mathcal{U}'$. Les morphismes $\phi$ et $\phi'$ induisent les
  morphismes $\phi^{q}$ et ${\phi'}^{q} : \sisi{{\phi'}^{q} :}{} H^q(\mathcal{U},
  M) \to H^q(\mathcal{U}', M)$. Les
  morphismes $\phi^q$ et ${\phi'}^{q}$ sont égaux. En particulier si
  les familles $\mathcal{U}$ et $\mathcal{U}'$ sont équivalentes
  (\ie s'il existe un morphisme de $\mathcal{U}$ dans $\mathcal{U}'$
  et un morphisme de $\mathcal{U}'$ dans $\mathcal{U}$) les
  $A$-Modules $H^{q} (\mathcal{U}, M)$ et $H^q(\mathcal{U}', M)$ sont
  canoniquement isomorphes.
\end{enonce*}

\begin{enonce*}[remark]{Preuve}
En effet, soient $R$ et $R'$ les cribles de $X$ engendrés par les
familles $\mathcal{U}$ et $\mathcal{U}'$. Les morphismes $\phi$ et
$\phi'$ définissent un même morphisme de $R$ dans $R'$ et induisent
donc deux morphismes homotopes entre les résolutions projectives
$A_{\mathcal{U}∙}$ et $A_{\mathcal{U}'∙}$.
\end{enonce*}

\begin{enonce*}[remark]{Exercice 2.3.6}\etiquette[2.3.6]{V.exer2.3.6}
(Résolution standard)
\begin{enumerate}
\renewcommand{\theenumi}{\alph{enumi}}
\renewcommand{\labelenumi}{\theenumi)}
\item Soit $C$ une petite catégorie. Pour tout entier $n > 0$, on
  note $\mathrm{Fl}^{n}(C)$ l'ensemble des suites de morphismes de $C : (u_1,
  \ldots, u_n)$ telles que pour tout $i$, $0 < i < n$, le but de $u_i$
  soit égal à la source de $u_{i + 1}$ de sorte que les
  morphismes $u_i$ et $u_{i +1}$ sont\pageoriginale composables. Définir un
  ensemble semi-simplicial
\[
\xymatrix@M=.25cm@C=.75cm{
ES(C) = \ob \;C  & \mathrm{Fl}^{1}(C) \ar@<3pt>[l]\ar@<-3pt>[l]
 & \mathrm{Fl}^{2}(C) \ar@<6pt>[l] \ar[l] \ar@<-6pt>[l] & \mathrm{Fl}^{3}(C)\ar@<9pt>[l]
\ar@<3pt>[l] \ar@<-3pt>[l] \ar@<-9pt>[l] \quad \cdots ,
}
\]
dont les opérateurs faces $s_i : \mathrm{Fl}^n(C) \to \mathrm{Fl}^{n-1}(C)$ sont les
suivants:
\begin{align*}
s_{0}(u_1, \ldots, u_n) &= (u_2, \ldots, u_n),\\
s_{i}(u_1, \ldots, u_n) &= (u_1, \ldots, u_{i+1} u_i, \ldots, u_n), \hskip1em 0 %% hskip=bidouille : harmoniser \XXX
< i < n ,\\
s_{n}(u_1, \ldots, u_n) &= (u_1, \ldots, u_{n-1}).
\end{align*}

\item Montrer que lorsque $C$ possède un objet initial ou un objet
  final, le complexe $ES(C)$ est homotopiquement trivial.

\item Pour tout objet $X$ de $C$, on note ${}_{X\backslash}C$ la catégorie
  des flèches de source $X$. Définir un préfaisceau
  semi-simplicial $ES(C)$ dont la valeurs en tout objet $X$ de $C$ soit
  $ES({}_{X\backslash}C)$. On note $\ZZ_{ES(C)}$ le préfaisceau
  semi-simplicial
  abélien libre engendré par $ES(C)$. Il est muni d'une
  augmentation canonique $\ZZ_{ES(C)}\to \ZZ$ dans le préfaisceau
  constant de valeur $\ZZ$. Montrer que, en passant au complexe
  différentiel associé, le complexe $\ZZ_{ES(C)}$ est une
  résolution de $\ZZ$ et que les composants de ce complexe sont des
  préfaisceaux abéliens projectifs.

\item Pour tout préfaisceau abélien $M$, on pose $ST^{∙}(M) =
  \Hom(\ZZ_{ES(C)}, M)$. Expliciter les composants de
  $ST^{∙}(M)$. Montrer que pour tout entier $q \geq 0$,
  $H^q(ST^{∙}(M)) = H^q(C^{\wedge}, M)$.

\item Montrer que pour tout faisceau $S$ sur $C$ les foncteurs $H^q(S,⋅)$
 commutent aux restrictions des scalaires.

\item Remarquer que lorsque $C$ est un groupe, $\ZZ_{ES(C)}$ est la
  résolution standard du module trivial $\ZZ$, \cite{SGA4-05:3}.

\item Montrer que pour tout préfaisceau abélien $M$ sur $C$,
  $\ZZ_{ES(C)} \otimes M$ est une résolution (à gauche) de
  $M$. Définir le foncteur $\varinjlim_{C}$ sur les %% displaystyle ? \XXX
  préfaisceaux. Montrer qu'il est exact à droite. Noter $H_q(C,
  M)$ ses foncteurs dérivés à gauche. Montrer que les composants
  du complexe $\ZZ_{ES(C)} \otimes M$ sont acycliques pour
  $\varprojlim_{C}$. En déduire, en notant $ST_{∙}(M)$ le complexe
  $\varprojlim_{C}(\ZZ_{ES(C)} \otimes M)$, des isomorphismes
  $H_q(ST_{∙}(M) \sim H_q(C, M))$.

\item Soit\pageoriginale $u : C^{\wedge} \to \Ens$
  l'unique morphisme du topos $C^{\wedge}$ dans le topos
  ponctuel. On note $u_{!} : C^{\wedge}_{\ZZ} \to \Ab$
  l'adjoint à gauche du foncteur $u^\ast : \Ab \to
  C^{\wedge}_{\ZZ}$. Montrer que pour tout préfaisceau
  abélien $M$, $u_{!}M =\varinjlim_{C} M$.

\item On note dorénavant $ST^{∙}(C, M)$ et $ST_{∙}(C, M)$ les
  complexes notés $ST^{∙}(M)$ et $ST_{∙}(M)$. Un
  préfaisceau $M$ sur $C$ est dit localement constant s'il
  transforme tous  les morphismes de $C$ en isomorphismes. Associer
  à tout préfaisceau localement constant $M$ un préfaisceau
  localement constant $M$ sur la catégorie $C^{\circ}$ (catégorie
  opposée à $C$) tel que pour tout morphisme $u$ de $C$,
  $\check{M}(u) = M(u)^{-1}$. Trouver un isomorphisme canonique entre
  les complexes $ST^{∙}(C, M)$ et $ST^{∙}(C^{\circ},
  \check{M})$ (\resp $ST_{∙}(C, M)$ et $ST.(C^{\circ}, \check{M})$)

\item Soit $C$ une petite catégorie possédant un objet initial
  (\resp une petite catégorie filtrante). Montrer que pour tout
  préfaisceau constant $M$, $H^{q} (C^{\wedge}, M) =
  0$ pour $q > 0$ (\resp $H_q(C, M) = 0$ pour $q > 0$).
\end{enumerate}
\end{enonce*}

\subsection{Cas des petits sites, Cohomologie de Čech.}\etiquette[2.4]{V.subsec2.4}

\subsubsection{}\etiquette[2.4.1]{V.subsubsec2.4.1}
Soient $(C, A)$ un $𝒰$-site annelé, $C^{\wedge}$ le
topos des faisceaux sur $C$, $\epsilon: C \to
C^{\wedge}$ le foncteur canonique qui associe à un
objet de $C$ le faisceau associé au préfaisceau représenté par
cet objet. Par abus de notation, pour tout objet $X$ de $C$ et tout
faisceau de $A$-modules $M$ nous poserons $H^q(X, M) = H^q(X, M)$ % un epsilon manque probablement \XXX
(\cf 2.3.1). Rappelons que lorsque la topologie de $C$ est moins fine
que la topologie canonique, ce qui est toujours le cas dans la
pratique, le foncteur $\epsilon$ est pleinement fidèle et permet
d'identifier $C$ avec son image par $\epsilon$.

\subsubsection{}\etiquette[2.4.2]{V.subsubsec2.4.2}
On note $\Hh^0 : C^{\sim{\phantom{a}}}_A \to
C^{\wedge}_A$ le foncteur d'inclusion des faisceaux de
$A$-modules dans la catégorie des préfaisceaux de
$A$-modules. Pour tout faisceau de $A$-modules $M$ on a donc par
définition:
\begin{equation*}
\Hh^{0}(M)(X) = H^{0}(X, M) = M(X){\quoi},\tag{2.4.2.1}
\end{equation*}\etiquette[2.4.2.1]{V.2.4.2.1}
pour tout objet $X$ de $C$. Le foncteur $\Hh^{0}$ est
exact à gauche. Ses foncteurs dérivés à droite sont notés
$\Hh^{q}$. Comme pour tout objet $X$ de $C$, le foncteur «
section sur $X$ » est exact dans la catégorie des préfaisceaux,
on a
\begin{equation*}
\mathcal{H}^q(M)(X) = H^q(X, M){\quoi}, \tag{2.4.2.2}
\end{equation*}\etiquette[2.4.2.2]{V.2.4.2.2}
pour \pageoriginale tout objet $X$ de $C$ et tout faisceau de
$A$-modules $M$, de sorte que le préfaisceau $\mathcal{H}^q(M)$
n'est autre que le préfaisceau $X \mapsto H^q(X, M)$.

\subsubsection{}\etiquette[2.4.3]{V.subsubsec2.4.3}
On suppose que $(C, A)$ est un \emph{petit site} annelé de sorte que
$C^{\wedge}$ est un \emph{topos} auquel on peut
appliquer les résultats de \hyperref[V.subsec2.3]{2.3}. Soit $X$ un objet de $C$ et $R
\hookrightarrow X$ un crible couvrant. Pour tout \emph{préfaisceau}
de $A$-modules $G$, les groupes $H^{q}(R, G)$ (qui sont donc
calculés dans le topos $C^{\wedge}$) sont appelés
\emph{les groupes de cohomologie de Čech du préfaisceau
  $G$ relatifs au crible couvrant $R$.}
Lorsque $R \hookrightarrow X$ est le crible engendré par une famille
couvrante $\mathcal{U} = (U_i \to X)_{i \in I}$ de
morphismes quarrables ces groupes peuvent se calculer à l'aide du
complexe $C^{∙}(\mathcal{U}, G)$ (\hyperref[V.subsubsec2.3.3]{2.3.3}) appelé \emph{complexe de
  Čech de $G$ relatifs à la famille couvrante $\mathcal{U}$.} (ou
du recouvrement $\Uu$). Les groupes
$H^{q}(\Uu,G)=H^{q}(R,G)$ (\hyperref[V.prop2.3.4]{2.3.4}) sont alors appelés
\emph{groupes de cohomologie de Čech de $G$ relatifs à la famille
  couvrante $\Uu$}.

\subsubsection{}\etiquette[2.4.4]{V.subsubsec2.4.4}
Soit $M$ un faisceau de $A$-modules sur $C$. Les groupes
$H^q(\mathcal{U}$, $\Hh^{0}(M))$ sont le plus souvent
notés, abusivement, $H^q(\mathcal{U}, M)$ et appelés groupes de
cohomologie de \sisi{Cech}{Čech} du faisceau $M$ relatifs à la famille couvrante
$\mathcal{U}$.

\subsubsection{}\etiquette[2.4.5]{V.subsubsec2.4.5}
On note $\check{\Hh}^0 : C^{\wedge}_A \to
C^{\wedge}_A$ l'extension naturelle aux préfaisceaux
de $A$-modules, du foncteur $L$ décrit en II. On a donc, par
définition, pour un préfaisceau $G$ et un objet $X$ de $C$:
\begin{equation}
\check{\Hh}^0 (G) (X) = \varinjlim_{R \hookrightarrow X}
G(R){\quoi},\tag{2.4.5.1}
\end{equation}\etiquette[2.4.5.1]{V.2.4.5.1}
la limite inductive étant prise suivant les cribles couvrant $X$. Il
résulte de \eqhyperref[V.2.4.5.1]{2.4.5.1} que le foncteur $\check{\Hh}^{0}$
est exact à gauche. Les foncteurs dérivés à droite de
$\check{\Hh}^{0}$ sont notés
$\check{\Hh}^{q}$. Comme les foncteurs « section sur $X$ »
et « limite inductive filtrante » sont exacts, il résulte de
\eqhyperref[V.2.4.5.1]{2.4.5.1} qu'on a
\begin{equation}
\check{\Hh}^{q} (G) (X) = \varinjlim_{R \hookrightarrow
  X} H^q(R, G){\quoi},\tag{2.4.5.2}
\end{equation}\etiquette[2.4.5.2]{V.2.4.5.2}
la limite étant prise suivant les cribles couvrant $X$.

Les \pageoriginale préfaisceaux $\check{\Hh}^{q}$ sont appelés les
\emph{préfaisceaux de cohomologie de Čech}. Pour tout objet
  $X$ de $C$, on pose
\begin{equation*}
{\check{H}}^q (X, G) = \check{\Hh}^{q} (G) (X){\quoi}.\tag{2.4.5.3}
\end{equation*}\etiquette[2.4.5.3]{V.2.4.5.3}
Les groupes ${\check{H}}^q(X, G)$ sont appelés les \emph{groupes de
  cohomologie de Čech}. Lorsque la topologie de $C$ est
  définie par une prétopologie, ce qui est le plus souvent le cas
  dans la pratique, on a, compte tenu de \hyperref[V.prop2.3.4]{2.3.4},
\begin{equation*}
{\check{H}}^q (X, G) \simeq \varinjlim_{\mathcal{U}} {H}^{q}
(\mathcal{U}, G){\quoi},\tag{2.4.5.4}
\end{equation*}\etiquette[2.4.5.4]{V.2.4.5.4}
la limite inductive étant prise suivant les familles couvrantes
quarrables préordonnées par la relation d'ordre naturelle sur les
cribles qui leur correspondent (\hyperref[V.cor2.3.5]{2.3.5}).

\subsubsection{}\etiquette[2.4.6]{V.subsubsec2.4.6}
Soit $M$ un faisceau de $A$-modules. On pose, abusivement
\begin{equation}
{\check{H}}^q (X, M) = {\check{H}}^q (X, \Hh^{0}(M)),
\check{\Hh}^{q} (M) = \check{\Hh}^{q}(\check{\Hh}^{0} (M)){\quoi},\tag{2.4.6.1}
\end{equation}\etiquette[2.4.6.1]{V.2.4.6.1}
et les groupes ${\check{H}}^q(X, M)$ sont appelés groupes de
cohomologie de \sisi{Cech}{Čech} du faisceau $M$. Signalons que si les
foncteurs ${\check{H}}^q$ sont des foncteurs dérivés sur la
catégorie des préfaisceaux, ils ne forment pas, en général, un
$\delta$-foncteur sur la catégorie des faisceaux.

\subsection{Changement d'univers. Cohomologie de Čech dans le cas des $𝒰$-sites}\etiquette[2.5]{V.subsec2.5}

\subsubsection{}\etiquette[2.5.1]{V.subsubsec2.5.1}
Soient $(C, A)$ un $𝒰$-site annelé et $𝒱$ un univers
contenant $𝒰$. Le site $(C, A)$ est alors un
$𝒱$-site et on a un $𝒰$-topos
$C^{\sim}_{𝒰}$, un $𝒱$-topos
$C^{\sim}_{𝒱}$ et un foncteur canonique
d'inclusion $\epsilon : C^{\sim}_{𝒰}
\hookrightarrow C^{\sim}_{𝒱}$. Le
foncteur $\epsilon$ est exact et pleinement fidèle sur les
catégories de Modules et transforme les Modules injectifs en Modules
injectifs (\hyperref[V.prop1.9]{1.9}). Pour tout couple de $U$-faisceaux de
$A$-modules on a donc des isomorphismes canoniques
\begin{equation*}
\Ext^q_A(C^{\sim}_{𝒰}; M, N) \simeq
\Ext^q_{\epsilon A}(C^{\sim}_{𝒱}; \epsilon M,
\epsilon N),\quad q \geq 0{\quoi}.\tag{2.5.1.1}
\end{equation*}\etiquette[2.5.1.1]{V.2.5.1.1}
En particulier, pour tout $𝒰$-faisceau d'ensembles $R$ sur $C$, on a
des isomorphismes canoniques (\hyperref[V.subsubsec2.1.1]{2.1.1})
\begin{equation*}
H^q(R, M) \simeq H^q(\epsilon R, \epsilon M),\quad q \geq
0{\quoi},\tag{2.5.1.2}
\end{equation*}\etiquette[2.5.1.2]{V.2.5.1.2}
et \pageoriginale plus particulièrement encore, pour tout objet $X$
de $C$, on a des isomorphismes canoniques (\hyperref[V.subsubsec2.4.1]{2.4.1})
\begin{equation*}
H^q(X, M) \simeq H^q(X, \epsilon M){\quoi}.\tag{2.5.1.3}
\end{equation*}\etiquette[2.5.1.3]{V.2.5.1.3}
En termes vagues, on peut donc dire que la cohomologie des faisceaux
ne dépend pas du choix des univers et on peut toujours, pour la
nécessité d'une démonstration \sisi{où}{ou} d'une construction, augmenter
l'univers pour calculer la cohomologie d'un faisceau.

\subsubsection{}\etiquette[2.5.2]{V.subsubsec2.5.2}
Soient $(C, A)$ un $𝒰$-site annelé et $𝒱$ un univers
contenant $𝒰$. Notons $C^{\sim}_{A}$ % \sim plus joli que le \sim -> à uniformiser \XXX
la catégorie des $𝒰$-faisceaux $\widehat{\epsilon} :
C^{\wedge}_{A, 𝒰} \to
C^{\wedge}_{A, 𝒱}$ le foncteur d'inclusion
des $𝒰$-préfaisceaux dans les $𝒱$-préfaisceaux de
$A$-modules. Le foncteur $\widehat{\epsilon}$ est exact et par suite
les foncteurs dérivés du foncteur $\widehat{\epsilon} \Hh^{0} :
C^{\sim}_A \to C^{\wedge}_{A,
𝒱} (\phantom{ABCD})$ %\XXX référence à ajouter
sont les foncteurs $\widehat{\epsilon}
\Hh^q$, $q \geq 0$. Par abus de notation, nous noterons encore
$\Hh^q :C^{\sim}_A \to
C^{\wedge}_{A, 𝒱}$, $q \geq 0$, les foncteurs
$\widehat{\epsilon} \Hh^q$. Cet agrandissement de l'univers présente
lorsque $C$ est $𝒱$-petit l'avantage suivant: La catégorie
$C^{\wedge}_{𝒰}$ n'est pas en général un
$𝒰$-topos et les $𝒰$-préfaisceaux de $A$-modules ne
sont pas nécessairement des sous-modules de $𝒰$-préfaisceaux
injectifs, alors que la catégorie des $𝒱$-préfaisceaux est
un topos (un $𝒱$-topos) et que par suite tout
$𝒱$-préfaisceau de $A$-modules est un sous-objet d'un
$𝒱$-préfaisceau injectif (\hyperref[V.subsubsec0.1.1]{0.1.1}). Ainsi
pour tout $𝒱$-préfaisceau d'ensembles $R$ (et en particulier
lorsque $R$ est un $𝒰$-préfaisceau) et pour tout
$𝒰$-faisceau de $A$-modules $M$, les groupes $H^q(R, \Hh^q
(M))$ sont définis par \hyperref[V.subsubsec2.1.1]{2.1.1} et il résulte de
\eqhyperref[V.2.5.1.2]{2.5.1.2} que ces groupes ne dépendent pas de l'univers
$𝒱$ considéré. De même, pour tout couple d'entiers $\geq 0$
$p$ et $q$, les préfaisceaux $\check{\Hh}^p (\Hh^q (M))$ sont
définis par \hyperref[V.subsubsec2.4.5]{2.4.5} et il résulte de
\eqhyperref[V.2.5.1.2]{2.5.1.2} et de \eqhyperref[V.2.4.5.4]{2.4.5.4} que ces préfaisceaux ne
dépendent pas de l'univers $𝒱$ utilisé pour les définir.

\section{La suite spectrale de Cartan-Leray relative à un
  recouvrement}\etiquette[3]{V.sec3}\pageoriginale

\begin{enonce*}{Proposition 3.1}\etiquette[3.1]{V.prop3.1}
Soient $(C, A)$ un $𝒰$-site annelé, $𝒱$ un univers contenant
$𝒰$. Le foncteur $\Hh^{0} : C^{\sim}_A
\to C^{\wedge}_{A, 𝒱}$ (\hyperref[V.subsubsec2.5.2]{2.5.2}) transforme les
$A$-Modules injectifs en préfaisceaux injectifs. Pour tout entier $q
> 0$, et tout $A$-Module $M$, le faisceau associé au préfaisceau
$\Hh^q (M)$ est nul.
\end{enonce*}

Notons $\underline{a}_{𝒱}$ le foncteur « faisceau associé » pour les
$𝒱$-préfaisceaux et $\epsilon :C^{\sim}_A
\hookrightarrow C^{\sim}_{A, 𝒱}$ le foncteur
d'inclusion des $𝒰$-faisceaux dans les $𝒱$-faisceaux. On a
$\underline{a}_{𝒱}\Hh^{0} = \epsilon$ (II 3.6) et comme les
foncteurs $\underline{a}_{𝒱}$ et $\epsilon$ sont exacts, on a
$\underline{a}_{𝒱}
\Hh^q = 0$ pour tout entier $q > 0$. Pour tout $𝒰$-faisceau
$M$ et tout $𝒱$-préfaisceau $N$ on a un isomorphisme fonctoriel
$\Hom_{C^{\wedge}_{A, 𝒱}} (N, \Hh^{0} (M))
\simeq \Hom_{C^{\sim}_{A, 𝒱}}(\underline{a}_{𝒱}N, \epsilon M)$. Lorsque $M$ est injectif, $\epsilon M$ est
injectif (\hyperref[V.prop1.9]{1.9}) et comme le foncteur $\underline{a}_{𝒱}$ est exact, le foncteur
$N \mapsto \Hom_{C_{A, 𝒱}} (N, \Hh^{0} (M))$ est
exact. Par suite $\Hh^{0} (M)$ est injectif.

\begin{enonce*}{Théorème 3.2}\etiquette[3.2]{V.thm3.2}
Soient $(C, A)$ un $𝒰$-site annelé, $R$ un $𝒰$-préfaisceau
d'ensembles sur $C$, $M$ un faisceau de $A$-Modules. Il existe une
suite spectrale, fonctorielle en $R$ et en $M$:
\begin{equation*}
H^{p+q} (\underline{a} R, M) \Longleftarrow E^{p, q}_2 = H^p (R, \Hh^q
(M)){\quoi}.
\tag{3.2.1}
\end{equation*}\etiquette[3.2.1]{V.3.2.1}
(Lorsque $C$ n'est pas $𝒰$-petit, le terme $H^p(R, \Hh^q
(M))$ doit être interprété comme la cohomologie du préfaisceau
$\Hh^q(M)$ dans le topos $C^{\wedge}_{𝒱}$ où $𝒱$
est un univers contenant $𝒰$ tel que $C$ soit
$𝒱$-petit (\hyperref[V.subsubsec2.5.2]{2.5.2})).
\end{enonce*}

Par définition du foncteur « faisceau associé » (\cf II) on a un
isomorphisme de foncteur $H^{0}(\underline{a} R, M) \simeq H^0(R,
\Hh^{0}(M))$. Le foncteur $\Hh^{0}$ transforme
les objets injectifs en objets injectifs. La suite spectrale des
foncteurs composés (\hyperref[V.prop0.3]{0.3}) est la suite spectrale cherchée.

\begin{enonce*}{Corollaire 3.3}\etiquette[3.3]{V.cor3.3}
Soient $X$ un objet de $C$ et $\mathcal{U} = (U_i \to X)$, $i \in I$,
une famille couvrante telle que pour tout $i \in I$, le morphisme $U_i
\to X$ soit quarrable. On a alors une suite spectrale (dite de
Cartan-Leray):
\begin{equation}
H^{p+q}(X, M) \Longleftarrow E^{p, q}_2 = H^p(\mathcal{U},
\Hh^q (M)){\quoi}.\tag{3.3.1}
\end{equation}\etiquette[3.3.1]{V.3.3.1}
\end{enonce*}

Soit \pageoriginale $R \hookrightarrow X$ le crible engendré par
$\mathcal{U}$. Comme ce crible est couvrant, le faisceau associé à
$R$ est le faisceau associé à $X$ (II 5.2). Compte tenu de \hyperref[V.subsubsec2.4.1]{2.4.1},
on a $H^{p + q}(\underline{a} R, M) = H^{p + q}(X, M)$. Le corollaire résulte
alors de \hyperref[V.prop2.3.4]{2.3.4}.

\begin{enonce*}{Corollaire 3.4}\etiquette[3.4]{V.cor3.4}
Il existe une suite spectrale fonctorielle en le faisceau $M$ et
l'objet $X$ de $C$
\begin{equation}
H^{p + q}(X, M) \Longleftarrow E^{p, q}_2 = \check{H}^p(X,
\Hh^q (M)){\quoi}.\tag{3.4.1}
\end{equation}\etiquette[3.4.1]{V.3.4.1}
Cette suite spectrale fournit, lorsque $X$ varie dans $C$, la suite
spectrale de préfaisceaux:
\begin{equation}
\Hh^{p + q} (M)\Longleftarrow E^{p,q}_2 =
\check{\Hh}^p(\Hh^q(M)){\quoi}.
\tag{3.4.2}
\end{equation}\etiquette[3.4.2]{V.3.4.2}
Ces suites spectrales fournissent des morphismes fonctoriels
(morphismes de coin):
\begin{align*}%\etiquette[3.4.3]{V.3.4.3}
\Phi^q(M) &: \Hh^q(M) \to \Hh^q(M){\quoi},\tag{3.4.3}\\
\Phi^q_X(M) &: H^q (X, M) \to H^q (X, M){\quoi}.\tag{3.4.4}
\end{align*}
 Les morphismes $\Phi^q$ et $\Phi^q_X$ sont des isomorphismes lorsque
$q$ égale $0$ ou $1$. Les morphismes $\Phi^2$ et $\Phi^2_X$ sont des
monomorphismes. Plus généralement, lorsque les préfaisceaux
$\Hh^i(M)$ sont nuls pour $0 < i < n$, les morphismes $\Phi^q(M)$
et $\Phi^q_X(M)$ sont des isomorphismes pour $0 \leq q \leq n$ et des
monomorphismes pour $q = n + 1$.
\end{enonce*}

La première suite spectrale s'obtient en passant à la limite
inductive dans la suite spectrale \eqhyperref[V.3.2.1]{3.2.1} sur les cribles $R
\hookrightarrow X$ couvrant $X$ \eqhyperref[V.2.4.5.2]{2.4.5.2}. La deuxième suite
spectrale s'en déduit aussitôt \eqhyperref[V.2.4.5.3]{2.4.5.3}. Les faisceaux
associés aux préfaisceaux $\Hh^q (M)$ sont nuls lorsque $q > 0$
(\hyperref[V.prop3.1]{3.1}). On a donc $\check{\Hh}^{0}
\check{\Hh}^{0} \check{\Hh}^{q} (M) = 0$ pour tout $q > 0$
(II.3.4). Par suite $\check{\Hh}^{0} \Hh^{q} (M)$ est un préfaisceau séparé dont le faisceau associé est nul. On a donc $\check{\Hh}^{0} \Hh^{q} (M) = 0$
pour $q > 0$ (II 3.2). Les assertions sur les morphismes $\Phi^q$ et
$\Phi^q_X$ s'en déduisent aussitôt.

\begin{enonce*}{Corollaire 3.5}\etiquette[3.5]{V.cor3.5}
Soient $(E, A)$ un topos annelé et $M$ un $A$-Module (à gauche pour
fixer les idées). Notons $ℳ$
le Groupe abélien
sous-jacent à $M$. Le foncteur $M \mapsto ℳ$ est exact et par
suite, pour tout objet $X$ de $E$, l'isomorphisme canonique
$$
H^{0} (X, M) \xrightarrow{\sim} H^{0} (X, ℳ)
$$
se\pageoriginale prolonge en des morphismes
$$
H^q(X, M) \to H^q(X, ℳ)\qquad q \geq 0{\quoi}.
$$
Ces morphismes sont des isomorphismes.
\end{enonce*}

Pour tout objet $Y$ de $E$, on a \eqhyperref[V.2.4.5.4]{2.4.5.4}:
$$
H^p(Y, M) = \varinjlim_{\mathcal{U}} H^p(\mathcal{U}, M) \quav
H^p(Y, ℳ) = \varinjlim_{\mathcal{U}} H^p(\mathcal{U},
ℳ)
$$
la limite inductive étant prise sur les familles épimorphiques
$\mathcal{U} = (Y_i \to Y)$, $i \in I$.

Par suite, l'homomorphisme canonique $\check{H}^p(Y, M) \to
\check{H}^p(Y, ℳ)$ est un isomorphisme (\hyperref[V.prop2.3.4]{2.3.4}). Donc \eqhyperref[V.2.4.5.3]{2.4.5.3},
l'homomorphisme canonique $\check{\Hh}^p (M) \to
\check{\Hh}^p (ℳ)$ est un isomorphisme. Si $M$ est un
$A$-Module injectif, on a $\check{\Hh}^p (ℳ) = 0$ pour $p >
0$; d'où $\Hh^1 (ℳ) = 0$ et par récurrence sur $p$,
$\Hh^p (ℳ) = 0$, $p > 0$ (\hyperref[V.cor3.4]{3.4}). Donc $H^p(X,
ℳ) = 0$ pour
$p > 0$ \eqhyperref[V.2.4.2.2]{2.4.2.2} et \hyperref[V.subsec2.5]{2.5}. Par suite, le foncteur $M \mapsto ℳ$
transforme les objets injectifs en objets acycliques pour le foncteur
$H^{0} (X, ⋅)$, d'où l'isomorphisme annoncé.

\begin{enonce*}[remark]{Exercice 3.6}\etiquette[3.6]{V.exer3.6}
Soient $G$ un groupe topologique et $B_G$ son topos classifiant (IV
2.5). Notons $E_G$ l'objet de $B_G$ constitué par l'espace
topologique sous-jacent à  $G$ muni de l'opération de translation
à gauche par les éléments de $G$. Le morphisme canonique de
$E_G$ dans l'objet final $e_G$ de $B_G$ est un épimorphisme; d'où
un recouvrement $\mathcal{U} = (E_G \to e_G)$ et, pour tout faisceau
abélien $F$ de $B_G$, une suite spectrale
\begin{equation*}
E^{p, q}_2 = H^p(\mathcal{U}, \Hh^q (F)) \Longrightarrow H^{p
  +q}(B_G, F){\quoi}, \tag{3.6.1}
\end{equation*}\etiquette[3.6.1]{V.3.6.1}
qu'on se propose d'étudier.
\end{enonce*}

\begin{enumerate}
\renewcommand{\labelenumi}{\theenumi)}
\item Montrer que pour tout entier $n$, le topos
$$
B_{G/ E_G \times E_G \times \ldots \times E_G}\quad (n\quad \mbox{facteurs})
$$
est canoniquement équivalent au topos (IV 2.5)
$$
\TOP (G \times G \times G \times \ldots \times G)\quad (n - 1\quad
\mbox{facteurs})
$$
Pour\pageoriginale tout entier $n$, notons $F_n$ le faisceau sur l'espace
topologique $G \times \ldots \times G$ ($n$ facteurs) induit par $F$.

\item En déduire que le terme $E^{p, q}_2$ de \eqhyperref[V.3.6.1]{3.6.1} est
  canoniquement isomorphe au $p$-ème groupe de cohomologie d'un
  complexe du type
$$
H^{q}(e, F_{0}) \to H^{q}(G, F_1) \to H^{q}(G \times G, F_2) \to
\cdots {\quoi},
$$
qu'on explicitera ($e$ désigne l'espace topologique réduit à un
point).

\item En déduire que la cohomologie du topos classifiant $B_G$ à
  valeur dans les faisceaux localement constants est isomorphe à la
  cohomologie singulière correspondante des espaces classifiants utilisés en topologie lorsque pour les espaces topologiques du type $G \times G \times
  \ldots \times G$, la cohomologie des faisceaux localement constants
  coïncide avec la cohomologie singulière correspondante (Ce qui
  est le cas lorsque, par exemple, $G$ est localement contractile);
  (Pour les espaces classifiants, on pourra consulter \cite{SGA4-05:15}).
% problème de ponctuation, dont parenthésage... \XXX
\end{enumerate}

\section{Faisceaux acycliques}\etiquette[4]{V.sec4}

\begin{enonce*}{Définition 4.1}\etiquette[4.1]{V.def4.1}
\emph{Soient $(E, A)$ un topos annelé, $F$ un $A$-Module, $S$ une famille
topologiquement génératrice de $E$. On dit que $F$ est
$S$-acyclique si pour tout objet $X$ de $S$, et tout entier $q > 0$,
on a $H^q(X, F) = 0$. Lorsque $S$ est égal à $\ob E$, les faisceaux
$S$-acycliques sont appelés les} faisceaux flasques.
\end{enonce*}

\setcounter{subsection}{1}
\subsection{}\etiquette[4.2]{V.subsec4.2}
Soient $(C, A)$ un $𝒰$-site annelé, $F$ un faisceau de $A$-modules
sur $C$. On dit que $F$ est $C$-\emph{acyclique} si $F$ est $S$-acyclique
où $S$ est la famille des faisceaux associés aux objets de $C$.

\begin{enonce*}{Proposition 4.3}\etiquette[4.3]{V.prop4.3}
Soient $(C, A)$ un $𝒰$-site annelé, $F$ un faisceau de
$A$-modules. Notons $\mathcal{H}^{0} :
C^{\sim}_{A} \to C^{\wedge}_{A}$
le foncteur d'inclusion des $A$-Modules dans la catégorie des
préfaisceaux de $A$-modules (il s'agit ici des
$𝒱$-préfaisceaux où $𝒱$ est un univers tel que
$C$ soit $𝒱$-petit). Les conditions suivantes sont
équivalentes:
\begin{enumerate}
\renewcommand{\theenumi}{\roman{enumi}}
\renewcommand{\labelenumi}{\theenumi)}
\item $F$ \pageoriginale est $C$-acyclique.

\item Pour tout objet $X$ de $C$ et tout crible couvrant $R
  \hookrightarrow X$, on a:
$$
H^q(R, \Hh^{0} (F)) = 0 \mbox{ pour tout } q > 0{\quoi}.
$$

\item Pour tout objet $X$ de $C$, on a:
$$
\check{H}^q (X, F) = 0 \mbox{ pour tout } q > 0{\quoi};
$$
\end{enumerate}
\end{enonce*}

\begin{itemize}
\item[i $\Rightarrow$ ii)]: Comme $F$ est $C$-acyclique, les préfaisceaux
$\Hh^q(F)$ sont nuls pour $q > 0$. La suite spectrale \eqhyperref[V.3.2.1]{3.2.1}
fournit alors un isomorphisme $H^q (R, \Hh^{0} (F)) \simeq
H^q(X, F)$ pour tout $q$. Par suite $H^q(R, \Hh^{0} (F)) =
0$ pour $q > 0$.
\item[ii $\Rightarrow$ iii)]: clair par passage à la limite
  inductive.
\item[iii $\Rightarrow$ i)]: on démontre par récurrence sur $q$
  que les préfaisceaux $\Hh^q (F)$ sont nuls pour $q >
  0$. Pour cela, on utilise \eqhyperref[V.3.4.2]{3.4.2}.
\end{itemize}


\setcounter{subsection}{3}
\subsection{}\etiquette[4.4]{V.subsec4.4}
Il résulte du critère \hyperref[V.prop4.3]{4.3} ii) et de \hyperref[V.cor3.5]{3.5} que la propriété de
$S$-acyclicité ne dépend que du faisceau abélien sous-jacent. En
particulier, un faisceau de $A$-modules est flasque si et seulement si
le faisceau abélien sous-jacent est flasque.

\begin{enonce*}{Corollaire 4.5}\etiquette[4.5]{V.cor4.5}
Soient $(E, A)$ un topos annelé et $F$ un $A$-Module. Les
propriétés suivantes sont équivalentes:
\begin{enumerate}
\renewcommand{\theenumi}{\roman{enumi}}
\renewcommand{\labelenumi}{\theenumi)}
\item $F$ est flasque;

\item Pour toute famille épimorphique $𝒳 = (X_i \to X)_{i
\in I}$,
$$
H^q (𝒳 , F) = 0 \mbox{ pour tout } q > 0{\quoi}.
$$
\end{enumerate}
\end{enonce*}

Dans 4.2, on prend pour $C$ le topos $E$ lui-même. Le corollaire
résulte alors de l'équivalence i) $\Leftrightarrow$ ii) de
\hyperref[V.subsec4.2]{4.2}.

\setcounter{subsection}{5}
\subsection{}\etiquette[4.6]{V.subsec4.6}
Les faisceaux injectifs sont flasques. Les faisceaux flasques sont
$S$-acycliques pour toute famille topologiquement génératrice $S$.
Les faisceaux flasques ne sont pas nécessairement injectifs (prendre
pour topos $E$ le topos des ensembles). Les faisceaux $S$-acycliques
ne sont pas nécessairement flasques (exercice 4.1.3).

\begin{enonce*}{Proposition 4.7}\etiquette[4.7]{V.prop4.7}
Soient $(E, A)$ un topos, $F$ un $A$-Module flasque, $X$ un objet de
$E$. Pour tout sous-objet $Y$ de $X$ l'homomorphisme canonique
$H^{0} (X, F) \rightarrow H^{0}(Y, F)$ est \pageoriginale surjectif.
\end{enonce*}

Soit $Y$ un sous-objet de $X$ tel que le morphisme $H^{0}(X, F)
\to H^{0}(Y, F)$ ne soit pas surjectif. Soit $Z$ l'objet obtenu
en recollant deux copies $X_1$ et $X_2$ de $X$ le long de $Y$.
L'objet $Z$ est recouvert par les deux sous-objets $X_1$ et $X_2$ on
a $X_1 ×_Z X_2 = Y$. Notons $𝒳  = (X_1, X_2)$ le
recouvrement ainsi obtenu. On constate aussitôt que
$H^1(𝒳 , F) \neq 0$. Contradiction.

\setcounter{subsection}{7}
\subsection{}\etiquette[4.8]{V.subsec4.8}
La propriété décrite dans \hyperref[V.prop4.7]{4.7} ne caractérise pas, dans le cas
des topos généraux, les faisceaux flasques (exer. \hyperref[V.exer4.15]{4.15}). Elle la
caractérise cependant dans le cas des topos engendrés par leurs
ouverts et en particulier dans le cas des topos associés aux espaces
topologiques (exer. \hyperref[V.exer4.16]{4.16}). La terminologie de flasque adoptée ici
coïncide dans le cas des espaces topologiques avec la terminologie
de \cite{SGA4-05:7} (exer. \hyperref[V.exer4.16]{4.16}).

\begin{enonce*}{Proposition 4.9}\etiquette[4.9]{V.prop4.9}
Soient $u : (E, A) \to (E', A')$ un morphisme de topos annelés.
\begin{enumerate}
\renewcommand{\labelenumi}{\theenumi)}
\item Le foncteur $u_\ast$ transforme les $A$-Modules flasques en
  $A'$-Modules flasques.

\item Soient $S$ et $S'$ des familles topologiquement génératrices
  de $E$ et $E'$ respectivement telles que $u^\ast(S') \subset S$. Le
  foncteur $u_\ast$ transforme les $A$-Modules $S$-acycliques en
  $A'$-Modules $S'$-acycliques.

\item Lorsque $u$ est un morphisme plat (\hyperref[V.def1.8]{1.8}) le foncteur $u_\ast$
  transforme les $A$-Modules injectifs en $A'$-Modules injectifs.
\end{enumerate}
\end{enonce*}

Soient $X$ un objet de $E'$, $𝒳  = (X_i \to X)_{i \in I}$ une
famille épimorphique, $F$ un $A$-Module flasque, $C^{∙}(𝒳 ,
u_\ast F)$ le complexe de Čech du recouvrement
$𝒳 $. On a un isomorphisme canonique
$C^{∙}(𝒳 , u_\ast
F) \simeq C^{∙}(u^\ast (𝒳 ), F)$ en utilisant l'adjonction de
$u_\ast$ et de $u^\ast$ et le fait que $u^\ast$ commute aux produits
fibrés. De plus $u^\ast$ commute aux limites inductives et par suite
$u^\ast(𝒳 ) = (u^\ast X_i \to u^\ast X)_{i \in I}$ est une famille
épimorphique. Comme $F$ est flasque, on a $H^q(u^\ast(𝒳 ),
F)$ pour $q > 0$ et par suite $H^q(𝒳 , u_\ast F) = H^q(C^{∙}
(𝒳 , u_\ast F)) = H^q(C^{∙} (u^\ast(𝒳 , F)) = H^q
(u^\ast(𝒳 ), F)) = 0$ pour $q > 0$. Donc $u_\ast F$ est flasque.

Démontrons 2). Lorsque $S'$ est stable par produits fibrés une
démonstration analogue à celle qui précède permet de démontrer 2).
Dans le cas général, nous utiliserons la suite spectrale du
morphisme $u$ \eqhyperref[V.5.3.2]{5.3.2}). L'assertion 2 ne sera pas utilisée
avant\pageoriginale \eqhyperref[V.5.3.2]{5.3.2}. Soit $F$ un faisceau
$S$-acyclique. Les faisceaux $R^q u_\ast F$ sont les faisceaux
associés aux préfaisceaux $X \mapsto H^q(u^\ast X, F)$. Comme $S'$
est une famille topologiquement génératrice et comme $F$ est
$S$-acyclique, on a $R^q u_\ast F = 0$ pour tout $q > 0$. La suite
spectrale \eqhyperref[V.5.3.2]{5.3.2} fournit alors un isomorphisme, pour tout
objet $X$ de $E' : H^q (X, u_\ast F) \simeq H^q (u^\ast X, F)$ et
par suite, pour tout $X$ dans $S'$, $H^q (X, u_\ast F) = 0$.

Démontrons 3). Lorsque le morphisme $u$ est plat, le foncteur $u^\ast$
pour les Modules est exact (\hyperref[V.def1.8]{1.8}). Par suite le foncteur $u_\ast$ adjoint
à droite de $u^\ast$ transforme les objets injectifs en objets
injectifs (\hyperref[V.prop0.2]{0.2}).

\begin{enonce*}{Proposition 4.10}\etiquette[4.10]{V.prop4.10}
Soient $F$ un $A$-Module du topos annelé $(E, A)$, $G$ un
$A$-Module injectif.
\begin{enumerate}
\renewcommand{\labelenumi}{\theenumi)}
\item Le foncteur $F \mapsto \Hhom_A (F, G)$ est exact.

\item Le groupe abélien $\Hhom_A (F, G)$ est flasque.
\end{enumerate}
\end{enonce*}

\begin{enonce*}[remark]{Preuve}
Montrons
\begin{enumerate}
\renewcommand{\labelenumi}{\theenumi)}
\item Soit :
$$
0 \to F' \to F \to F'' \to 0
$$
une suite exacte, Il nous faut montrer que la suite :
$$
0 \to \Hhom_A (F'', G) \to \Hhom_A (F, G) \to
\Hhom_A (F', G) \to 0
$$
est exacte et pour cela il suffit de montrer que pour tout objet $H$ de
$E$, la suite:
\begin{align*}
0 \to \Hom_E(H, \Hhom_A (F'', G))  \to \Hom_E (H,
\Hhom_A (F , G))\hspace{3cm} \\
\to \Hom_E(H, \Hhom_A (F', G)) \to 0
\end{align*}
est exacte. Or (IV 6.12) cette dernière suite est isomorphe à la
suite
$$
0 \to \Hom_A (A_H \otimes_A F'', G) \to \Hom_A (A_H \otimes_A F, G)
\to \Hom_A (A_H \otimes_A F', G) \to 0
$$
et le $A$-Module $A_H$ est plat (\hyperref[V.prop1.3.1]{1.3.1}).

\item Montrons que $\Hhom_A(F, G)$ est flasque. Soit:
$$
𝒳  = (X_i \to X), i \in I{\quoi},
$$
une famille épimorphique de $E$, et
\[
\xymatrix@M=.25cm@C=1cm{
C_∙(𝒳 ) : \cdots \ar@<5pt>[r] \ar[r] \ar@<-5pt>[r] &
\displaystyle{\mathop{\coprod}_{i,j}}\ZZ_{X_{i}\times X_j}
\ar@<3pt>[r]\ar@<-3pt>[r]& \displaystyle{\mathop{\coprod}_{i}} \ZZ_{X_i}
}
\]
le\pageoriginale complexe défini en \hyperref[V.prop1.5]{1.5}. Ce complexe est une
résolution plate de
l'objet $\ZZ_X$ qui est lui-même plat (\hyperref[V.subsec1.4]{1.4} et
1.3).

\noindent
Calculons alors les groupes:
$$
H^q(𝒳 , \Hhom_A(F, G)) \xrightarrow{\sim}
H^q(\Hom_{\ZZ}^{∙}(C_∙ (𝒳 ), \Hhom_A(F, G))
\xrightarrow{\sim} H^q(\Hom_A (C_∙ (𝒳 ) \otimes_{\ZZ} F, G)){\quoi}.
$$
Le complexe
$$
C_∙ (𝒳 ) \otimes_{\ZZ} F
$$
est acyclique en degré $\neq 0$, car $C_∙ (𝒳 )$ est une
résolution plate d'un module plat. Par suite $\Hom(C_∙ (𝒳 )
\otimes_{\ZZ} F, G)$ est acyclique en degré $\neq 0$.
\end{enumerate}
\end{enonce*}

\begin{enonce*}{Proposition 4.11}\etiquette[4.11]{V.prop4.11}
Soient $(E, A)$ un topos annelé, $F$ un faisceau flasque (\resp
injectif) de $A$-modules.
\begin{enumerate}
\renewcommand{\labelenumi}{\theenumi)}
\item Pour tout objet $X$ de $E$ le $A | X$-Module $j_X^\ast F$ est
  flasque (\resp injectif)

\item Pour tout fermé $Z$ de $E$, le faisceau des sections de $F$
  à supports dans $Z$ \emph{(IV 14)} est flasque (\resp injectif).
\end{enumerate}
\end{enonce*}

L'assertion 1) résulte de \hyperref[V.cor4.5]{4.5} lorsque $F$ est flasque. Le foncteur
$j^\ast_X$ admet un adjoint à gauche exact $j_{X!}$ (IV 11). Par suite,
lorsque $F$ est injectif, $j^\ast_X F$ est injectif (\hyperref[V.prop0.2]{0.2}). Démontrons
2). Notons $i : Z \to E$ le morphisme d'inclusion. Le faisceau des
sections de $F$ à support dans $Z$ est alors le faisceau $i_\ast i^{!} F$ (IV 14). Le foncteur $i_\ast i^{!}$ est
adjoint à droite au foncteur $i_\ast i^\ast$ qui est exact (IV 14). Par
suite il transforme faisceau injectif en faisceau injectif. Soient $U$
l'ouvert complémentaire de $Z$, $j : U \to E$ le morphisme
d'inclusion et $F$ un faisceau flasque. On a une suite exacte (IV 14):
\begin{equation*}
0 \to i_\ast i^{!} F \to F \to j_\ast j^\ast F{\quoi}.\tag{4.11.1}
\end{equation*}\etiquette[4.11.1]{V.4.11.1}
Pour tout objet $X$ de $E$, on a $j_\ast j^\ast F(X) = F(X \times U)$ et le
morphisme $F(X) \to j_\ast j^\ast F(X)$ induit par la dernière flèche de
\eqhyperref[V.4.11.1]{4.11.1} provient de l'injection canonique $X \times U \hookrightarrow
X$. Comme $F$ est flasque, ce morphisme est surjectif (\hyperref[V.prop4.7]{4.7}) et par
suite, la dernière flèche de \eqhyperref[V.4.11.1]{4.11.1} est un épimorphisme de
préfaisceaux. Pour tout objet $X$ de $E$, la suite exacte de
cohomologie déduite de \eqhyperref[V.4.11.1]{4.11.1} fournit $H^q(X, i_\ast i^{!}F) = 0$ pour
$q > 0$ et par suite $i_\ast i^{!}F$ est flasque.

\setcounter{subsection}{11}
\subsection{}\etiquette[4.12]{V.subsec4.12}
La\pageoriginale propriété pour un faisceau d'être flasque (\resp
injectif) se localise \eqhyperref[V.4.11.1]{4.11.1}. Mais ce n'est pas, en
général, une propriété de caractère local (exer. \hyperref[V.exer4.15]{4.15}).
Cependant, c'est une propriété de caractère local dans le cas des
topos engendrés par leurs ouverts et en particulier dans le cas des
topos associés aux espaces topologiques (exer. \hyperref[V.exer4.16]{4.16}).

\begin{enonce*}[remark]{Exercice 4.13}\etiquette[4.13]{V.exer4.13}
Soient $X$ un espace localement compact et $F$ un faisceau $c$-mou
\cite{SGA4-05:7}. En utilisant le caractère local de la mollesse (\loccit),
montrer que la restriction de $F$ à tout ouvert paracompact de $X$
est un faisceau mou. Montrer que les ouverts paracompacts forment
une base de la topologie de $X$. En déduire que, en notant $S$ la
famille des ouverts paracompacts de $X$, le faisceau $F$ est
$S$-acyclique. Montrer que le faisceau des fonctions continues sur
$X$ est $S$-acyclique mais n'est pas flasque lorsque $X$ n'est pas
discret.
\end{enonce*}

\begin{enonce*}[remark]{Problème 4.14}\etiquette[4.14]{V.prob4.14}
Étudier les topos totalement acycliques, i.e. les topos tels que
$H^q (X, F) = 0$ pour tout $q > 0$, tout objet $X$, tout faisceau
abélien $F$.
\end{enonce*}

\begin{enonce*}[remark]{Exercice 4.15}\etiquette[4.15]{V.exer4.15}
Soient $G$ un groupe discret, $B_G$ son topos classifiant (IV
2.4). Montrer que pour tout faisceau abélien $F$ et tout
monomorphisme $X \hookrightarrow Y$, l'homomorphisme $F(Y) \to F(X)$
est surjectif. Soit $E(G)$ le groupe $G$ considéré comme espace
homogène sous lui-même. C'est un objet de $B_G$. Le topos
$B_{G/E(G)}$ est équivalent au topos ponctuel (IV 8). Le morphisme
$E(G) \to e$ ($e$ objet final de $B_G$) est un épimorphisme. Pour
tout faisceau abélien $F$ de $B_G$, le faisceau $F | E(G)$ est
flasque. En déduire que la propriété d'être flasque ou
injectif n'est pas de caractère local.
\end{enonce*}

\begin{enonce*}[remark]{Exercice 4.16}\etiquette[4.16]{V.exer4.16}
On dit qu'un topos $E$ est engendré par ses ouverts si les ouverts
de $E$ (\ie les sous-objets de l'objet final $e$ de $E$) forment une
famille génératrice (I 7). Un tel topos possède la propriété
suivante:

\noindent
(P) Toute famille épimorphique $X_i \to X$, $i \in I$, est
majorée par une famille épimorphique $U_j \to X$, $j \in J$, où
les $U_j \to X$ sont des monomorphismes.

\begin{enumerate}
\renewcommand{\theenumi}{\alph{enumi}}
\renewcommand{\labelenumi}{\theenumi)}
\item Existe-t-il\pageoriginale des topos qui possèdent la propriété (P) et
  qui ne sont pas engendrés par leurs ouverts? Les topos associés
  aux espaces topologiques sont engendrés par leurs ouverts. Si $E$
  est engendré par ses ouverts (\resp jouit de (P)), pour tout objet
  $X$ de $E$, $E_{/X}$ est engendré par ses ouverts (\resp jouit de
  (P)). Tout sous topos d'un topos engendré par ses ouverts est
  engendré par ses ouverts. La propriété (P) n'est pas une
  propriété de caractère local

\item Soient $E$ un topos, $\Ouv(E)$ la catégorie des ouverts de $E$
  munie de la topologie induite. Le foncteur d'inclusion $\Ouv(E)\to
  E$ est un morphisme de sites, d'où un morphisme de topos
  $\Pi:E\to\Ouv(E)^{\sim}$. Le foncteur $\Pi^{\ast}$ est pleinement
  fidèle. Le morphisme $\Pi$ possède vis à vis de la 2-catégorie des
  topos engendrés par leurs ouverts une propriété universelle que le
  lecteur explicitera.

\item Soient $E$ un topos engendré par ses ouverts et $\Point(E)$
  l'ensemble des points à isomorphismes près de $E$. Montrer que
  $\Point(E)$ est petit. Mettre une topologie sur $\Point(E)$ et définir
  un morphisme $\Top(\Point(E)) \to \Ouv(E)^{\sim}$ faisant de
  $\Top(\Point(E))$ un sous-topos de $\Ouv(E)^{\sim}$. Pour tout espace
  topologique $X$, montrer que tout morphisme de $\Top(X)$ dans $E$
  fournit un morphisme de $\Top(X)$ dans $\Top(\Point(E))$. Montrer
  qu'un topos engendré par ses ouverts est équivalent à un topos
  $\Top(X)$ ($X$ espace topologique) si et seulement s'il possède
  suffisamment de points. Montrer qu'il existe des topos engendrés
  par leurs ouverts qui ne sont pas équivalents à des topos
  $\Top(X)$ où $X$ est un espace topologique.

\item Soit $E$ un topos possédant la propriété $(P)$. Pour qu'un
  faisceau abélien $F$ sur $E$ soit flasque, il faut et il suffit
  que pour tout monomorphisme $X \hookrightarrow Y$, l'homomorphisme
  $F(Y) \to F(X)$ soit surjectif.

\item Soit $E$ un topos possédant la propriété $(P)$. Montrer
  que la propriété pour un faisceau $F$ d'être flasque est une
  propriété de nature locale.

\item Soient $E$ un topos engendré par ses ouverts, $e$ l'objet
  final de $E$. Pour qu'un faisceau abélien $F$ soit flasque il faut
  et il suffit que pour tout ouvert $U$ de $e$, le morphisme canonique
  $F(e) \to F(U)$ soit surjectif.
\end{enumerate}
\end{enonce*}

\begin{enonce*}[remark]{Exercice 4.17}\etiquette[4.17]{V.exer4.17}
(Faisceaux\pageoriginale flasques et changement d'univers)

Soient $C$ un $𝒰$-site (par exemple un $𝒰$-topos) et
$𝒱$ un
univers contenant $𝒰$. Notons $\epsilon : C^{\sim}_{𝒰} \to
C^{\sim}_{𝒱}$ les catégories de faisceaux correspondantes et
l'injection canonique. Soit $F$ un $𝒰$-faisceau abélien flasque
sur $C$. On se propose de montrer que $\epsilon F$ est flasque. On
remarque tout d'abord que pour tout objet $X$ de $C^{\sim}_{𝒰}$ on a
$H^q(\epsilon X, \epsilon F) = 0$ pour $q > 0$. Tout objet $Y$
de $C^{\sim}_{𝒱}$ admet une famille épimorphique $\epsilon X_i
\to Y$, $i \in I$ où les $\epsilon X_i \to Y$ sont des
monomorphismes. On en conclut que $H^q (Y, \epsilon F) =
\displaystyle \varprojlim_{\epsilon X \hookrightarrow Y}^{q} F(X)$. Pour
montrer que ces $\varprojlim^q$ sont nuls pour $q \neq 0$, on
peut se ramener au topos $\Ouv(C_{𝒱/Y})^{\sim}$ et utiliser le fait
que pour ce topos, un faisceau est flasque s'il l'est localement.
\end{enonce*}

\section{Les $R^qu_*$ et la suite spectrale de Cartan-Leray relative à un morphisme de topos}\etiquette[5]{V.sec5}

\setcounter{subsection}{-1}
\subsection{}\etiquette[5.0]{V.subsec5.0}
Soient $u : (E, A) \to (E', A')$ un morphisme de topos annelés, $u_\ast
= E \to E'$ le foncteur image directe. La notation $u_\ast$ désignera
encore l'extension aux Modules du foncteur image directe. Le foncteur
$u_\ast$ pour les Modules (à gauche pour fixer les idées) est exact
à gauche. Ses foncteurs dérivés droits sont notés
$R^qu_{\ast}$, $q \geq 0$.

\begin{enonce*}{Proposition 5.1}\etiquette[5.1]{V.prop5.1}
\begin{enumerate}
\renewcommand{\labelenumi}{\theenumi)}
\item Pour tout $A$-Module $M$, le faisceau  $R^qu_\ast M$ est le
  faisceau associé au préfaisceau $X' \mapsto H^q(u^\ast X', M)$ $(X'
  \in \ob E')$.

\item La formation des foncteurs $R^{q}u_{\ast}$ commute aux
  restrictions des scalaires.

\item La formation des $R^qu_\ast$ commute aux
  localisations. De manière précise, pour tout objet $X'$ de $E'$,
  si on désigne par $u_{/X} : E_{/u^{\ast}X'} \to E'_{/X'}$ le
  morphisme déduit de $u$ par localisation \emph{(IV 8)}, on a, pour tout
  $A$-Module $M$, un isomorphisme canonique
\begin{equation*}
R^q(u_{/X'})_\ast (M |u^\ast X') \simeq  R^qu_\ast (M) | X' \qquad q \geq
0{\quoi}.\tag{5.1.1}
\end{equation*}\etiquette[5.1.1]{V.5.1.1}
\end{enumerate}
\end{enonce*}

Désignons par $u^{\wedge}_\ast :
E^{\wedge} \to E^{\prime\wedge}$ le foncteur
image directe pour les $𝒰$-préfaisceaux
$(u^{\wedge}_\ast M = M ∘ u^{\ast})$.
Comme $u^\ast$ et $u_\ast$ sont adjoints, on a un
isomorphisme
$$
u_\ast \simeq \underline{a} u^{\wedge}_\ast {}^{0}{\quoi}, % \XXX sens ?
$$
où\pageoriginale $\underline{a}$ est le foncteur faisceau associé pour
$E'$. Comme les
foncteurs $\underline{a}$ et $u^{\wedge}_\ast$ sont exacts, on a
$$
R^q u_\ast\simeq \underline{a}
u^{\wedge}_\ast\Hh^q{\quoi},
$$
ce qui est une autre manière d'énoncer 1). L'assertion 2) résulte
alors de 1) et de \hyperref[V.cor3.5]{3.5}. Par définition du morphisme
$u_{/X'}$ on a un isomorphisme canonique \eqhyperref[V.5.1.1]{5.1.1} pour $q = 0$. Le
cas général s'en déduit en remarquant que les foncteurs de
localisation (IV 8) sont exacts et transforment les objets injectifs
en objets injectifs (\hyperref[V.prop4.11]{4.11}).

\begin{enonce*}{Proposition 5.2}\etiquette[5.2]{V.prop5.2}
Soient $u : E \to E'$ un morphisme de topos et $S'$ une famille
génératrice de $E'$. Les faisceaux $M$ acycliques pour les
foncteurs $H^{0}(u^\ast X', ⋅)$, $X' \in S'$, sont acycliques pour le
foncteur $u_\ast$. En particulier, les faisceaux flasques sont acycliques
pour $u_\ast$.
\end{enonce*}

Résulte de \hyperref[V.prop5.1]{5.1}. 1).

\begin{enonce*}{Proposition 5.3}\etiquette[5.3]{V.prop5.3}
Soient $u : E \to E'$ un morphisme de topos et $M$ un Groupe abélien
de $E$. On a une suite spectrale:
\begin{equation*}
E^{p, q}_2 = H^p(E', R^q u_\ast M) \Longrightarrow H^{p + q} (E,
M){\quoi}.\tag{5.3.1}
\end{equation*}\etiquette[5.3.1]{V.5.3.1}
Plus généralement, pour tout objet $X'$ de $E'$, on a une suite
spectrale:
\begin{equation}
E^{p, q}_2 = H^p(X', R^q u_\ast M) \Rightarrow H^{p + q} (u^\ast X',
M){\quoi}.\tag{5.3.2}
\end{equation}\etiquette[5.3.2]{V.5.3.2}
\end{enonce*}

Par définition des foncteurs images directe et réciproque, on a un
isomorphisme $H^{0} (X', u_\ast M) \simeq H^{0} (u^\ast X', M)$. Le
foncteur $u_\ast$ transforme les objets injectifs en faisceaux flasques
(\hyperref[V.subsec4.6]{4.6} et \hyperref[V.prop4.9]{4.9}). Les suites spectrales proposées sont donc des suites
spectrales de foncteurs composés (\hyperref[V.prop0.3]{0.3}).

\begin{enonce*}{Proposition 5.4}\etiquette[5.4]{V.prop5.4}
Soient $u : E \to E'$ et $v : E' \to E''$ deux morphismes de
topos. On a une suite spectrale
$$
R^p v_\ast R^q u_\ast \Longrightarrow R^{p+q} (v \circ u)_\ast {\quoi}.
$$
\end{enonce*}

On a $v_\ast u_\ast \simeq (vu)_\ast$ et le foncteur $u_\ast$
transforme les objets injectifs en\pageoriginale faisceaux flasques
(\hyperref[V.prop4.9]{4.9}) donc acycliques pour $v_\ast$ (\hyperref[V.prop5.2]{5.2}).
On a donc une suite spectrale des foncteurs composés
(\hyperref[V.prop0.3]{0.3}).

\section{Ext locaux et cohomologie à supports}\etiquette[6]{V.sec6}

\setcounter{subsection}{-1}
\subsection{}\etiquette[6.0]{V.subsec6.0}

Soient $(E, A)$ un topos annelé, $M$ un $A$-Module, à gauche pour
fixer les idées. Le foncteur $N \mapsto \Hhom_A(M, N)$, sur
la catégorie des $A$-Modules à gauches et à valeurs dans la
catégorie des Groupes abéliens est exact à gauche (IV 12). Ses
foncteurs dérivés droits sont notés:
\begin{equation*}
\Eext_A^q(M, N){\quoi}.\tag{6.0.1}
\end{equation*}\etiquette[6.0.1]{V.6.0.1}
En particulier, on a
\begin{equation*}
\Eext_A^{0}(M, N) = \Hhom_A(M, N){\quoi}.\tag{6.0.2}
\end{equation*}\etiquette[6.0.2]{V.6.0.2}
Par définition, on a des isomorphismes canoniques (\hyperref[V.subsec2.1]{2.1} et IV 12).
\begin{equation*}
H^{0} (E, \Eext_A^{0}(M, N)) = \Ext^{0}_A (E;M,
N)=\Hom_A(M, N){\quoi}.\tag{6.0.3}
\end{equation*}\etiquette[6.0.3]{V.6.0.3}
Plus généralement, pour tout objet $X$ de $E$, on a (\hyperref[V.subsec2.2]{2.2})
\begin{equation*}
H^{0} (X, \Eext_A^{0}(M, N)) = \Ext^{0}_A (X;M,
N)=\Hom_{A|X}(M|X, N|X){\quoi}.\tag{6.0.4}
\end{equation*}\etiquette[6.0.4]{V.6.0.4}

\begin{enonce*}{Proposition 6.1}\etiquette[6.1]{V.prop6.1}
\begin{enumerate}
\renewcommand{\labelenumi}{\theenumi)}
\item La formation des foncteurs $\Eext^q_A$ commute aux
  localisations. De manière précise, pour tout objet $X$ de $E$,
  on a des isomorphismes fonctoriels:
\begin{equation*}
\Eext^q_A(M, N) | X \simeq \Eext^q_{A|X}(M|X, N|X){\quoi}.\tag{6.1.1}
\end{equation*}\etiquette[6.1.1]{V.6.1.1}

\item Le faisceau $\Eext^q_A(M, N)$ est isomorphe au faisceau associé
  au préfaisceau $X \mapsto \Ext^q_A(X;M, N)$.

\item Il existe une suite spectrale
\begin{equation*}
\Ext^{p+q}_A (E;M, N) \Longleftarrow E^{p, q}_2
 = H^p(E, \Eext^q_A(F, G)){\quoi}.\tag{6.1.2}
\end{equation*}\etiquette[6.1.2]{V.6.1.2}
\end{enumerate}
Plus\pageoriginale généralement, pour tout objet $X$ de $E$, on a
une suite spectrale fonctorielle en $X$ et en les arguments $M$ et $N$
\begin{equation*}
\Ext^{p+q}_A (X;M, N) \Longleftarrow E^{p, q}_2
 = H^p(X, \Eext^q_A(M, N)){\quoi}.\tag{6.1.3}
\end{equation*}\etiquette[6.1.3]{V.6.1.3}
\end{enonce*}

Par définition, on a un isomorphisme \eqhyperref[V.6.1.1]{6.1.1} lorsque $q = 0$ (IV
12). Le cas général s'en déduit compte tenu du fait que les
foncteurs de localisation sont exacts et transforment les Modules
injectifs en Modules injectifs (\hyperref[V.subsec2.2]{2.2}). Le foncteur $N \mapsto
\Hhom_A(M, N) = \Eext^0_A(M, N)$ transforme les Modules
injectifs en faisceaux flasques donc en faisceaux acycliques pour
$H^{0}(X, ⋅)$ (\hyperref[V.prop4.10]{4.10}). Les suites spectrales \eqhyperref[V.6.1.2]{6.1.2} et \eqhyperref[V.6.1.3]{6.1.3}
sont des suites spectrales de foncteurs composés compte tenu de
\eqhyperref[V.6.0.3]{6.0.3} et \eqhyperref[V.6.0.4]{6.0.4}. Lorsque $X$ varie dans $E$, la suite spectrale
\eqhyperref[V.6.1.3]{6.1.3} fournit une suite spectrale de préfaisceaux, d'où en
passant aux faisceaux associés, une suite spectrale de
faisceaux. Comme les faisceaux associés aux préfaisceaux $X
\mapsto H^p(X, ⋅)$ sont nuls lorsque $p \neq 0$ (\hyperref[V.prop3.1]{3.1}), cette suite
spectrale dégénère et fournit l'isomorphisme annoncé dans 2).

\begin{enonce*}{Proposition 6.2}\etiquette[6.2]{V.prop6.2}
Les foncteurs $(M, N) \mapsto \Eext^q_A(M, N)$, $q \geq 0$, forment un
$\delta$-foncteur en la variable $M$ et la variable $N$. De manière
explicite, pour toute suite exacte $0 \to N' \to N \to N'' \to 0$
(\resp $0 \to M' \to M \to M'' \to 0$) on a des longues suites exactes:
\begin{equation*}
\cdots \to \Eext^q_A(M, N') \to \Eext^q_A(M, N) \to \Eext^q_A(M, N'')
\xrightarrow{\delta}\Eext^{q+1}_A(M, N') \to \cdots \tag{6.2.1}
\end{equation*}\etiquette[6.2.1]{V.6.2.1}
 (\resp
\begin{equation*}
\cdots \to \Eext^q_A(M'', N) \to \Eext^q_A(M, N) \to \Eext^q_A(M', N)
\xrightarrow{\delta}\Eext^{q+1}_A(M'', N) \to \cdots ){\quoi}.\tag{6.2.2}
\end{equation*}\etiquette[6.2.2]{V.6.2.2}
\end{enonce*}

Il résulte des propriétés générales des foncteurs $\Ext^q$ que pour
tout objet $X$ de $E$, les foncteurs $(M, N) \mapsto \Ext^q_A(X;M,
N)$, $q \geq 0$, forment un $\delta$-foncteur en chacun des
arguments, d'où l'assertion, en faisant varier $X$ et en prenant
le faisceau associé (\hyperref[V.prop6.1]{6.1}).

\setcounter{subsection}{2}
\subsection{}\etiquette[6.3]{V.subsec6.3}
Soient $(E, A)$ un topos annelé, $M$ un $A$-Module, $Z$ un fermé
de $E$ (IV 9), $U$ l'ouvert complémentaire. On note $H^{0}_Z(E,
M)$ le groupe des sections de $M$ dont le support\pageoriginale est
contenu dans $Z$ (IV 14) et $ℋ^{0}_Z(M)$ le sous-faisceau de $M$
défini par les sections de $M$ « à supports dans $Z$ » (IV
14). Les foncteurs $H^{0}_Z(E, ⋅)$ et $ℋ^{0}_Z(⋅)$ sont
exacts à gauche (IV 14). Les foncteurs dérivés sont notés
$H^q_Z(E, ⋅)$ et $ℋ^q_Z(⋅)$ respectivement et appelés les
\emph{groupes (\emph{\resp} faisceaux) de cohomologie de $M$ à support dans $Z$}\footnote{Comparer avec \SGA 2 I pour le cas des espaces topologiques ordinaires, ainsi que l'exposé de \textsc{Hartshorne}
      cité p-80 plus bas.}.

On a des isomorphismes canoniques (IV 14)
\begin{align*}
H^{0}_Z(E, M) & \simeq \Hom_A(A_Z, M) \simeq \Ext^{0}_A (E;A_Z,
M){\quoi},\tag{6.3.1}\\
ℋ^{0}_Z(M) & \simeq \Hhom_A(A_Z, M) \simeq \Eext^{0}_A
(A_Z, M){\quoi},\tag{6.3.2}
\end{align*}
d'où des isomorphismes pour tout $q \geq 0$
\begin{align*}%\etiquette[6.3.3]{V.6.3.3}
H^q_Z(E, M) &\simeq \Ext^q_A(E; A_Z, M){\quoi},\tag{6.3.3}\\
ℋ^q_Z(M) & \simeq \Eext^q_A(A_Z, M){\quoi}.\tag{6.3.4}
\end{align*}
On remarquera que $A_Z$ étant un biModule, les faisceaux
$\Eext^q_A(A_Z,M) \simeq ℋ^q_Z(M)$ sont munis canoniquement de
structures de $A$-Module.

Pour tout objet $X$ de $E$, notons $Z_{/X}$ le sous-topos fermé de
$E_{/X}$ déduit de $Z$ par localisation (c'est le complémentaire
de $U \times X$). Par définition, on a des isomorphismes canoniques
\begin{equation*}
H^{0}(X, ℋ^{0}_Z(M)) \simeq ℋ^{0}_{Z_{/X}}
(E_{/X}, M |
X){\quoi}.\tag{6.3.5}
\end{equation*}\etiquette[6.3.5]{V.6.3.5}
On pose
\begin{equation*}
H^q_Z(X, M) \simeq H^q_{Z_{/X}} (E_{/X}, M | X){\quoi}.\tag{6.3.6}
\end{equation*}\etiquette[6.3.6]{V.6.3.6}
Compte-tenu de (6.3.2), on a des isomorphismes canoniques
\begin{equation*}
H^q_Z(X, M) \simeq \Ext^q_A(X;A_Z, M){\quoi}.\tag{6.3.7}
\end{equation*}\etiquette[6.3.7]{V.6.3.7}

\begin{enonce*}{Proposition 6.4}\etiquette[6.4]{V.prop6.4}
\begin{enumerate}
\renewcommand{\labelenumi}{\theenumi)}
\item La\pageoriginale formation des foncteurs $ℋ^q_Z$ commute à la
  localisation. De manière précise, pour tout objet $X$ de $E$, et
  tout $A$-Module $M$, on a des isomorphismes canoniques
\begin{equation*}
ℋ^{0}_Z(M) | X \simeq ℋ^q_{Z_{/X}}(M |
X){\quoi}.\tag{6.4.1}
\end{equation*}\etiquette[6.4.1]{V.6.4.1}

\item Le faisceau $ℋ^q_Z(M)$ est le faisceau associé au
  préfaisceau $X \mapsto H^q_Z(X, M)$.

\item Il existe une suite spectrale:
\begin{equation*}
H^{p+q}_Z(E, M) \Longleftarrow H^p(E, ℋ^q_Z(M)){\quoi}.\tag{6.4.2}
\end{equation*}\etiquette[6.4.2]{V.6.4.2}
Plus généralement, pour tout objet $X$ de $E$, il existe une suite
spectrale
\begin{equation*}
H^{p + q}_Z(X, M)\Longleftarrow H^p(X, ℋ^q_Z(M)){\quoi}.\tag{6.4.3}
\end{equation*}\etiquette[6.4.3]{V.6.4.3}
\end{enumerate}
\end{enonce*}

On ne fait que traduire la proposition \hyperref[V.prop6.1]{6.1} à l'aide du dictionnaire \hyperref[V.subsec6.3]{6.3}.

\begin{enonce*}{Proposition 6.5}\etiquette[6.5]{V.prop6.5}
Avec les notations de \hyperref[V.subsec6.3]{6.3}, notons $j : U \to E$ le morphisme
canonique. Pour tout $A$-Module $M$, il existe une suite exacte de
faisceaux:
\begin{equation*}
0 \to ℋ^{0}_Z(M) \to M \to j_\ast (M | U) \to
ℋ^1_Z(M) \to 0{\quoi},\tag{6.5.1}
\end{equation*}\etiquette[6.5.1]{V.6.5.1}
et des isomorphismes pour $q \geq 2$:
\begin{equation*}
ℋ^q_Z(M)\simeq \Eext^{q-1}_A (A_U, M) \simeq R^{q-1}
j_\ast(M|U){\quoi}.\tag{6.5.2}
\end{equation*}\etiquette[6.5.2]{V.6.5.2}
On a de plus, une longue suite exacte
\begin{equation*}
\cdots\to H^q_Z(E, M) \to H^q(E, M) \to H^q(U, M) \to H^{q+1}_Z(E, M)
\to \cdots\footnote{cette suite exacte précise le rôle des
    invariants cohomologiques globaux $H^q_Z(E, N)$ comme des «
    groupes de cohomologie de $E$ \emph{modulo} l'ouvert $U$, à
    coefficient dans $M$ ».} \tag{6.5.3}
\end{equation*}\etiquette[6.5.3]{V.6.5.3}
et plus généralement, pour tout objet $X$ de $E$, on a une longue
suite exacte
\begin{equation*}
\cdots\to H^q_Z(X, M) \to H^q(X, M) \to H^q(X \times U, M) \to
H^{q+1}_Z(X, M) \to \cdots.\tag{6.5.4}
\end{equation*}\etiquette[6.5.4]{V.6.5.4}
\end{enonce*}

Par définition de $A_Z$, on a une suite exacte (IV 14)
\begin{equation*}
0 \to A_U \to A \to A_Z \to 0{\quoi}.\tag{6.5.5}
\end{equation*}\etiquette[6.5.5]{V.6.5.5}
Les\pageoriginale foncteurs $\Eext^q_A(A, ⋅)$ sont nuls pour $q >
0$. On a $\Hhom_A(A_U, M) \simeq j_\ast(M|U)$ (IV 14) et par
suite (\hyperref[V.subsec2.2]{2.2}) $\Eext^q_A(A_U, M) \simeq R^q j_{\ast} (M|U)$. Enfin
$\Eext^q_A(A_Z, M)\simeq ℋ^q_Z(M)$ (6.3.4). La longue
suite exacte
\eqhyperref[V.6.2.2]{6.2.2} fournit dans ce cas \eqhyperref[V.6.5.1]{6.5.1} et
\eqhyperref[V.6.5.2]{6.5.2}. Les suites exactes
\eqhyperref[V.6.5.3]{6.5.3} et \eqhyperref[V.6.5.4]{6.5.4} résultent du dictionnaire
\hyperref[V.subsec6.3]{6.3} et de la longue
suite exacte du $\delta$-foncteur $\Ext^q_A(X; \cdot ,M), q \geq 0$,
associée à \eqhyperref[V.6.5.5]{6.5.5}.

\begin{enonce*}{Proposition 6.6}\etiquette[6.6]{V.prop6.6}
\begin{enumerate}
\renewcommand{\labelenumi}{\theenumi)}
\item Les faisceaux flasques sont acycliques pour les foncteurs
  $ℋ^{0}_Z$ et $H^{0}_Z(X, ⋅)$.

\item Les foncteurs $ℋ^q_Z$ et $H^q_Z(X, ⋅)$ commutent
  aux restrictions des scalaires.
\end{enumerate}
\end{enonce*}

Soit $M$ un faisceau flasque. La suite exacte \eqhyperref[V.6.5.4]{6.5.4} fournit des
égalités $H^q_Z(X, M) = 0$ pour tout $X$ et tout $q \geq 2$ et une
suite exacte
$$
0 \to H^{0}_Z(X, M) \to H^{0}(X, M) \to H^{0}(X \times U,
M) \to H^1_Z(X, M) \to 0{\quoi}.
$$
Mais le faisceau $M$ étant flasque, le morphisme $H^{0}(X, M)
\to H^{0}(X \times U, M)$ est surjectif \hyperref[V.sec4]{4}. Par suite $H^1_Z(X,
M) = 0$ et $M$ est acyclique pour $H^{0}_Z(X;M)$. En passant aux
faisceaux associés, on en déduit que $M$ est acyclique pour le
foncteur $ℋ^{0}_Z$ (\hyperref[V.prop6.4]{6.4}). Il est clair que les foncteurs
$ℋ^{0}_Z$ et $H^{0}_Z(X, ⋅)$ commutent aux restrictions des
scalaires et comme les faisceaux flasques sont acycliques pour ces
deux foncteurs, la propriété 2) en résulte.

\begin{enonce*}{Proposition 6.7}\etiquette[6.7]{V.prop6.7}
Soient $(E, A)$ un topos annelé, $Z$ un fermé de $E$, $U$ l'ouvert
complémentaire, $M$ et $N$ deux $A$-Modules. Il existe des
isomorphismes fonctoriels en $M$ et $N$ compatibles avec les
changements de fermés:
$$
H^{0}_Z(\Hhom_A(M, N)) \simeq \Hhom_A(M,
ℋ^0_Z(N)) \simeq \Hhom_A(M \otimes_A A_Z, N){\quoi}.
$$
\end{enonce*}

Résulte de (IV 14) compte-tenu de (6.3.2).

\setcounter{subsection}{7}
\subsection{}\etiquette[6.8]{V.subsec6.8}
On pose
\begin{equation*}
\Hhom_{A, Z}(M, N) = ℋ^{0}_Z(\Hhom_A (M, N)){\quoi}.\tag{6.8.1}
\end{equation*}\etiquette[6.8.1]{V.6.8.1}
Le foncteur $N \mapsto \Hhom_{A, Z}(M, N)$ est exact à
gauche. Ses foncteurs dérivés sont notés $\Eext^q_{A, Z}(M, N)$:
ce sont \emph{les faisceaux} $\Eext$ \emph{à supports dans} $Z$. On a
donc
\begin{equation*}
\Eext^{0}_{A, Z}(M, N) = \Hhom_{A, Z}(M, N){\quoi}.\tag{6.8.2}
\end{equation*}\etiquette[6.8.2]{V.6.8.2}
Posons\pageoriginale $M_Z = M \otimes_A A_Z$. Il résulte de IV 14 que $M_Z$ est
l'image directe sur $E$ de l'image réciproque sur $Z$ de $M$. On a
donc, compte tenu de \hyperref[V.prop6.7]{6.7} des isomorphismes
\begin{equation*}
\Eext^{q}_{A, Z}(M, N) \simeq \Eext^{q}_{A}(M_Z, N){\quoi}.\tag{6.8.3}
\end{equation*}\etiquette[6.8.3]{V.6.8.3}
Passons maintenant aux invariants globaux. On pose
\begin{equation*}
\Hom_{A, Z}(M, N) = \Ext^{0}_{A, Z}(E; M, N) = H^{0}_Z(E,
\Hhom_A(M, N)){\quoi}.\tag{6.8.4}
\end{equation*}\etiquette[6.8.4]{V.6.8.4}
Le groupe $\Hom_{A, Z}(M, N)$ est le sous-groupe du groupe des
morphismes de $M$ dans $N$ dont le support est dans $Z$ (IV 14)
i.e. qui sont nuls sur $U$. Plus généralement, pour tout objet $X$
de $E$, on pose
\begin{equation*}
\Ext^{0}_{A, Z}(X;M, N) = H^{0}_Z (X, \Hhom_A(M, N)){\quoi}.\tag{6.8.5}
\end{equation*}\etiquette[6.8.5]{V.6.8.5}
Les foncteurs $N \mapsto \Ext^{0}_{A, Z} (X;M, N)$ sont exacts à
gauche. Les foncteurs dérivés sont notés $\Ext^q_{A, Z}(X, M,
N)$. Ce sont les \emph{groupes $\Ext$ à supports dans} $Z$. Les
définitions \eqhyperref[V.6.3.6]{6.3.6}, \eqhyperref[V.6.8.1]{6.8.1} et \eqhyperref[V.6.8.5]{6.8.5} et les isomorphismes \hyperref[V.prop6.7]{6.7}
fournissent des isomorphismes
\begin{equation*}
\Ext^{0}_{A, Z}(X;M, N) \simeq
\begin{cases}
H^{0} (X, \Hhom_{A, Z}(M, N)){\quoi},\\[3pt]
\Ext^{0}_A (X; M, ℋ^{0}_Z (N)){\quoi},\\[3pt]
\Ext^{0}_A (X; M_Z, N){\quoi}.
\end{cases}\tag{6.8.6}
\end{equation*}\etiquette[6.8.6]{V.6.8.6}
Le dernier isomorphisme de \eqhyperref[V.6.8.6]{6.8.6} fournit des isomorphismes
\begin{equation*}
\Ext^{q}_{A,Z}(X;M,N)\simeq \Ext^{q}_{A}(X;M_{Z},N){\quoi}.\tag{6.8.7}
\end{equation*}\etiquette[6.8.7]{V.6.8.7}

\begin{enonce*}{Proposition 6.9}\etiquette[6.9]{V.prop6.9}
\begin{enumerate}
\renewcommand{\labelenumi}{\theenumi)}
\item Il existe deux suites spectrales fonctorielles en $M$ et $N$
  compatibles avec les changements de fermés
\begin{equation*}
\Eext^{p+q}_{A, Z}(M, N) \Longleftarrow \tag{6.9.1}
\begin{cases}
E^{p, q}_2 = ℋ^p_Z(\Eext^q_A(M, N)){\quoi},\\[3pt]
'E^{p, q}_2 = \Eext^p_A(M, ℋ^q_Z(N)){\quoi}.
\end{cases}
\end{equation*}\etiquette[6.9.1]{V.6.9.1}

\item Il \pageoriginale existe trois suites spectrales fonctorielles
  en $X$, $M$ et
  $N$ compatibles avec les changements de fermés
\begin{equation*}
\Ext^{p+q}_{A, Z}(M, N) \Longleftarrow \tag{6.9.2}
\begin{cases}
E^{p, q}_2 = H^p_Z (X, \Eext^q_A(M, N)){\quoi},\\[3pt]
'E^{p, q}_2 = H^p(X, \Eext^q_{A, Z}(M, N)){\quoi},\\[3pt]
``E^{p, q}_2 = \Ext^p_A (X; M, ℋ^q_Z(N)).
\end{cases}
\end{equation*}\etiquette[6.9.2]{V.6.9.2}

\item Les faisceaux $\Eext^q_{A, Z}(M, N)$ sont canoniquement isomorphes
  aux faisceaux associés aux préfaisceaux $X \mapsto \Ext^q_{A,
  Z}(X;M, N)$.
\end{enumerate}
\end{enonce*}

Lorsque $N$ est injectif, le faisceau $\Hhom_A(M, N)$ est flasque
\hyperref[V.prop4.10]{4.10} donc acyclique pour $ℋ^{0}_Z$
(\hyperref[V.prop6.6]{6.6}). La première suite spectrale de \eqhyperref[V.6.9.1]{6.9.1} est
une suite spectrale de foncteurs composés \eqhyperref[V.6.8.1]{6.8.1}. De même,
lorsque $N$ est injectif, $ℋ^{0}_Z(N)$ est injectif
(\hyperref[V.prop4.11]{4.11}) et la deuxième suite spectrale de \eqhyperref[V.6.9.1]{6.9.1}
est une suite spectrale de foncteurs composés déduite de
\hyperref[V.prop6.7]{6.7}. Les suites spectrales de \eqhyperref[V.6.9.2]{6.9.2} sont des
suites spectrales de foncteurs composés déduites de la définition
\eqhyperref[V.6.8.5]{6.8.5} et les deux premiers isomorphismes de \eqhyperref[V.6.8.6]{6.8.6}.
Enfin, en faisant varier $X$ dans la deuxième suite spectrale de
\eqhyperref[V.6.9.2]{6.9.2} et en prenant les faisceaux associés, on obtient une
suite spectrale de faisceaux qui dégénère grâce à \hyperref[V.prop3.1]{3.1} et
qui fournit les isomorphismes de 3).

\begin{enonce*}{Proposition 6.10}\etiquette[6.10]{V.prop6.10}
Les foncteurs $(M, N) \mapsto \Eext^q_{A, Z}(M, N)$, $q \geq 0$, et
$(M, N) \mapsto \Ext^q_{A, Z}(X;M, N)$, $q \geq 0$, sont des
$\delta$-foncteurs en chacune des variables $M$ et $N$. En notant
$M_U$ le faisceau $M \otimes_A A_U$ \emph{(\cf IV 11)}, on a une longue suite
exacte
\begin{equation}
\ldots\xrightarrow{\delta}\Eext^q_{A, Z} (M, N) \to \Eext^q_A(M, N) \to
\Eext^q_A(M_U, N) \xrightarrow{\delta} \Eext^{q+1}_{A, Z}(M, N) \to
\ldots \tag{6.10.1}
\end{equation}\etiquette[6.10.1]{V.6.10.1}
et une suite exacte analogue par les groupes $\Ext$.
\end{enonce*}

Les foncteurs $\Eext^q_A(⋅ , ⋅)$ et $\Ext^q_A(X; . , ⋅)$ forment des
$\delta$-foncteurs par rapport à chacune des variables (\hyperref[V.prop6.2]{6.2}) et le
foncteur $M \mapsto M_Z$ est exact car $A_Z$ est plat (\hyperref[V.cor1.3.3]{1.3.3}). La
première assertion résulte donc des isomorphismes \eqhyperref[V.6.8.3]{6.8.3} et
\eqhyperref[V.6.8.7]{6.8.7}. La suite exacte \eqhyperref[V.6.10.1]{6.10.1} et la suite exacte analogue pour les
groupes $\Ext$ est la longue suite exacte du $\delta$-foncteur
$$
\Eext^q_A(⋅ , N), q \geq 0 (\resp \Ext^q_A(X; . , N), q \geq 0)
$$
relative\pageoriginale à la suite exacte $0 \to M_U \to M \to M_Z \to
0$, compte
tenu de \eqhyperref[V.6.8.3]{6.8.3} et \eqhyperref[V.6.8.7]{6.8.7}.

\setcounter{subsection}{10}
\subsection{}\etiquette[6.11]{V.subsec6.11}
Indiquons brièvement comment on peut étendre ces résultats au
cas des familles de supports. Soit $E$ un topos. Pour tout objet $X$
de $E$ désignons par $\Fer(X)$ l'ensemble des fermés du topos
$E_{/X}$. Cet ensemble est en correspondance biunivoque, par passage
au complémentaire, avec l'ensemble des ouverts du topos $E_{/X}$,
i.e. avec l'ensemble des sous-objets de $X$ (IV 8). Le préfaisceau
$\Fer : X \mapsto \Fer(X)$ est en fait un $𝒰$-faisceau ainsi qu'on le
vérifie immédiatement, il est donc représentable. En d'autres
termes, il existe un objet de $E$ noté encore $\Fer$ et un fermé
$Z_{\Fer}$ de $E_{/\Fer}$ tel que pour tout $X$, tout élément de
$\Fer(X)$ se déduise de $Z_{\Fer}$ par un changement de base par un
morphisme $u_Z : X \to \Fer$ uniquement déterminé par $Z$.

\begin{enonce*}{Définition 6.12}\etiquette[6.12]{V.def6.12}
On appelle famille de supports de $E$ un sous-ensemble $\Phi$ de
l'ensemble des fermés de $E$ qui possède les propriétés
suivantes:
\begin{itemize}
\item[{\rm (S1)}] La réunion d'une famille finie d'éléments de $\Phi$
  appartient à $\Phi$
\item[{\rm (S2)}] Tout fermé de $E$ contenu dans un élément de $\Phi$
  est un élément de $\Phi$.
\end{itemize}
\end{enonce*}

\setcounter{subsection}{12}
\setcounter{subsubsection}{0}
\subsubsection{}\etiquette[6.12.1]{V.subsubsec6.12.1}
Soient $E$ un topos, $\Phi$ une famille de supports de $E$, $X$ un
objet de $E$. On désigne par $\Phi(X)$ la plus petite famille de
supports de $E_{/X}$ qui contient les fermés de $E_{/X}$ déduits
de fermés de $\Phi$ par le changement de base $X \to e$ ($e$ objet
final de $E$). Le foncteur $X \mapsto \Phi(X)$ est un sous-préfaisceau
du faisceau $\Fer$. Il est donc séparé.

\subsubsection{}\etiquette[6.12.2]{V.subsubsec6.12.2}
On dit qu'une famille $\Phi$ de supports de $E$ est \emph{de
  caractère local} si elle possède la propriété suivante:

\noindent
(CL) Pour toute famille épimorphique $(X_i \to e)$, $i \in I$ (où
  $e$ est l'objet final de $E$) la suite d'ensembles
$$
\Phi \to \prod_i \Phi(X_i) \rightrightarrows \prod_{i, j} \Phi(X_i
\times X_j)
$$
est exacte.

Soit $\vieuxsheaf{a\Phi}$ le faisceau associé au préfaisceau $X \mapsto
\Phi(X)$ (\hyperref[V.subsubsec6.12.1]{6.12.1}). La condition (CL) est
équivalente à la condition que le morphisme canonique $\Phi \to
\vieuxsheaf{a\Phi} (e)$ soit\pageoriginale une bijection. (\cf la
construction du faisceau associé dans II dans le cas d'un
préfaisceau séparé). Pour vérifier (CL) on peut donc se limiter à
une famille finale de familles épimorphiques $(X_i \to e)$, $i \in
I$.

\begin{enonce*}[remark]{Exemple 6.12.3}\etiquette[6.12.3]{V.exem6.12.3}
\begin{enumerate}
\renewcommand{\labelenumi}{\theenumi)}
\item Soit $Z$ un fermé de $E$. L'ensemble des fermés de $E$
  contenus dans $Z$ est une famille de supports de $E$. Elle est de
  caractère local.

\item Soit $T$ un espace topologique et $\Phi$ une famille de supports
  paracompactifiants de $T$ \cite{SGA4-05:7}. La famille $\Phi$ n'est pas, en
  général, de caractère local.

\item Soient $T$ un espace topologique et $p$ un entier. La famille de
  $\Phi_p$ des fermés de $T$ de codimension de Krull $\geq p$, est
  une famille de supports de $T$. Elle est de caractère local.

\item Soit $E$ un topos possédant la propriété suivante: toute
  famille épimorphique $(X_i \to e)$, $i \in I$, est majorée par
  une famille épimorphique finie. Alors toute famille de supports de
  $E$ est de caractère local. En effet, en utilisant l'exemple 1),
  il suffit de montrer qu'une limite inductive filtrante
  $\Phi_\lambda$ de familles de caractère local est une famille de
  caractère local, ce qui résulte immédiatement du passage à
  la limite inductive sur la suite d'ensembles
$$
\Phi_\lambda \to \prod_i \Phi_{\lambda} (X_i) \rightrightarrows
\prod_{i, j} \Phi_\lambda (X_i \times X_j){\quoi},
$$
où $(X_i \to e)$, $i \in I$, est une famille épimorphique
finie. En effet d'après \hyperref[V.subsubsec6.12.2]{6.12.2} et l'hypothèse sur $E$, on peut se
limiter à  des familles épimorphiques finies pour vérifier les
conditions (CL) et comme les limites inductives filtrantes commutent
aux limites projectives finies (I 2), l'exactitude de ces suites
d'ensembles est conservée par passage à la limite inductive.

\item Soit $E$ un topos cohérent (\hyperrefext[VI][VI.def2.3]{2.3}). Alors d'après ce qui
  précède toute famille de support de $E$ est de caractère
  local.
\end{enumerate}
\end{enonce*}

\setcounter{subsection}{12}
\subsection{}\etiquette[6.13]{V.subsec6.13}
Soient $(E, A)$ un topos annelé, $\Phi$ une famille de supports de
$E$, $N$ un $A$-Module (à gauche pour fixer les idées), $X$ un
objet de $E$. On pose
\begin{align*}
ℋ^{0}_{\Phi}(N) &= \varinjlim_{Z \in \Phi}
ℋ^{0}_Z(N){\quoi},\tag{6.13.1}\\
H^{0}_{\Phi}(E, N) &= \varinjlim_{Z \in \Phi}
H^{0}_Z(E, N){\quoi}.\tag{6.13.2}
\end{align*}
Soit\pageoriginale $M$ un $A$-Module à gauche, on pose:
\begin{gather*}
\Eext^{0}_{A, \Phi} (M, N) = \Hhom_{A, \Phi}(M, N) =
\varinjlim_{Z \in \Phi} \Eext^{0}_{A, Z}(M, N){\quoi},\tag{6.13.3}\\
\Ext^{0}_{A, \Phi}(X;M, N) = \Hom_{A | X, \Phi(X)}(M|X, N|X) =
\varinjlim_{Z \in \Phi} \Ext^{0}_{A, Z}(X; M, N){\quoi}.\tag{6.13.4}
\end{gather*}
On a des isomorphismes canoniques:
\begin{equation*}%
\Eext^{0}_{A, \Phi} (M, N) \simeq
ℋ^{0}_{\Phi}(\Eext^{0}_A(M, N)){\quoi},\tag{6.13.5}
\end{equation*}\etiquette[6.13.5]{V.6.13.5}

\begin{equation*}%
\Ext^{0}_{A, \Phi} (X;M, N) \simeq
H^{0}((X, \Eext^{0}_A(M, N)){\quoi}.\tag{6.13.6}
\end{equation*}\etiquette[6.13.6]{V.6.13.6}

Lorsque $\Phi$ est la famille des fermés contenus dans un fermé
$Z$, on a des isomorphismes:
\begin{equation*}
\begin{cases}
H^{0}_{\Phi} \simeq ℋ^{0}_Z\\
H^{0}_{\Phi} \simeq H^{0}_Z\\
\Eext^{0}_{A, {\Phi}} (⋅ , ⋅) \simeq \Eext^{0}_{A, Z}(⋅ , ⋅)\\
\Ext^{0}_{A, {\Phi}} (X;⋅ , ⋅) \simeq \Ext^{0}_{A, Z}(X;⋅ , ⋅){\quoi}.
\end{cases}\tag{6.13.7}
\end{equation*}\etiquette[6.13.7]{V.6.13.7}
Les limites inductives qui définissent les foncteurs précédents
sont filtrantes. Par suite ces foncteurs sont exacts à gauche. Leurs
foncteurs dérivés droites sont notés:
\begin{equation}
ℋ^q_{\Phi}, H^q_{\Phi}, \Eext^q_{A, \Phi}(⋅ , ⋅), \Ext^q_{A,
  {\Phi}}(X; . , ⋅){\quoi}.\tag{6.13.8}
\end{equation}\etiquette[6.13.8]{V.6.13.8}
Ils se calculent eux aussi par limites inductives des foncteurs
$ℋ^q_Z$, $H^q_Z$, \etcn pour $Z$ parcourant $\Phi$.

Des isomorphismes \eqhyperref[V.6.13.5]{6.13.5} et (6.13.6), on tire deux
suites spectrales par passage à la limite inductive sur la première
suite spectrale de \eqhyperref[V.6.9.1]{6.9.1} et la première suite spectrale de
\eqhyperref[V.6.9.2]{6.9.2}:
\begin{gather*}
\Eext^{p+q}_{A, \Phi} (M, N) \Leftarrow E^{p, q}_2 =
ℋ^p_{\Phi}(\Eext^q_A(M, N)){\quoi},\tag{6.13.9}\\
\Ext^{p+q}_{A, \Phi} (X;M, N) \Leftarrow E^{p, q}_2 =
H^p_{\Phi}(X, \Eext^q_A(M, N)){\quoi}.\tag{6.13.10}
\end{gather*}

\subsection{}\etiquette[6.14]{V.subsec6.14}\pageoriginale
Avec les notations de \hyperref[V.subsec6.13]{6.13}, on a, sans hypothèse sur $\Phi$, un
morphisme fonctoriel
\begin{equation*}
\theta : H^{0}_{\Phi}(N) \to H^{0}(E,
ℋ^{0}_{\Phi}(N)){\quoi}.\tag{6.14.1}
\end{equation*}\etiquette[6.14.1]{V.6.14.1}
Ce morphisme est toujours injectif mais n'est pas en général un
isomorphisme. Cependant, si $\Phi$ est de \emph{caractère local}, le
morphisme $\theta$ est un isomorphisme ainsi qu'on le vérifie
immédiatement. De même si $\Phi$ est \emph{de caractère local},
on a un isomorphisme
\begin{equation}
\theta^{'} : \Ext^{0}_{A, {\Phi}}(E;M, N) \simeq H^{0} (E,
\Eext^{0}_{A, \Phi} (M, N)){\quoi}.\tag{6.14.2}
\end{equation}\etiquette[6.14.2]{V.6.14.2}

\begin{enonce*}{Proposition 6.15}\etiquette[6.15]{V.prop6.15}
Soient $(E, A)$ un topos annelé, tel que $E$ soit cohérent \emph{(VI)}
$\Phi$ une famille de supports de $E$, $M$ et $N$ deux $A$-Modules. Il
existe deux suites spectrales:
\begin{align*}%\etiquette[6.15.1]{V.6.15.1}
H^{p + q}_{\Phi} (E, N) \Leftarrow E^{p, q}_2 &= H^p_{\Phi} (E,
ℋ^q_{\Phi}(N)){\quoi},\tag{6.15.1}\\
\text{\rm Exp}^{p + q}_{A, \Phi} (E;M, N) \Leftarrow E^{p, q}_2 &= H^p (E,
\Eext^q_{A, {\Phi}}(M, N)){\quoi}.\tag{6.15.2}
\end{align*}
\end{enonce*}

Ces suites spectrales se déduisent de la suite spectrale \eqhyperref[V.6.4.2]{6.4.2} et de
la deuxième suite spectrale \eqhyperref[V.6.9.2]{6.9.2} par passage à la limite
inductive sur les fermés $Z$ de $\Phi$, compte tenu de ce que la
cohomologie d'un topos cohérent commute aux limites inductives de
faisceaux (IV 5).

\setcounter{subsection}{15}
\subsection{}\etiquette[6.16]{V.subsec6.16}
Signalons, sans démonstration, qu'on peut étendre au cas des
famille de supports la deuxième suite spectrale de \eqhyperref[V.6.9.1]{6.9.1} et la
troisième suite spectrale de \eqhyperref[V.6.9.2]{6.9.2} (avec $X $= objet final de $E$)
en supposant que le topos $E$ est cohérent et que le Module $M$ est
parfait \cite{SGA4-05:13}.

Enfin on peut aussi généraliser la notion de familles de supports
en introduisant les préfaisceaux de familles de supports, les
faisceaux de familles de supports et les groupes \pageoriginale et
faisceaux de cohomologie correspondants\footnote{\cf \cite{SGA4-05:12} chap.IV, \S
  1 (Lecture Notes 20, Springer) pour le développement de ces
  notions sous forme de fugue avec variations.}.

\section{Appendice : Cohomologie de Čech}\etiquette[7]{V.sec7}

En développant une idée due à P. Cartier, on montre dans ce
paragraphe comment on peut dans un topos quelconque, calculer la
cohomologie d'un faisceau à l'aide de recouvrements. Pour la théorie
classique des espaces paracompacts, on renvoie à \cite{SGA4-05:7}; pour une
autre méthode qui s'applique à certains topos, et en particulier au
topos étale, voir \cite{SGA4-05:14}.

\subsection{Squelette et cosquelette}\etiquette[7.1]{V.subsec7.1}

\setcounter{subsubsection}{-1}
\subsubsection{}\etiquette[7.1.0]{V.subsubsec7.1.0}

Soit $𝚫$ la catégorie des simplexes types (les
objets de $𝚫$ sont les ensembles ordonnés
$\Delta_n = [0, \ldots, n]$, les morphismes sont les applications
croissantes). Soient $E$ un $𝒰$-topos et $𝚫(E)$ de
catégorie des préfaisceaux sur $𝚫$ à \sisi{valeur}{valeurs} % l'usage est bien d'utiliser le pluriel ? \XXX
dans $E$, autrement dit la catégorie des \emph{objets
  semi-simpliciaux} de $E$. Désignons par $𝚫[n]$ la
sous-catégorie pleine de $𝚫$ définie par
les objets $\Delta_p$, $p \leq n$, par $i_n : 𝚫[n]
\to \Delta$ le foncteur d'inclusion et par $𝚫
E[n]$ la catégorie des préfaisceaux sur $𝚫[n]$
à \sisi{valeur}{valeurs} dans $E$, autrement dit, la catégorie des \emph{objets
  semi-simpliciaux tronqués à l'ordre} $n$ de $E$. D'après I
5.1, on a une suite de trois foncteurs adjoints (I 5.3)
\[
\xymatrix@C=3cm@M=.5cm{
𝚫 E[n] \ar@<17pt>[r]^{i_{n!}} \ar@<-19pt>[r]^{i_{n\ast}}&
𝚫 E  \ar@<1pt>[l]_{i^{\ast}_n}\\
}
\]
où le foncteur $i^\ast_n$ est le foncteur restriction à la
catégorie $𝚫[n]$, et où les foncteurs $i_{n\ast}$
et $i_{n!}$ sont respectivement ses adjoints à droite et à
gauche. On notera que $𝚫 E$ et $𝚫 E[n]$ sont des
$𝒰$-topos (IV 1.2) et que $(i^{\ast}_n, i_{n\ast})$ est
un morphisme de topos (IV 3.1) qui est un plongement de
$𝚫 E[n]$ dans $𝚫 E$ (I 5.6. et IV
9.1.1).

\begin{enonce*}{Définition 7.1.1}\etiquette[7.1.1]{V.def7.1.1}
On\pageoriginale note $\sk_n$ (\resp $\cosk_n$) et on appelle
\emph{foncteur squelette d'ordre $n$} (\resp \emph{foncteur
  cosquelette d'ordre $n$}) le foncteur
$i_{n!}i^{\ast}_n$ (\resp $i_{n\ast}i^\ast_n$) de $𝚫 E$
dans $𝚫 E$.
\end{enonce*}

Les foncteurs squelette et cosquelette sont d'un usage constant en
théorie des ensembles semi-simpliciaux \cite{SGA4-05:3}.

\begin{enonce*}{Proposition 7.1.2}\etiquette[7.1.2]{V.prop7.1.2}
\begin{enumerate}
\renewcommand{\labelenumi}{\theenumi)}
\item Le morphisme d'adjonction $\sk_n \to \id$ est un
  monomorphisme. Les morphismes canoniques $\sk_n \sk_m \to \sk_n$ (\resp
  $\cosk_n \to \cosk_n \cosk_m$) sont des isomorphismes lorsque $n \leq
  m$ (\resp $n \geq m$).

\item Soit $u : E \to E'$ un morphisme de topos. Le foncteur $u^\ast$,
  prolongé aux objets semi-simpliciaux et semi-simpliciaux
  tronqués, commute aux foncteurs $i_{n!}$, $i^{\ast}_n$, $\sk_n$,
  $\cosk_n$.
\end{enumerate}
\end{enonce*}

La première assertion résulte immédiatement des
définitions. Pour démontrer 2), il suffit de constater que,
d'après I 5.1, les foncteurs $i_{n!} \ldots, \cosk_n$ se calculent
à l'aide de limites inductives et de limites projectives finies.

\subsection{Un lemme d'acyclicité}\etiquette[7.2]{V.subsec7.2}

\setcounter{subsubsection}{-1}
\subsubsection{}\etiquette[7.2.0]{V.subsubsec7.2.0}

Soit $M$ un groupe abélien de $E$.
L'\emph{homologie} d'un objet semi-simplicial $K$ de $E$, à coefficients dans $M$, se définit de
la manière usuelle: On considère d'abord $A(K)$, le groupe
semi-simplicial abélien libre engendré par $K$, puis on forme le
produit tensoriel $M \otimes_{\mathbb{Z}} A(K)$; on obtient ainsi un groupe
semi-simplicial abélien de $E$. On considère alors le complexe de
\sisi{groupe abélien}{groupes abéliens} associé (en formant la somme alternée des opérateurs
bord) et on en prend l'homologie. Les objets d'homologie sont notés
$H_i(K, M)$. Ce sont des groupes abéliens de $E$. Lorsque $M$ est
le groupe $\mathbb{Z}_E$ (groupe abélien libre engendré par l'objet final de
$E$), on note plus simplement $H_i(K)$. Les groupes $H_i(K)$ sont
appelés les \emph{groupes
  d'homologie} de $K$. On dit qu'un objet semi-simplicial $K$ est
\emph{acyclique} si pour tout groupe abélien $M$ de $E$, les groupes
$H_i(K, M)$ sont nuls $(i > 0)$ ou, ce qui est équivalent, si les
groupes $H_i(K, M)$ sont nuls $(i > 0)$. % deux fois le même énoncé... \XXX

La formation de l'homologie commute aux images réciproques par les
morphismes de topos\pageoriginale.

Le but de ce numéro est de prouver le lemme suivant:

\begin{enonce*}{Lemme 7.2.1}\etiquette[7.2.1]{V.lem7.2.1}
Soient $E$ un topos, $K_{∙}$ et $L_{∙}$ deux objets semi-simpliciaux de
$E$, $v_{∙} : K_{∙} \to L_{∙}$ un morphisme d'objets semi-simpliciaux, $n$ un
entier $\geq 0$. On suppose que le morphisme $v$ possède les
propriétés:
\begin{enumerate}
\renewcommand{\labelenumi}{\theenumi)}
\item Pour tout entier $p < n$ $(p \geq 0)$ le morphisme $v_p : K_p \to
  L_p$ est un isomorphisme.

\item Le morphisme $K_n \xrightarrow{v_n} L_n$ est un épimorphisme.

\item Les morphismes canoniques $K_{∙} \to \cosk_n (K_{∙})$, $L_{∙} \to \cosk_n
  (L_{∙})$ sont des isomorphismes.
\end{enumerate}
Alors, pour tout entier $p$, le morphisme $v_p$ est un épimorphisme
et le morphisme $v_{∙}$ induit un isomorphisme sur les objets d'homologie.
\end{enonce*}

\begin{enonce*}[remark]{Remarque}
La démonstration montre en fait, qu'en adoptant une définition
convenable de l'homotopie dans les topos, le morphisme $v_{∙}$ est une
équivalence d'homotopie.
\end{enonce*}

\begin{enonce*}[remark]{Preuve}
On peut supposer que $E$ est le topos des faisceaux sur un petit site
$C$ (IV 1) et que par suite $E$ est un sous-topos d'un topos $E'$
ayant suffisamment de foncteurs fibres (par exemple le topos
$C^{\wedge}$). Notons $a^\ast : E' \to E$ le foncteur
image inverse par le morphisme de plongement $E \hookrightarrow
E'$. Soit
$$
v _{∙}[n] : K _{∙}[n] \to L _{∙} [n]
$$
le morphisme d'objets semi-simpliciaux tronqués obtenu en tronquant
$v_{∙}$ à l'ordre $n$. Notons $L_{∙} [n]'$ l'image (dans $E'$)
de $K_{∙}[n]$ par $v_{∙}[n]$, $v_{∙}[n]' : K_{∙} [n] \to L_{∙}[n]'$
le morphisme induit par $v_{∙} [n]$. Posons $L! =
i_{n\ast} L_{∙}[n]'$, $K! = i_{n\ast} K_{∙}[n]$, $v! = i_{n\ast}
v_{∙}[n]'$ (le foncteur $i_{n\ast}$ est ici relatif à $E'$).
D'après 1), le morphisme $v! : K! \to L!$ possède les propriétés 1),
2), 3)  du lemme. De plus, d'après 2), 3) et \hyperref[V.prop7.1.2]{7.1.2}, le
morphisme $a^\ast v!$ n'est autre que $v_{∙}$. Il suffit donc de
démontrer le lemme pour $v!$ et par suite on peut supposer que $E$
possède suffisamment de foncteurs fibres; donc, en utilisant ces
foncteurs fibres, on peut se ramener au cas où $E$ est la
catégorie des ensembles.
\end{enonce*}

% continuer substitution L_∙ -> L_. voire, mieux, L_{\bullet} ou variante \XXX

On constate tout d'abord que les hypothèses du lemme sur le morphisme
$v_{∙}$ sont\pageoriginale stables par tout changement de bases $M_∙
\to L_∙$ où $M_∙$ est un cosquelette d'ordre $n$. Par suite la
fibre $P_{∙}$ de $v_{∙}$ en un point base quelconque de $L_∙$ est du type
$i_{n\ast} P_{∙} [n]$ où $P_{∙}[n]$ est un complexe non vide tronqué à
l'ordre $n$ de la forme
\begin{equation*}
P_n \to e \to e \to e \to \cdots \to e{\quoi}.\tag{7.2.2}
\end{equation*}\etiquette[7.2.2]{V.7.2.2}
Notons, pour tout entier $i$, $\delta \sheaf{\Delta}_i$ l'ensemble
semi-simplicial bord du simplexe $\sheaf{\Delta}_{i}$. Tout morphisme
de $\delta\sheaf{\Delta}_{i}$ dans $P$. se prolonge en un morphisme
de $\sheaf{\Delta}_{i}$ dans $P_{∙}$. En effet, la propriété est
évidente si $i \geq n$ car $P_{∙}$ est un cosquelette d'ordre $n$ et
elle est claire pour $i < n$ d'après la description \eqhyperref[V.7.2.2]{7.2.2}.
Comme $P_{∙}$ est un complexe de Kan, $P_{∙}$ est homotopiquement trivial
\cite{SGA4-05:6}. Notons que le morphisme $v_{∙}$ est une fibration au sens de
Kan \cite{SGA4-05:6}. Comme les fibres de cette fibration sont
homotopiquement triviales, $v_{∙}$ induit un isomorphisme sur les
groupes d'homotopie de $K_∙$ et $L_∙$ en tous les points bases de
$K_∙$. D'après le théorème d'Hurewitz, ceci implique que $v_{∙}$ induit
un isomorphisme sur l'homologie \cite{SGA4-05:6}. Ceci démontre la deuxième
assertion du lemme. Pour démontrer la première assertion, on
remarque que tout morphisme d'un complexe tronqué à l'ordre $n$ dans
$L_∙[n]$ se relève en un morphisme dans $K_∙[n]$. Par suite, d'après
3), tout morphisme d'un complexe dans $L_∙$ se relève en un morphisme
dans $L_∙$. Donc $v_{∙}$ est surjectif.

\subsection{Hyper-recouvrements}\etiquette[7.3]{V.subsec7.3}

\setcounter{subsubsection}{-1}
\subsubsection{}\etiquette[7.3.0]{V.subsubsec7.3.0}
Soit $C$ un site, appartenant à l'univers $U$, où les produits
fibrés et les produits de deux objets soient
représentables. Désignons par $C^{\wedge}$
le topos des $U$-préfaisceaux sur $C$ et par $C^{\sim}$ le topos des
$U$-faisceaux sur $C$. Un objet $K$ de $C^{\wedge}$ est
dit \emph{semi-représentable} s'il est isomorphe à une somme directe
de préfaisceaux représentables.

Soit $SR(C)$ la sous-catégorie pleine de $C^{\wedge}$
définie par les objets semi-représentables. Dans la catégorie
$SR(C)$ les limites projectives pour des catégories d'indices finies
non vides sont représentables, et le foncteur d'inclusion $SR(C)
\hookrightarrow C^{\wedge}$ commute à ces limites
projectives finies. D'après une remarque déjà faite, on en
déduit que si $K_∙$ est un objet semi-simplicial tronqué à
l'ordre $n$ de $C^{\wedge}$ dont tous les
objets\pageoriginale sont
semi-représentables, le prolongement $i^+_n(K_∙)$ possède la même
propriété. Donc si $K_∙$ est un objet semi-simplicial de
$C^{\wedge}$ dont tout les objets sont
semi-représentables, alors pour tout $n$, l'objet $\cosk_n (K_∙)$
possède la même propriété.

\subsubsection{Définitions et notation}\etiquette[7.3.1]{V.subsubsec7.3.1}

\paragraph{}\etiquette[7.3.1.1]{V.par7.3.1.1}

Soit $p > 0$ un entier. Un objet semi-simplicial $K_∙$ de
$C^{\wedge}$ est appelé un \emph{hyper-recouvrement
  de type} $p$ de $C$, ou lorsqu'aucune confusion n'en résulte, un
hyper-recouvrement de type $p$, s'il possède les propriétés
suivantes:
\begin{itemize}
\item[HR1)] Pour tout entier $n \geq 0$, $K_n$ est semi-représentable.

\item[HR2)] Le morphisme canonique $K_∙ \to \cosk_p (K_∙)$ est un isomorphisme.

\item[HR3)] Pour tout entier $n \geq 0$ le morphisme canonique de préfaisceaux:
$K_{n + 1} \to (\cosk_n(K_∙))_{n + 1}$ est un isomorphisme couvrant de
préfaisceaux (II 6.2).
Le morphisme canonique de préfaisceaux $K_{0} \to e$, où $e$
est l'objet final de $C^{\wedge}$, est un morphisme
couvrant de préfaisceaux.
\end{itemize}

\paragraph{}\etiquette[7.3.1.2]{V.par7.3.1.2}
Un hyper-recouvrement de type $p$ est un hyper-recouvrement de type
$q$ pour tout $q \geq p$. Un objet semi-simplicial $K_∙$ sera appelé
un \emph{hyper-recouvrement} (de $C$) s'il possède les propriétés HR
1) et HR 3).

\paragraph{}\etiquette[7.3.1.3]{V.par7.3.1.3}
On désignera par $HR_p$ (\resp $HR$) et on appellera catégorie
des hyperrecouvrements de type $p$ (\resp catégorie des
hyper-recouvrements) la catégorie suivante:
\begin{enumerate}
\renewcommand{\theenumi}{\alph {enumi}}
\renewcommand{\labelenumi}{\theenumi)}
\item Les objets de $HR_p$ (\resp $HR$) sont les hyper-recouvrements
  de type $p$ (\resp les hyper-recouvrements).

\item Soient $K_∙$ et $L_∙$ deux objets de $HR_p$ (\resp $HR$). Un
  morphisme de $HR_p$ (\resp $HR$) de source $K_∙$ et de but $L_∙$ est
  un morphisme $v. : K_∙ \to L_∙$ d'objets semi-simpliciaux de
  $C^{\wedge}$.
\end{enumerate}

\paragraph{}\etiquette[7.3.1.4]{V.par7.3.1.4}
Soient $C$ un site où les produits fibrés soient représentables
et $X$ un objet de $C$. Un hyper-recouvrement de $X$ (\resp un
hyper-recouvrement de type $p$ de $X$) sera un hyper-recouvrement du
site $C/X$ (\resp un hyper-recouvrement de type $p$ de $C/X$). On
définit de même la catégorie des hyper-recouvrements de $X$ (\resp
du\pageoriginale hyper-recouvrement de type $p$ de $X$). On a une
suite de foncteurs d'inclusion:
$$
\cdots HR_p \hookrightarrow HR_{p + 1} \hookrightarrow \cdots
\hookrightarrow HR{\quoi}.
$$
Ces foncteurs sont pleinement fidèles. Nous noterons $HR_{\infty}$
la limite inductive des catégories $HR_p$.

\setcounter{paragraph}{5}
\paragraph{}\etiquette[7.3.1.6]{V.par7.3.1.6}
Désignons, pour tout entier $n \geq 0$, par $\sheaf{\Delta}_n^c$ % problème de compilation \XXX
l'objet semi-simplicial type de dimension $n$ à valeur dans la
catégorie des ensembles, \ie le foncteur sur $𝚫$
représenté par l'ensemble ordonné $[0, n]$. On désigne par
$\sheaf{\Delta}_n^c$ le préfaisceau sur $C$ (à valeur dans la
catégorie des ensembles semi-simpliciaux constant de valeur
$\sheaf{\Delta}_n$. Le préfaisceau $\sheaf{\Delta}_n^c$ est
donc un préfaisceau d'ensemble semi-simplicial. Les deux injections
canoniques de $\Delta_0$ dans $\Delta_1$ définissent deux
morphismes de préfaisceaux semi-simpliciaux:
\[
\xymatrix@C=2cm@M=.25cm{
\underline{\Delta}_0^{c} \ar@<3pt>[r]^{e_{0}} \ar@<-3pt>[r]_{e_{1}} &
\underline{\Delta}_1^c ,
}
\]
et définissent, par suite, pour tout préfaisceau semi-simplicial
$K_∙$, deux injections canoniques:
\[
\xymatrix@C=2cm@M=.25cm{
K_{∙} \ar@<3pt>[r]^(.45){e_{0}} \ar@<-3pt>[r]_(.45){e_{1}}&  K_{∙}
\times \underline{\Delta}_1^c .
}
\]
Deux morphismes
$
\xymatrix@C=1cm@M=.25cm{
K_{∙} \ar@<3pt>[r]^(.45){u_{0}} \ar@<-3pt>[r]_(.45){u_{1}}&  L_{∙}
}
$
de préfaisceaux semi-simpliciaux
sont dits \emph{morphismes homotopes} s'il existe un morphisme
$$
v : K_∙ \times \sheaf{\Delta}_1^c \to L_∙
$$
tel que les diagrammes:
\[
\xymatrix@R=.5cm@C=1.25cm{
K_∙ \ar[rr]^{e_i} \ar[rdd]_{u_i}&   &K_{∙}\times \sheaf{\Delta}_1^c
\ar[ldd]^v\\
                                    &   & &\txt{i = 0,1}\\
                       & L_{∙} &
}
\]
soient\pageoriginale commutatifs. Le morphisme $v$ est appelé une \emph{homotopie
  reliant $u_0$ à $u_1$}. La relation : $u_0$ et $u_1$ sont deux
  morphismes homotopes de $K_∙$ dans $L_∙$, n'est pas, en général,
  une relation d'équivalence. Cependant la relation d'équivalence
  engendrée par cette relation est compatible avec la composition
  des morphismes. Cela nous permet donc de définir la catégorie
  des préfaisceaux semi-simpliciaux à homotopie près, en passant
  au quotient par la relation d'équivalence engendrée par la
  relation d'homotopie.

\paragraph{}\etiquette[7.3.1.7]{V.par7.3.1.7}
On définit ainsi les catégories $\sheaf{HR}_p$, $\sheaf{HR}_{\infty}$,
$\sheaf{HR}$,
\emph{catégories des hyper recouvrements à homotopie près}.

\begin{enonce*}{Théorème 7.3.2}\etiquette[7.3.2]{V.thm7.3.2}
Soit $C$ un site satisfaisant aux conditions \hyperref[V.subsubsec7.3.0]{7.3.0}.
\begin{enumerate}
\renewcommand{\labelenumi}{\theenumi)}
\item La catégorie $\sheaf{HR}^{0}_p$ (\resp $\sheaf{HR}^{0}$)
  est filtrante \emph{(I 2.7)}.

\item Soient $K_∙$ un objet de $HR_p$ (\resp de $HR$), $n$ un entier tel
  que $0 \leq n \leq p$ (\resp $0 \leq n$), et
$$
u : X \to K_n
$$
un morphisme couvrant de préfaisceaux. Il existe un objet $L_∙$ de
$HR_p$ (\resp de $HR$) et un morphisme $f : L_∙ \to K_∙$, tels que le
morphisme
$$
f_n : L_n \to K_n
$$
se factorise par $u$.

\item Les faisceaux semi-simpliciaux associés \emph{(II 3.5)} aux
  hyper-recouvrements sont acycliques (\hyperref[V.subsubsec7.2.0]{7.2.0}) en degrés strictement
  positifs. Le $0$-ème faisceau d'homologie est isomorphe au faisceau
  associé au faisceau constant $𝒵$.
\end{enumerate}
\end{enonce*}

\begin{enonce*}[remark]{Preuve}
Démontrons 3). Soit ${K}^{\sim}_{∙}$ le faisceau semi-simplicial
associé à $K_∙$. Le foncteur « faisceau associé » commute au
foncteur $\cosk_n$ (\hyperref[V.prop7.1.2]{7.1.2}. 2)). Par suite le faisceau
semi-simplicial ${K}^{\sim}_{∙}$ possède les propriétés suivantes:
\begin{enumerate}
\renewcommand{\theenumi}{\alph {enumi}}
\renewcommand{\labelenumi}{\theenumi)}
\item Le morphisme canonique $K_0 \to e$ est un épimorphisme de
  faisceaux.

\item Pour tout entier $n \geq 0$, le morphisme canonique $K_{n + 1}
  \to (\cosk_n(K_∙))_{n + 1}$ est un épimorphisme de faisceaux.
\end{enumerate}
\end{enonce*}

Posons alors, pour $n \geq 0$, ${L_∙}_n = \cosk_n {K}^{\sim}_{∙}$. On a
alors, pour tout $n$, une suite de morphismes de faisceaux
semi-simpliciaux:
$$
{K_∙}^{\sim} \xrightarrow{u_n}{L_∙}_n\xrightarrow{v_n} {L_∙}_{n - 1}
\cdots \to {L_∙}_{0} \xrightarrow{v_0} e\cdot{\quoi},
$$
et\pageoriginale pour tout $j$, $0 \leq j \leq n$, $v_j$ possède les
propriétés de \hyperref[V.lem7.2.1]{7.2.1}. De plus, le morphisme $u_n$ induit
un isomorphisme sur les composants de degré $\leq n$. On déduit
alors de \hyperref[V.subsec2.1]{2.1} par récurrence sur $n$, que $K\cdot$ est
acyclique.

Pour démontrer 1) et 2) nous introduirons une terminologie.

\begin{enonce*}{Définition 7.3.3}\etiquette[7.3.3]{V.def7.3.3}
Soit $f : X_{∙} \to Y_∙$ un morphisme de préfaisceaux
semi-simpliciaux. On dit que $f$ est \emph{spécial de type} $p$ (\resp
\emph{spécial}) si:
\begin{enumerate}
\renewcommand{\labelenumi}{\theenumi)}
\item Les objets $X_{∙}$ et $Y_∙$ sont canoniquement isomorphes à leurs
  cosquelettes d'ordre $p$ et $\cosk_p (f) = f$ (\resp pas de
  conditions sur $f$, \sisi{$X^{∙}$}{$X_{∙}$} et \sisi{$Y^{∙}$}{$Y_{∙}$}).

\item Pour tout entier $n$ tel que $0 \leq n \leq p$ (\resp $0 \leq n$),
  le morphisme $\Phi_{n + 1}$ figurant dans le diagramme ci-après
  est couvrant:
\[
\xymatrix@C=1cm@R=1.25cm{
X_{n+1} \ar[rrr]^{f_{n+1}} \ar[rd]^{\Phi_{n+1}} \ar[dd]       & & & Y_{n+1}\ar[dd]\\
                                    &P_{n+1} \ar[ld] \ar[rru]& &\\
(\cosk_n X_{∙})_{n+1}\ar[rrr]_{(\cosk_n f)_{n+1}}& & &
(\cosk_n Y_{∙})_{n+1} .
}
\]
(Les flèches verticales sont définies par les morphismes canoniques
$X_{∙} \to \cosk_n X_{∙}$ et $Y_∙ \to \cosk_n Y_∙$. L'objet $P_{n+1}$ est le
produit fibré et $\Phi_{n+1}$ est l'unique flèche rendant le
diagramme commutatif).

\item Le morphisme $f_0 : X_0 \to Y_0$ est couvrant.
\end{enumerate}

Un préfaisceau semi-simplicial $K_∙$ est dit spécial de type $p$
(\resp spécial) si le morphisme canonique:
$$
K_∙ \to e_{∙}\quad (e_{∙}\quad \mbox{\rm est l'objet semi-simplicial final})
$$
est spécial de type $p$ (\resp spécial) \ie si $K_∙$ satisfait les
conditions \emph{HR 2)} et \emph{HR 3)} (\resp \emph{HR 3)}) de \hyperref[V.par7.3.1.1]{7.3.1.1}.
\end{enonce*}

\begin{enonce*}{Lemme 7.3.4}\etiquette[7.3.4]{V.lem7.3.4}
\begin{enumerate}
\renewcommand{\labelenumi}{\theenumi)}
\item Le\pageoriginale composé de deux morphismes spéciaux (\resp spéciaux de
  type $p$) est un morphisme spécial (\resp spécial de type $p$).

\item Soient $K_∙$ un préfaisceau semi-simplicial spécial (\resp
  spécial de type $p$), $X_{∙}\xrightarrow{f}Y_∙$ un morphisme
  spécial (\resp spécial de type $p$) et $u : K_∙ \to Y_∙$ un
  morphisme de complexes. Alors le produit fibré $P_{∙} =
  K_∙ ×_{Y_∙} X_∙$ est spécial (\resp spécial de type $p$).
\end{enumerate}
\end{enonce*}

La preuve de ce lemme est laissée au lecteur.

\setcounter{subsubsection}{4}
\subsubsection{}\etiquette[7.3.5]{V.subsubsec7.3.5}
\emph{Démonstration de l'assertion} 2) \emph{de} \hyperref[V.thm7.3.2]{7.3.2}.:
On peut tout d'abord supposer que $X$ est semi-représentable. Pour
tout objet semi-simplicial $L_∙$, désignons par $\Hom_{(X)}(L_∙, K_∙)$
l'ensemble des morphismes d'objets semi-simpliciaux munis d'une
factorisation $L_n \to X \xrightarrow{u} K_n$. Le foncteur
$\Hom_{(X)}(⋅ , K_∙)$ est représentable. En effet ce foncteur
transforme les limites inductives en limites projectives, et la
catégorie $𝚫(C^{\wedge})$ est un topos. On
désignera par $P$ un objet qui représente ce foncteur. Pour un objet
$Z$ de $C^{\wedge}$, désignons de même par
$j^{+}_n(Z)$ l'objet qui représente le foncteur. $L_∙ \mapsto
\Hom(L_n, Z)$ sur $𝚫 C^{\wedge}$, dont la
construction par $\varprojlim$ est explicitée dans III 1.1. On
constate ici que cette construction ne fait intervenir que des
$\varprojlim$ finies, d'où on conclut aussitôt que $j^{+}_n$
transforme préfaisceaux semi-représentables en préfaisceaux
semi-simpliciaux à composantes semi-représentables, et morphismes
couvrants $Z' \to Z$ en morphismes spéciaux (\hyperref[V.def7.3.3]{7.3.3}).

Par définition de $P_{∙}$, on a un carré cartésien:
\[
\xymatrix@C=2cm@R=1.5cm{
P_{∙} \ar[r]\ar[d]& j_n^{+}(X) \ar[d]^{j_n^{+}(u)}\\
K_{∙} \ar[r]      &  j_n^{+}(K_n)\quad.
}
\]
D'après ce qu'on vient de signaler, les objets semi-simpliciaux
$j^+_{n}(X)$ et $j^+_{n}(K_n)$ sont semi-représentables et le
morphisme $j^+_{n}(u)$ est un morphisme spécial de type $n$. Il
suffit alors pour conclure d'appliquer \hyperref[V.cor3.4]{3.4}. et le sorite \hyperref[V.subsubsec7.3.0]{7.3.0}.

\setcounter{subsubsection}{6}
\setcounter{paragraph}{-1}
\paragraph{}\etiquette[7.3.6.0]{V.par7.3.6.0}
Soient\pageoriginale $M_∙$ et $N_∙$ deux préfaisceaux
semi-simpliciaux. Le foncteur:
$$
L_∙ \mapsto \Hom(L_∙ \times M_∙, N_∙)
$$
est représentable. Le préfaisceau semi-simplicial qui le
représente sera noté
$$
\Hhom_∙(M_∙, N_∙){\quoi}.
$$
Le préfaisceau composant de degré $n$ de cet objet sera noté
$\Hhom_n(M_∙, N_∙)$. On a un isomorphisme canonique
$$
\Hhom_n(M_∙, N_∙) \xrightarrow{\sim} \Hhom_0(M_∙ \times
\sheaf{\Delta}^c_n, N){\quoi}.
$$

\begin{enonce*}{Lemme 7.3.6}\etiquette[7.3.6]{V.lem7.3.6}
Soit
\[
\xymatrix@C=2cm@R=1.5cm{
M_{∙} \ar@{^{(}->}[r]& N_{∙}\\
\sk_n \sheaf{\Delta}_{n+1}^c \ar[u] \ar@{^{(}->}[r]^{u}& \sheaf{\Delta}_{n+1}^c
\ar[u]  & \kern -1.5cm \txt{\rm ($u$ l'injection canonique)}
}
\]
un diagramme co-cartésien de préfaisceaux semi-simpliciaux.

Pour tout hyper-recouvrement $L_∙$, le morphisme:
$$
\Hhom_{0}(N_∙, L_∙) \to \Hhom_{0}(M_∙, L_∙)
$$
est couvrant.
\end{enonce*}

\begin{enonce*}[remark]{Preuve}
On a un diagramme cartésien:
\[
\xymatrix{
\Hhom_0(N_{∙},L_{∙})\ar[d] \ar[r]&
\Hhom_{0}(M_{∙},L_{∙}) \ar[d]\\
\Hhom_0(\sheaf{\Delta}^c_{n+1}, L_{∙})\ar[r]^(.45){(x)} &
\Hhom_0(\sk_n \sheaf{\Delta}^c_{n+1}, L_{∙}).
}
\]
Il suffit donc de montrer que le morphisme $(\ast)$ est couvrant. Or le
morphisme $(\ast)$ est isomorphe au morphisme:
$$
L_{n+1} \to (\cosk_n L_∙)_{n + 1}
$$
qui est couvrant par hypothèse, \cqfd
\end{enonce*}

\subsubsection{}\etiquette[7.3.7]{V.subsubsec7.3.7}
Démonstration de l'assertion 1) de \hyperref[V.thm7.3.2]{7.3.2}. : Le produit de deux
objets $HR_p$ (\resp de $HR$) est encore un objet de $HR_p$ (\resp de $HR$)
(\hyperref[V.lem7.3.4]{7.3.4}). Pour démontrer que la catégorie $\sheaf{HR}^{0}_p$
(\resp $\sheaf{HR}^{0}$) est filtrante, il suffit de montrer
qu'étant donnés\pageoriginale deux morphismes de $HR_p$ (\resp $HR$)
\[
\xymatrix@C=2cm@M=.2cm{
K_{∙} \ar@<3pt>[r]^{u_0} \ar@<-3pt>[r]_{u_1}&L_∙
}
\]
il existe un morphisme $v : M_∙ \to K_∙$ de $HR_p$ (\resp de $HR$) tel que
les morphismes:
\[
\xymatrix@C=2cm@M=.2cm{
M_{∙} \ar@<3pt>[r]^{u_0 v} \ar@<-3pt>[r]_{u_1 v}&L_∙
}
\]
soient homotopes, \ie tel qu'il existe $w : M_∙ \times
\sheaf{\Delta}_1^c \to L_∙$ rendant commutatifs les diagrammes
\[\vcenter{
\xymatrix@C=2cm@R=1.5cm{
M_{∙} \ar[d]_{v} \ar[r]^(.4){e_i}& M_{∙} \times \underline{\Delta}^c_1 \ar[d]^{w}\\
K_{∙} \ar[r]^{u_i} & L_{∙} & \kern -1.5cm\txt{$i= 0,1$      .}
}}\leqno(\rm x)
\]
Soit alors, pour tout objet semi-simplicial $M_∙$ de
$C^{\wedge}$ (non nécessairement un
hyper-recouvrement), $F(M_∙)$ l'ensemble des couples $(v, w):v : M_∙
\to K_∙$, $w : M_∙ \times \sheaf{\Delta}_1^c \to L_∙$ tels que les
diagrammes (x) soient commutatifs. Le foncteur $M_∙ \mapsto F(M_∙)$
est un foncteur contravariant en $M_∙$ qui transforme les limites
inductives en limites projectives, et qui par suite est
représentable par un objet semi-simplicial $F_{∙}$. L'objet $F_{∙}$ est
évidemment le sommet d'un diagramme cartésien :
\[
\xymatrix@=2cm{
  F_{∙}\ar[r] \ar[d]& \Hhom_{∙}(\sheaf{\Delta}^{c}_{1},L_{∙})\ar[d]^{\pi}\\
K_{∙} \ar[r]_{u_0,u_1}& L_{∙} \times L_{∙} ,
}
\]
où $\pi$ est défini par les deux inclusions de
$\sheaf{\Delta}^c_{0}$ dans $\sheaf{\Delta}^c_{1}$. Il
suffit donc, pour démontrer l'assertion, de montrer que $F$. est un
objet de $HR_p$ (\resp de HR) et pour cela, d'après
\hyperref[V.lem7.3.4]{7.3.4} 2) et le
sorite \hyperref[V.subsubsec7.3.0]{7.3.0}, il suffit de montrer que

\begin{enumerate}
\renewcommand{\theenumi}{\alph{enumi}}
\renewcommand{\labelenumi}{\theenumi)}
\item $\Hhom_∙(\sheaf{\Delta}^c_1, L_∙)$\pageoriginale est
  semi-représentable,

\item le morphisme $\pi$ est spécial de type $p$ (\resp spécial).
\end{enumerate}

On vérifie tout d'abord immédiatement que
$\Hhom_∙(\sheaf{\Delta}^c_1, L_∙)$ est un objet semi-simplicial
semi-représentable. En effet les composantes de cet objet se
calculent par limites projectives finies à partir des $L_n$, et par
suite sont semi-représentables. Vérifions maintenant b). Tout
d'abord on montre que, lorsque $L_∙$ est spécial de type $p$, le
morphisme
$$
\Hhom_∙(\sheaf{\Delta}_1^c, L_∙) \to \cosk_p \Hhom_∙(\sheaf{\Delta}^c_1, L_∙)
$$
est un isomorphisme; ce qui permet de vérifier la propriété 1)
de \hyperref[V.def7.3.3]{7.3.3}. Ensuite, le morphisme:
$$
\pi_{0} : \Hhom_{0}(\sheaf{\Delta}^c_1, L_∙) \to
L_{0} \times L_{0}
$$
est isomorphe au morphisme canonique:
$$
L_1 \xrightarrow{(d_{0}, d_{1})} L_{0} \times L_{0}
$$
qui est couvrant par hypothèse ($L_∙$ est hyper-recouvrement).

Il reste donc à vérifier la propriété 2) de \hyperref[V.def7.3.3]{7.3.3}, \ie à
vérifier que $\forall n$, dans le diagramme
% \XXX double indice à clarifier
\[
\xymatrix@C=.75cm@R=1.25cm{
\Hhom._{n+1}(\sheaf{\Delta}_1^c, L_∙) \ar[rr]^{\pi_{n+1}} \ar[dd] \ar[rd]_{\Phi_{n+1}}& &
L_{n+1} \times L_{n+1} \ar[dd]\\
     & P_{n+1} \ar[ru] \ar[ld]&\\
(\cosk_n \Hhom_∙ (\sheaf{\Delta}_1^c,L_∙))_{n+1}
\ar[rr]^{(\cosk_n\pi)_{n+1}}&  &(\cosk_n L_∙)_{n+1}
\times (\cosk_n L_∙)_{n+1}
}
\]
$(P_{n + 1}$ est le produit fibré), le morphisme $\Phi_{n+1}$ est
couvrant. Or un « adjoint functors chasing » simple montre que:
\begin{align*}
P_{n+1} &\xrightarrow{\sim} \Hhom_{0} (\sk_{n+1}
(\sheaf{\Delta}_{n+1}^c \times \sheaf{\Delta}_1^c), L_∙),\\
\Hhom_{n + 1}(\sheaf{\Delta}_1^c, L_∙) &\xrightarrow{\sim}
\Hhom_{0}(\sheaf{\Delta}_{n+1}^c \times
\sheaf{\Delta}_1^c, L_∙),
\end{align*}
et\pageoriginale que le morphisme $\Phi_{n+1}$ est isomorphe au
morphisme
$$
\Hhom_{0}(\sheaf{\Delta}_{n+1}^c \times
\sheaf{\Delta}_1^c, L_∙) \to
\Hhom_{0}(\sk_{n+1}(\sheaf{\Delta}_{n+1}^c \times \sheaf{\Delta}_1^c), L_∙)
$$
provenant de l'injection
$$
\sk_{n+1} (\sheaf{\Delta}_{n+1}^c \times \sheaf{\Delta}_1^c)
\hookrightarrow \sheaf{\Delta}_{n+1}^c \times \sheaf{\Delta}_1^c {\quoi}.
$$
Or il existe une suite de sous-objets de $\sheaf{\Delta}_{n+1}^c
\times \sheaf{\Delta}_1^c$:
$$
\sk_{n+1}(\sheaf{\Delta}_{n+1}^c \times \sheaf{\Delta}_1^c) =
M_{0} \hookrightarrow M_1 \hookrightarrow \ldots \hookrightarrow
M_k = \sheaf{\Delta}_{n+1}^c \times \sheaf{\Delta}_1^c{\quoi},
$$
telle que pour tout $0 \leq i \leq k$, le morphisme
$$
M_i \hookrightarrow M_{i+1}
$$
s'insère dans un diagramme cocartésien:
\[
\xymatrix@C=1.5cm@R=1.5cm{
M_i \ar[r] \ar@{^{(}->}[r]& M_{i+1} &\\
\sk_{n+1}\underline{\Delta} {}^c_{n+2} \ar@{^{(}->}[r] \ar[u]&
\underline{\Delta} {}^c_{n+2} \ar[u]\; .
}
\]
(On ajoute l'un après l'autre les simplexes non dégénérés de
dimension $n+2$ de $\sheaf{\Delta}_{n+1} \times
\sheaf{\Delta}_1$). D'après \hyperref[V.lem7.3.6]{7.3.6} le morphisme $\Phi_{n+1}$ est
un composé de morphismes couvrants et par suite est lui-même
couvrant, ce qui achève la démonstration de \hyperref[V.thm7.3.2]{7.3.2}.

\subsection{Le théorème d'isomorphisme}\etiquette[7.4]{V.subsec7.4}

\setcounter{subsubsection}{-1}
\subsubsection{}\etiquette[7.4.0]{V.subsubsec7.4.0}

Soit $F$ un préfaisceau en groupes abéliens sur un site $C$
satisfaisant à la condition \hyperref[V.subsubsec7.3.0]{7.3.0}. Pour tout hyper-recouvrement
$K_∙$, on posera:
$$
H^q (K_∙ , F) = H^q(\Hom_{C^{\wedge}}(K_∙, F)){\quoi}.
$$
Les $H^q(K_∙, F)$ forment un $\delta$-foncteur sur la catégorie des
préfaisceaux abéliens sur $C$, mais ils ne sont pas, en général, les
foncteurs dérivés du foncteur $H^{0}(K_∙, F)$.

Posons\pageoriginale alors:
\begin{equation*}
\Hvr{q} (C, F) = \varinjlim_{HR_{r}}
H^q(⋅ , F)\text{ } (r \mbox{ peut être infini}){\quoi},\tag{7.4.0.1}
\end{equation*}\etiquette[7.4.0.1]{V.7.4.0.1}
et

\begin{equation*}
\Hvv{q}(C, F) = \varinjlim_{HR}
H^q(⋅ , F){\quoi}.\tag{7.4.0.2}
\end{equation*}\etiquette[7.4.0.2]{V.7.4.0.2}
Les $\Hvr{q} (C, F)$ (\resp
$\Hvv{q}(C, F))$ forment encore un
$\delta$-foncteur car la catégorie $\sheaf{HR}^{0}_p$ (\resp
$\sheaf{HR}^{0})$ est filtrante (\hyperref[V.thm7.3.2]{7.3.2}). Comme les catégories
$\sheaf{HR}_{p}$ ne sont pas nécessairement des $U$-catégories, ces
$\delta$-foncteurs sont à valeurs dans la catégorie des
$V$-groupes abéliens, pour un univers $V$ convenable.

Supposons maintenant que le préfaisceau $F$ soit un faisceau, Comme
le faisceau semi-simplicial associé à un hyper-recouvrement est
acyclique (\hyperref[V.thm7.3.2]{7.3.2}), on a une suite spectrale fonctorielle en $F$ et en $K_∙$:
\begin{equation*}
H^p(K_∙, \Hh^q(F)) \Rightarrow
H^{p+q}(C^{\sim}, F){\quoi}.\tag{7.4.0.3}
\end{equation*}\etiquette[7.4.0.3]{V.7.4.0.3}
D'où, en passant à la limite, des suites spectrales:
\begin{equation*}
\begin{split}
\Hvr{p} (C, \Hh^q(F)) & \Rightarrow
H^{p+q}(C^{\sim}, F){\quoi},\\
\Hvv{p} (C, \Hh^q(F)) &\Rightarrow
H^{p+q}(C^{\sim}, F){\quoi}.
\end{split}\tag{7.4.0.4}
\end{equation*}\etiquette[7.4.0.4]{V.7.4.0.4}

\begin{enonce*}{Théorème 7.4.1}\etiquette[7.4.1]{V.thm7.4.1}
Soient $C$ un site satisfaisant la condition \hyperref[V.subsubsec7.3.0]{7.3.0}, $F$ un faisceau
abélien sur $C$.
\begin{enumerate}
\renewcommand{\labelenumi}{\theenumi)}
\item Les suites spectrales \eqhyperref[V.7.4.0.4]{7.4.0.4} définissent des isomorphismes
$$
\Hvr{q} (C, F) \xrightarrow{\sim}
H^q(C^{\sim}, F),\quad q \leq r+1
$$
et un monomorphisme
$$
\Hvr{r+2} (C, F) \longrightarrow
H^{r+2}(C^{\sim}, F){\quoi}.
$$

\item On a\pageoriginale un isomorphisme de $\delta$-foncteurs:
$$
\Hvi{q} (C, F) \xrightarrow{\sim}
\Hvv{q} (C, F)
\xrightarrow{\sim} H^q(C^{\sim}, F)\quad (\text{pour tout } q){\quoi}.
$$

\item Soient $G$ un préfaisceau de groupes abéliens et $F$ le
  faisceau associé. Il existe des isomorphismes de foncteurs:
$$
\Hvr{q} (C, G) \xrightarrow{\sim}
H^q(C^{\sim}, F),\quad q \leq r-1{\quoi},
$$
et un monomorphisme
$$
\Hvr{r}(C, G) \longrightarrow
H^r(C^{\sim}, F){\quoi}.
$$
Lorsque $G$ est un préfaisceau séparé, ce dernier morphisme est un
isomorphisme et il existe un monomorphisme:
$$
\Hvr{r+1}(C, G) \to H^{r+1}(C^{\sim},F){\quoi}.
$$

\item On a des isomorphismes:
$$
\Hvi{q}(C, G) \xrightarrow{\sim}
\Hvv{q}(G, G)
\xrightarrow{\sim}H^q(C^{\sim}, F)\quad (\text{tout } q){\quoi}.
$$
\end{enumerate}
\end{enonce*}

\begin{enonce*}[remark]{Preuve}
Il suffit de montrer que si $N$ est un préfaisceau abélien dont le
faisceau associé est nul, on a
$\Hvr{q}(C, N) = 0$ pour $q \leq r$ (\resp
$\Hvv{q}(C, N) = 0$ pour tout
 $q$); ce qui se fait immédiatement en utilisant \hyperref[V.thm3.2]{3.2}.
\end{enonce*}

\begin{enonce*}[remark]{Remarque 7.4.2}\etiquette[7.4.2]{V.rem7.4.2}
\begin{enumerate}
\renewcommand{\labelenumi}{\theenumi)}
\item On notera que pour les sites satisfaisants à la condition de
  \hyperref[V.subsubsec7.3.0]{7.3.0}, les hyperrecouvrements de type 0 sont les recouvrements
  ordinaires et les foncteurs
  $\Hvi{q}$ ne sont autres que les
  foncteurs $\Hv{q}$ introduits en \eqhyperref[V.2.4.5.1]{2.4.5.1}.

\item Soit $C$ un site à limites projectives finies représentables
  tel que pour tout objet $X$ de $C$ et toute famille couvrante $(X_i
  \to X)$, $i \in I$, il existe un $i_{0} \in I$ tel que
  $X_{i_{0}} \to X$ soit couvrant. On peut montrer alors que les
  hyper-recouvrements de type $p$ dont les composants sont
  représentables sont cofinaux dans $\sheaf{HR}_p$. De même, les
  hyper-recouvrements $K_∙$ tels qu'il existe un entier $p$ tel que
  $K_∙$ soit de type $p$ et tels que les composants de $K_∙$ soient
  représentables, sont cofinaux dans \pageoriginale $\sheaf{HR}$. (On peut
  le démontrer en s'inspirant de \hyperref[V.cor3.5]{3.5}). Ces
  hyper-recouvrements à
  composants représentables suffisent donc, dans le cas envisagé,
  pour calculer la cohomologie des faisceaux associés aux
  préfaisceaux.
\end{enumerate}
\end{enonce*}

\section{Appendice. Limites inductives locales (par P. Deligne)}\etiquette[8]{V.sec8}

Le rédacteur recommande au lecteur d'éviter, en principe, de lire
cet appendice. Il expose une technique qui permet parfois d'étendre
à des topos n'ayant pas assez de points des assertions que
l'existence de points rend triviales. Cette technique permet d'obtenir
une variante faisceautique du théorème de D. Lazard affirmant que
les modules plats sur un anneau sont les limites inductives de modules
libres de type fini.

\subsection{Catégories localement filtrantes}\etiquette[8.1]{V.subsec8.1}

\setcounter{subsubsection}{-1}
\subsubsection{}\etiquette[8.1.0]{V.subsubsec8.1.0}
Si $\Bb$ est une catégorie et si $p : \Aa \to
\Bb$ est une catégorie sur $\Bb$, on utilisera les
notations suivantes
\begin{itemize}
\item Pour $U \in \Ob$, $\Aa_U$ est la catégorie fibre
  $p^{-1}(U)$

\item Pour $f : U \to V$ dans $\Bb$ et $\lambda$, $\mu \in \Ob
  \Aa$ tels que $p(\lambda) = U$ et $p(\mu) = V$, on pose
$$
\Hom_f(\lambda, \mu) = p^{-1}(f), \mbox{ où } p : \Hom(\lambda, \mu)
\to \Hom(U, V){\quoi}.
$$
\end{itemize}

\begin{enonce*}{Définition 8.1.1}\etiquette[8.1.1]{V.def8.1.1}
Soit $\Ss$ un site. On appelle \emph{catégorie localement co-filtrante (ou
localement filtrante à gauche)} sur $\Ss$ une catégorie $p :
\Ll \to \Ss$ sur $\Ss$ telle que:
\begin{itemize}
\item[($\Ll0$)] Quels que soient $f:U \to V$ dans $\Ss$ et $\mu
  \in \Ob \Ll_V$, il existe $\lambda \in \Ob \Ll_V$ tel
  que $\Hom_f(\lambda, \mu) \neq ∅$,

\item[($\Ll1$)] Quels que soient $U \in \Ob \Ss$ et la famille
  finie $(\Ss_i)$ d'objets de $\Ll_U$, il existe un recouvrement
  $f_j : U_j \to U$ de $U$ et des objets $\mu_j \in \Ob
  \Ll_{U_{j}}$ tels que pour tout $i$ et $j$, on ait
  $\Hom_{f_{j}}(\mu_j, \lambda_i) \neq ∅$,

\item[($\Ll2$)] Quelles\pageoriginale que soient $f:U \to V$ dans $\Ss$ et la
  double flèche $(\phi_0, \phi_1): \Ss \rightrightarrows \mu$
  au-dessus de $f$, il existe un recouvrement $f_j : V_j \to V$ de
  $V$ et des flèches $\psi_j$ de but $\lambda$ telles que $p(\psi_j)
  \neq f_j$ et que $\phi_0 \psi_j = \phi_1, \psi_j$,
\end{itemize}
\end{enonce*}

\setcounter{subsubsection}{1}
\paragraph{}\etiquette[8.1.1.1]{V.par8.1.1.1}
Cette définition se simplifie lorsque $\Ll$ est fibrée
  sur $\Ss^{\ast}$: l'axiome
\begin{itemize}
\item[$(\Ll_\ast 0)$] est alors satisfait, et $(\Ll_1)$,
  $(\Ll_2)$, peuvent s'énoncer:

\item[$(\Ll'1)$] Quels que soient $U \in \Ob S$ et la famille finie
  $\lambda_i$ d'objets de $\Ll_U$, localement sur $U$, il
  existe $\mu$ s'envoyant dans tous les $\lambda_i$.

\item[$(\Ll'2)$] Quel que soit $U \in \Ob S$, toute double flèche
  $(\phi_0, \phi_1) : \lambda \rightrightarrows \mu$ dans
  $\Ll_U$ peut, localement sur $U$, être égalisée par
  une flèche de but $\lambda$.
\end{itemize}

\subsubsection{}\etiquette[8.1.2]{V.subsubsec8.1.2}
Soient $\Ss$ un site, $\Ll$ une catégorie sur $\Ss$ et
$\mathsf{e}$ un champ sur $\Ss$. On désigne par $\sheaf{\Gamma}\mathsf{e}$ la
catégorie des sections globales $\Hhom^{\cart}_{\Ss}(\Ss, C)$ de
$\mathsf{e}$. Si $\Ss$ a un objet final $S$,
$\sheaf{\Gamma}\mathsf{e}$ « n'est
autre » que $\mathsf{e}_S$.

\begin{enonce*}{Définition 8.1.2.1}\etiquette[8.1.2.1]{V.def8.1.2.1}
La catégorie des \emph{systèmes projectifs locaux d'objets de $C$,
indexés par} $\Ll$ est la catégorie $\Hhom_{\Ss}(\Ll, \mathsf{e})$.
\end{enonce*}

Désignons par $c$ (« système projectif local constant
associé ») le foncteur composé
$$
\sheaf{\Gamma} C = \Hhom^{\cart}_{\Ss} (\Ss, C) \hookrightarrow \Hhom_{\Ss}(\Ss, C)
\xrightarrow{u \mapsto u \circ p} \Hom_{\Ss}(\Ll, C)
$$
Lorsqu'on devra expliciter la dépendance en $\Ll$, on
écrira plutôt $c_{\Ll}$.

\begin{enonce*}{Définition 8.1.3}\etiquette[8.1.3]{V.def8.1.3}
  Le foncteur limite projective locale, noté $\varprojlim_{\lambda
    \in \Ll}$, est le foncteur partiellement défini
  adjoint à droite au foncteur $c$.
\end{enonce*}

Le foncteur
$$
\varprojlim : \Hhom_{\Ss}(\Ll, C) \to \sheaf{\Gamma} \mathsf{e}
$$
vérifie donc
$$
\Hom(X, \varprojlim_{\lambda \in \Ll} X_{\lambda}) \simeq
\Hom(c(X), (X_{\lambda})_{\lambda \in \Ll}){\quoi}.
$$

\begin{enonce*}{Définition 8.1.4}\etiquette[8.1.4]{V.def8.1.4}
Un $\Ss$-foncteur $F : \Ll \to \Mm$ entre catégories
localement cofiltrantes sur $\Ss$ est dit cofinal s'il vérifie les
deux conditions suivantes:
\begin{itemize}
\item[$(C \mathfrak{F} 1)$] Quels\pageoriginale que soient $U \in
\Ob \Ss$ et $\mu \in \Ob \Mm_U$, il existe un recouvrement $f_j :
U_j \to U$ de $U$ et des objets $\lambda_j \in \Ob \Ll_{U_{j}}$
tels que $\Hom_{f_j} (F (\lambda_j), \mu) \neq ∅$.

\item[$(C \mathfrak{F} 2)$] Quels que soient $f : U \to V$ dans $\Ss$ et la
double flèche $(\phi_0, \phi_1) : f(\lambda) \rightrightarrows \mu$
au-dessus de $f$, il existe un recouvrement $f_j : U_j
\rightrightarrows U$ de $U$ et des flèches $\psi_j$ de but $\lambda$
dans $\Ll$ telles que $p(\psi_j) = f_j$ et que $\phi_0
F(\psi_j) = \phi_1 F(\psi_j)$.
\end{itemize}
\end{enonce*}

Le composé de deux foncteurs cofinaux est un foncteur cofinal; les
équivalences de catégories localement filtrantes sont cofinales
(la démonstration est laissée au lecteur).

\begin{enonce*}{Proposition 8.1.5}\etiquette[8.1.5]{V.prop8.1.5}
Soient $F : \Ll \to \Mm$ un foncteur cofinal de catégories
localement cofiltrantes sur $\Ss$, $\mathsf{e}$ un champ sur $\Ss$
et $(X_{\mu})_{\mu \in \Mm}$ un système projectif local indexé par
$\Mm$. Alors, le morphisme de foncteurs en $X$ de $\Gamma C$ dans
\emph{(ens)}:
$$
F^{\ast} : \Hom(c_{\Mm}({}_X), (X_{\mu})_{\mu \in \Mm}) \to
\Hom (c_{\Ll} ({}_X), X_{F(\lambda)})_{\lambda \in
  \Ll} )
$$
est un isomorphisme, de sorte que
$$
F^{\ast} : \varprojlim_{\Mm} X_{\mu} \to \varprojlim_{\Ll} X_{F(\lambda)}
$$
est un isomorphisme, les deux membres étant simultanément
définis ou non définis.
\end{enonce*}

La démonstration utilise le lemme suivant, laissé au lecteur:

\begin{enonce*}{Lemme 8.1.6}\etiquette[8.1.6]{V.lem8.1.6}
Si $F : \Ll \to \Mm$ est cofinal, alors, quels que
soient $f : U \to V$ dans $\Ss$ et les flèches au-dessus de $f$ $\phi_i
: F(\lambda'_i) \to \mu$ $(i = 1, 2)$ il existe un recouvrement $f_j :
U_j \to U$ de $U$, des objets $\lambda'_j \in \Ob \Ll_{U_j}$
et des diagrammes commutatifs
\[
\xymatrix@=1.25cm{
              & F(\lambda_1)   \ar[rd]^{\phi_1}&\\
F(\lambda'_j) \ar@{--}[ru] \ar@{--}[rd]&                & \mu\\
              & F(\lambda_2) \ar[ru]_{\phi_2}  &
}
\]
$$
U_{j}\xrightarrow{f_{j}}U\to V
$$
de\pageoriginale projection $(f_j, f)$ dans $\Ss$.
\end{enonce*}

Quels que soient $U \in \Ob \Ss$ et $\mu \in \Mm_U$, il existe
par $(C \mathfrak{F} 1)$ un recouvrement $f_j : U_j \to U$ de $U$ et
des flèches $\phi_j : F(\lambda_j) \to \mu$ au-dessus des $f_j$ . Si
$$
\psi = (\psi_{\lambda}) \in \Hom(c_{\Ll} X,
(X_{F(\lambda)})_{\lambda \in \Ll}) : \psi_{\lambda} : X |
  p(\lambda) \to X_{F(\lambda)}
$$
est image de
$$
\phi = (\phi_{\mu}) \in \Hom(c_{\Mm} X,
(X_{\mu})_{\mu \in \Ll}) : \phi_{\mu} : X |
  p(\mu) \to X_{\mu},
$$
les diagrammes
\[
\xymatrix@=1.5cm{
                               & X_{F(\lambda_j)}
                               \ar[ld]_{\psi_{\lambda_j}} \ar[rd]^{\phi_j}&        \\
X|U_j \ar[rr]_{\phi_{\mu}|U_j} &                  & X_{\mu}|U_j
}
\]
sont commutatifs, de sorte que les $\phi_{\mu}$ sont déterminés,
localement, par $\psi$; puisque $\mathsf{e}$ est un champ, on
conclut que $F^{\ast}$ est injectif. Pour prouver $F^{\ast}$
surjectif, partons de $\psi$ et construisons $\phi$ tel que $\psi =
F^{\ast} \phi$. La construction précédente fournit les flèches
composées:
$$
\phi_{\mu, j} : \phi_j \circ \psi_{\lambda j} : X | U_j \to X_{\mu} | U_j{\quoi}.
$$
Vérifions que ces flèches se recollent, de sorte qu'il existe
$\phi_{\mu} : X | U \to X_{\mu}$ tel que $\phi_{\mu} | U_{j} =
\phi_{\mu, j}$. Soit un diagramme commutatif dans
\[
\xymatrix@=1cm{
                 & U_{ij}\ar[ld]_{g_i} \ar[rd]^{g_j}     &            \\
U_i  \ar[rd]_{f_i}            &            & U_j \ar[ld]^{f_j}\\
  %   & \rdTo_{f_i}&        & \ldTo_{f_j} & \\
                 &  U     &             \\
}
\]
Pour\pageoriginale vérifier que $\phi_{\mu, j} | U_{ij} = \phi_{\mu, i} | U_{ij}$,
il suffit de le faire localement sur $U_{ij}$, ce qui permet,
d'après 8.1.6 et $(\Ll0)$, de se limiter au cas où il
existe un diagramme commutatif
\[
\xymatrix@=1.5cm{
             & F(\lambda_{ij}) \ar[ld]_{F(a_{ij})} \ar[rd]^{F(a_{ji})} &\\
F(\lambda_i) \ar[rd]_{\phi_i}&                 & F(\lambda_{j}) \ar[ld]^{\phi_j}\\
             & \mu             &
}
\]
au-dessus du précédent. Le diagramme
\[
\xymatrix@R=1.5cm@C=2cm{
  & X_{F(\lambda_{ij})} \ar[ld] \ar[rd] \ar@/_1.68pc/[dd]&\\
X_{F(\lambda_i)}\ar[rd] & X|U_{i j}\ar[u]\ar[l]\ar[r]           & X_{F(\lambda_j)} \ar[ld]\\
                 & X_{\mu}             &
}
\]
est alors commutatif, et il existe $(\phi_{\mu})$ tel que, quel que
soit $\phi : F(\lambda) \to \mu$ au-dessus de $f : V \to U$, le
diagramme
\[
\xymatrix@=1.5cm{
                              & X_{F(\lambda)} \ar[rd]^{\phi}      &    \\
X|V \ar[rr]_{\phi_{\mu}|V}  \ar[ru]^{\psi_{\lambda}}  &
                              &     X_{\mu}|V\\}
\]
soit commutatif. Si $\sigma : \mu \to \mu'$ est une flèche de
$\Mm$, on aura $(\sigma | V) (\phi_{\mu} | V) = (\sigma | V)
\phi \cdot \psi_{\lambda} = (\sigma\phi | V) \cdot \psi_{\lambda} =
\phi_{\mu'} | V$ et dès lors, $\phi = (\phi_{\mu})$ est un morphisme
de foncteur tel que $\psi = F^{\ast}$ comme requis.

Si\pageoriginale $E$ est un ensemble ordonné, on désignera encore par $E$ la
catégorie, ayant $E$ pour ensemble d'objets, telle que $\Hom(i, j)$
est réduit à $1$ élément ou vide selon que $i \leq j$ ou non.

\begin{enonce*}{Proposition 8.1.7}\etiquette[8.1.7]{V.prop8.1.7}
Soit $\Ll$ une catégorie localement cofiltrante sur un site
$\Ss$. Il existe un ensemble ordonné $E$ et un foncteur $F$ de $E$
dans $\Ll$ tel que
\begin{enumerate}
\renewcommand{\theenumi}{\roman {enumi}}
\renewcommand{\labelenumi}{(\theenumi)}
\item $E$, regardé comme catégorie sur $\Ss$ à l'aide de $p^F : E
  \to \Ss$, est localement cofiltrante.
\item Le foncteur $F$ est cofinal.
\end{enumerate}
\end{enonce*}

Soit $\Ll'$ la catégorie localement cofiltrante produit de
$\Ll$ et de la catégorie définie par l'ensemble ordonné
$\ZZ$. La projection de $\Ll'$ dans $\Ll$ est
cofinale, de sorte qu'il suffit de prouver (\hyperref[V.prop8.1.7]{8.1.7}) pour $\Ll'$

Si $X$ est un ensemble fini de flèches dans $\Ll'$, on
désignera par $X^{0}$ l'ensemble des extrémités des
flèches dans $X$. Soit $E$ l'ensemble des parties finies non vides
de $\mathrm{Fl}(\Ll')$, telles qu'il existe $\Ss_X \in \Ob \Ll'$
vérifiant
\begin{enumerate}
\renewcommand{\theenumi}{\alph{enumi}}
\renewcommand{\labelenumi}{\theenumi)}
\item aucune flèche de $X$, sauf $1_{\lambda_{X}}$, n'aboutit à $\lambda_X$,

\item si $\mu \in X^{0}$, alors $1_{\mu} \in X$ et $\Hom(\Ss_{X},
  \mu) \cap X \neq ∅$,

\item Tout diagramme du type
\[
\xymatrix@=1.75cm{
\lambda_{X} \ar[r] \ar[rd]& \mu \ar[d]\\
            & \mu'
}
\]
de flèches de $X$ est commutatif.
\end{enumerate}

On ordonne $E$ par la relation opposée à la relation
d'inclusion. Si $X \in E$, et si $\mu \in X^{0}$, il existe une et
une seule flèche de $\lambda_X$ dans $\mu$ qui se trouve dans
$X$. Si $X$, $Y \in E$ et $X \supset Y$, soit $\lambda_{X, Y}$
l'unique flèche dans $X$ de $\lambda_X$  vers
$\lambda_Y$. Les fonctions $X \mapsto \lambda_X$ et $(X \supset Y)
\mapsto \lambda_{X, Y}$ forment un foncteur de $E$ dans $\Ll'$.

\begin{enonce*}{Lemme 8.1.8}\etiquette[8.1.8]{V.lem8.1.8}
La catégorie\pageoriginale $E$ est localement cofiltrante sur $\Ss$ et le foncteur
$\lambda : E \to \Ll'$ est cofinal.
\end{enonce*}

Il faut vérifier successivement les axiomes:
\begin{itemize}
\item[(i) \emph{axiome} ($\Ll0$)] Soit $f : V \to U$ et $X \in
  E_U$. Il existe $\lambda \in \Ob \Ll'_V$ et $\phi \in
  \Hom_f(\lambda, \lambda_X)$. De plus, on peut choisir $\lambda$ tel
  que $\lambda \not\in X^{0}$ (ici sert $\Ll' =
  \Ll \times \ZZ$). Soit $X' = X \cup \{1_{\lambda}\}
  \cup Y$ où $Y$ est l'ensemble des flèches composées $\lambda
  \xrightarrow{\phi}\lambda_{X} \xrightarrow{\psi}\mu$ $(\psi \in
  X)$. Alors, $X'
  \in E$, $\lambda_{X'} = \lambda$ et $\Hom(X', X)\neq ∅$.

\item[(ii) \emph{axiome} ($\Ll1$)] Soit $X_i$ une famille finie
  d'objets de $E_U$. Il existe un recouvrement $f_j : U_j \to U$, des
  objets $\lambda_j \in \Ob \Ll_{U_j}$ et des flèches
  $\phi_{ji} \in \Hom_{f_{j}} (\lambda_j, \lambda_{X_{i}})$. Si $x \in
  X^{0}_a \cap X^{0}_b$, il existe pour chaque $j$ un
  recouvrement $g_{jk} : U_{jk} \to U_j$ et des flèches $\psi_{jk} :
  \lambda_{jk} \to \lambda_j$ telles que $p(\psi_{jk}) = g_{jk}$ et
  qui égalisent la double flèche $\lambda_j \rightrightarrows x$,
  où une flèche appartient à $X_a$, l'autre à $X_b$. On peut
  s'arranger pour que les $\lambda_{jk}$ n'appartiennent à aucun des
  $X^{0}_i$. Remplaçons les $(f_j, \lambda_j, \phi_{ji})$ par
  les $(f_jg_{jk}, \lambda_{jk}, \phi_{ji} \psi_{jk})$, et
  répétons cette construction pour tous les $x$ dans une
  intersection de $X_i$. On obtient un nouveau système $(f_j,
  \lambda_j, \phi_{ji})$; posons
$$
X_j = \bigcup_i X_i \cup \{1_{\lambda_{j}}\} \cup Y_j
$$
où $Y_j$ est l'ensemble des flèches composées
$$
\lambda_j \xrightarrow{\phi_{ji}} \lambda_{X_{i}}
\xrightarrow{\psi} x\quad (\psi \in X_i){\quoi}.
$$
Alors, $X_j \in E_{V_{j}}$, et ces $X_j$ vérifient $(\Ll1)$.

\item[(iii) \emph{axiome} $(\Ll3)$] Trivial, faute de doubles
  flèches non triviales!

\item[(iv) \emph{axiome} $(C \mathfrak{F} 1)$] Trivial, car le foncteur
  $\lambda$ est surjectif.

\item[((v) \emph{axiome} $(C \mathfrak{F} 2)$] Soient $f : U \to V$ et
 une double
 flèche $(\phi_{0}, \phi_{1}) : \lambda_X \rightrightarrows \lambda$
 au-dessus de $f$. Il existe un recouvrement $f_j : U_j \to U$ de $U$
 et des flèches $\psi_j : \lambda_j \to \lambda_X$ telles que
 $p(\psi_j) = f_j$, $\phi_{0}$, $\psi_j = \phi_1 \psi_j$ et $\lambda_j
 \not\in X^{0}$. Soit $X_j = X \cup \{1_{\lambda_{j}}\} \cup
 Y_j$, où $Y_j$ est l'ensemble des flèches composées
$$
\lambda_j \xrightarrow[]{\psi_j} \lambda_X \xrightarrow{\psi} \mu\quad
(\psi \in X){\quoi}.
$$
Alors,\pageoriginale $X_j \in E_{V_{j}}$, et $F(\lambda_{X_{j}, X}) = \psi_j$
égalise $(\phi_0, \phi_1)$.
\end{itemize}

\setcounter{subsubsection}{9}
\setcounter{paragraph}{-1}
\paragraph{}\etiquette[8.1.9.0]{V.par8.1.9.0}

Supposons que $\Ss$ soit (le site des ouverts non vides de) l'espace
topologique réduit à un point. Une catégorie localement
cofiltrante sur $\Ss$ n'est alors autre qu'une catégorie filtrante à
gauche (\ie telle que $\Ss^{0}$ soit filtrante au sens de I 2.7),
les foncteurs cofinaux correspondant aux foncteurs cofinaux de
catégories filtrantes à gauche (I 8). La proposition \hyperref[V.prop8.1.7]{8.1.7} se
reformule

\begin{enonce*}{Proposition 8.1.9}\etiquette[8.1.9]{V.prop8.1.9}
Pour toute catégorie cofiltrante $\Ll$ il existe un ensemble ordonné
cofiltrant $E$ et un foncteur cofinal de $E$ dans $\Ll$.
\end{enonce*}

Dans ce cas particulier, la démonstration de \hyperref[V.prop8.1.7]{8.1.7} se réduit
essentiellement à la démonstration de (\hyperref[V.prop8.1.9]{8.1.9}) donnée dans I 8.

\subsubsection{}\etiquette[8.1.10]{V.subsubsec8.1.10}
Soient $q : \Tt \to \Ss$ un morphisme de sites et $p : \Ll
\to \Ss$ une catégorie localement cofiltrante sur $\Ss$. Désignons par
$\Ll^{\ast}$ la catégorie suivante:
\begin{enumerate}
\renewcommand{\theenumi}{\roman {enumi}}
\renewcommand{\labelenumi}{(\theenumi)}
\item un objet de $\Ll^{\ast}$ est un quadruple $(V, U, \phi,
  \lambda)$ tel que $V \in \Ob \Tt$, $U \in \Ob \Ss$, $\phi \in
  \Hom(V, q^{\ast} U)$ et $\lambda \in \Ob \Ll_U$.

\item une flèche $f: (V, U, \phi, \lambda) \to (V',\phi', U',
  \lambda')$ est un triple $(f_1, f_2, f_3)$ avec $f_1 \in \Hom(V,
  V')$, $f_2 \in \Hom(U, U')$, $f_3 \in \Hom (\lambda, \lambda')$ et
  $q^{\ast} f_2 \circ \phi = \phi' \circ f_1$, $p(f_3) = f_2$.

Le foncteur $f\mapsto f_1$ fait de $\Ll^{\ast}$ une catégorie sur $\Tt$.
\end{enumerate}

\begin{enonce*}{Lemme 8.1.11}\etiquette[8.1.11]{V.lem8.1.11}
La catégorie $\Ll^{\ast}$ est localement cofiltrante sur
$\Tt$.
\end{enonce*}

L'axiome $(\Ll0)$ est trivial. Vérifions $(\Ll1)$:

Soit une famille finie $L_i = (V, U_i, \phi_i, \lambda_i)$. Les
$\phi_i$ définissent un morphisme de $V$ dans le faisceau $q^{\ast} a
\prod U_i = a \prod q^{\ast} U_i$, de sorte qu'il existe des diagrammes
commutatifs
\[
\xymatrix@C=2.5cm@R=1.5cm{
V_j    \ar[r]^{f_j}   \ar[d]^{\psi_j}      & V \ar[d]^{\phi_i}\\
q^{\ast}U_j  \ar[r]_{q^{\ast}(P_{ij})}  & q^{\ast}U_i
}
\]
tels\pageoriginale que les $f_j$ recouvrent $V$.

Soit $\lambda_{ji}$ vérifiant $\Hom p_{ij} (\lambda_{ji},
\lambda_i) \neq ∅$. D'après $(\Ll1)$, quitte à
raffiner $U_j$ (et $V_j$), on peut prendre $\lambda_{ji} = \lambda_j$
indépendant de $i$, et les $L_j = (V_j, U_j, \psi_j, \lambda_j)$
vérifient $(\Ll1)$.

Vérifions $(\Ll2)$. Soit $(f_1, f_2) : L^1 \rightrightarrows
L^2$, avec $L^i = (V^i, U^i, \phi^i, \lambda^i)$ et $f_i = (f, g_i,
h_i)$. Les $\phi_i$ définissent un morphisme de $V$ dans le faisceau
$q^{\ast} a \Ker(U^1 \rightrightarrows U^2) = a \Ker(q^{\ast} U^1 \rightrightarrows
q^{\ast} U_2)$, de sorte qu'il existe des diagrammes commutatifs
\[
\xymatrix@C=2cm@R=1.5cm{
V_j \ar[r]^{f_j} \ar@{-}[d]^{\psi_j}& V^1 \ar[r] \ar@{-}[d]^{\phi^{1}}
& V^2 \ar@{-}[d]^{\phi^{2}}\\
q^{\ast}U_j \ar[r]_{p_j} & q^{\ast}U^{1} \ar@<2pt>[r] \ar@<-2pt>[r]&  U^{2} &\kern-1cm  (\mathsf{g}_1p_j =  \mathsf{g}_2p_j),
}
\]
tels que les $f_j$ recouvrent $V$. Soit $\chi_j : \lambda_j \to
\lambda^1$ tel que $p(\chi_j) = p_j$. D'après $(\Ll2)$ appliqué aux
$(\lambda_1 \chi_j, h_2 \chi_j)$, quitte à raffiner $U_j$ (et
$V_j$), on peut se débrouiller pour que $h$, $\chi_j = h_2 \chi_j$,
et les $L_j = (V_j, U_j, \psi_j, \lambda_j)$ vérifient ($\Ll_2$).

\begin{enonce*}{Proposition 8.1.12}\etiquette[8.1.12]{V.prop8.1.12}
Soit $t^{\ast} : \Ss^{\ast} \to \Ss$ une catégorie localement cofiltrante sur $\Ss$,
et munissons $\Ss^{\ast}$ de la topologie induite par celle de $\Ss$ à l'aide
de $t^{\ast}$. Alors, $t^{\ast}$ est une équivalence de sites $t : \Ss \to \Ss^{\ast}$.
\end{enonce*}

Par construction, $t_x : \Ff \mapsto \Ff \circ t^{\ast}$ transforme
faisceaux sur $\Ss$ en faisceaux sur $\Ss^{\ast}$. D'après $\Ll0$,
et $\Ll1$ appliqué à la famille vide, tout objet « assez petit »
de $\Ss$ est dans l'image de $t^{\ast}$ (\ie l'image de $t^{\ast}$
est un \emph{crible couvrant} dans $\Ss$). donc par le « lemme de
comparaison » (III 4) l'inclusion dans $\Ss$ de sa sous-catégorie % XXX hyperref manquant
pleine définie par $t^{\ast}(\Ob \Ss^{\ast})$ est une équivalence de
sites, ce qui permet de supposer $t^{\ast}$ surjectif sur les
objets.

Soit $\Ff$ un faisceau sur $\Ss^{\ast}$.

\begin{enumerate}
\renewcommand{\theenumi}{\roman {enumi}}
\renewcommand{\labelenumi}{(\theenumi)}
\item Soient $f_j : U_j \to U$ un recouvrement dans $\Ss$, $\Ss^1$ et
  $\Ss^2$ dans $\Ss^{\ast}U$ et $f^i_j \in \Hom_{f_j}(\mu_j, \lambda^i)$ pour
  $i = 1, 2$. Alors, $\Ff(\Ss^1)$ et $\Ff(\Ss^2)$ ont
  même image dans $\displaystyle\mathop{\prod}_j
  \Ff(\mu_j)$.

Soit\pageoriginale en effet $x = (x_j)$ dans l'image de
$\Ff(\Ss^1)$. Pour
que $x$ soit dans celle de $\Ff(\lambda^2)$, il suffit que
pour tout diagramme commutatif
\[
\xymatrix@C=2.5cm@R=1.5cm{
\mu \ar[r] \ar[d]& \mu_j \ar[d]\\
\mu_k \ar[r]& 2
}
\]
$x_j$ et $x_k$ aient même image dans $\Ff(\mu)$. Puisque
$\Ff$ est un faisceau, il suffit de vérifier cela après
avoir remplacé $\mu$ par les différents objets d'un de ses
recouvrements; d'après $(\Ll2)$, ceci permet de supposer
les diagrammes analogues au précédent, avec $\Ss^2$ remplacé par
$\Ss^1$, également commutatifs, auquel cas l'assertion est évidente.

\item Soient $f:U\to V$ dans $\Ss^{\ast}$, $\lambda \in \Ob
  \Ss^{\ast}_U$ et $\mu \in
  \Ob \Ss^{\ast}_V$. Il existe alors un recouvrement $f_j : U_j \to U$, un
  objet $\mu' \in \Ob {\ }^{\ast}_U$ et des flèches
\[
\xymatrix@C=2.5cm@R=1cm{
           & \Ss &\\
\lambda_j  \ar[ru]^{\phi_j} \ar[rd]_{\psi_j}&   & \mu\\
           & \mu \ar[ru]_{\psi}&
}
\]
sur
$$
U_{j}\xrightarrow{f_{j}}U\xrightarrow{f}V{\quoi}.
$$
D'après (i), il existe une et une seule flèche $\overline{f}$
rendant commutatif le diagramme suivant:
\[
\xymatrix@C=2.5cm@R=1.5cm{
\Ff(\mu)  \ar[r]^{\overline{f}} \ar[d]_{\psi}&
\Ff(\lambda) \ar[d]^{\phi_j}\\
\Ff(\mu') \ar[r]_{\psi_j}       & \Ff(\Ss_j)
}
\]
La flèche $\overline{f}$ ne change pas quand on remplace $(f_j)$
par un recouvrement plus fin donné par $f_k : U_k \to U_{j(k)}$,
qu'on remplace $\lambda_j$ par $\Ss_k$, muni de $\phi_k \in \Hom_{f_k}
(\lambda_k, \lambda_{j(k)})$,\pageoriginale et qu'on remplace $\phi_j(\psi_j)$ par
$\phi_{j(k)} \circ \phi_k (\psi_{j(k)} \circ \phi_k)$.

On laisse au lecteur le soin d'en déduire que $\overline{f}$ ne
dépend que de $f$.

\item Prouvons que $\overline{f}$ est fonctoriel en $f$. Soient donc
$$
W \xrightarrow{h} V \xrightarrow{g}U
$$
dans $\Ss$ et $\lambda^{0}$, $\mu^{0}$, $\nu^{0}$ dans $\Ob
\Ss_U$, $\Ob \Ss_{V}$ et $\Ob \Ss_W$ respectivement. Il existe un diagramme commutatif
\[
\xymatrix@C=.5cm@R=.5cm{
\nu^{0}  &         &          &         &        &\mu^{0} &        & \lambda^{0}\\
         &         &\nu^{1}\ar[rrru]   &         &        & &\mu^{1}\ar[ru] &            \\
         &         &          & {\nu^{2}}\ar[rrru] &       &&         &\\
         &        &         &         &       & &        & \\
         &\nu_{jk} \ar[luuuu]^(.7){e^{0}_{jk}} \ar[uuur]^(.7){e^{1}_{jk}} \ar[rruu]_{e^{2}_{jk}} \ar[rrrr]&  &        &         & \mu_j \ar[uuur]_{f^{1}_j}   \ar[uuuu]^(.3){f^{0}_j}  &         &
}
\]
tel que

\begin{enumerate}
\renewcommand{\theenumii}{\alph {enumii}}
\renewcommand{\labelenumii}{\theenumii)}
\item $t^{\ast} (\nu^i)$, $t^{\ast} (\mu^i)$, $t^{\ast} (e^i_{jk})$, $t^{\ast} (f^i_j)$ ne
  dépendent pas de $i$,

\item $i$ étant fixé, les $e^i_{jk}$ recouvrent $W$ et les $f^i_j$
  recouvrent $V$.
\end{enumerate}

On construit tout d'abord $\mu^1$, $\nu^1$ et $\nu^2$ par
$(\Ll0)$, et les $\mu_j$ comme en (ii). Soit $W'_{jk}$ un
recouvrement de $W$ s'envoyant dans le recouvrement $V_j = t^{\ast}(\mu_j)$
comme requis par le diagramme. Raffinant ${W'}_{jk}$ et appliquant
$(\Ll0)$ et $(\Ll1)$, on obtient un diagramme non
nécessairement commutatif, du type requis, vérifiant a) et
b).

Raffinant encore et appliquant $(\Ll2)$, on le rend
commutatif.

Reste alors à contempler le diagramme commutatif suivant:
\[
\xymatrix@C=1.5cm@R=.25cm{
  \Ff(\lambda^{0}) \ar[r]^{\bar{g}}& \Ff(\mu^0) \ar[rd]\ar[rr]^{\bar{h}}& &\Ff(\nu^{0})\ar[dddd] \\                                  &            &\Ff(\nu^1)\ar[rddd]   &   \\
\Ff(\mu^{1}) \ar[r]& \prod\Ff(\mu_j)\ar[rrdd]     &            &
  &  \\
 & & &\\
\Ff(\nu^2) \ar[rrr]&      &             &   \prod\Ff(\nu_{jk})
}
\]

\item Il\pageoriginale est clair que $\overline{g}$ est fonctoriel en
  $\Ff$: tout faisceau $\Ff$ sur $\Ss^{\ast}$ est donc
  image directe par $t$ d'un préfaisceau $\Ff_{0}$
  (uniquement déterminé) sur $\Ss$, et tout morphisme de faisceau $f :
  \Ff \to \mathfrak{g}$ est image d'un et d'un seul morphisme
  de préfaisceau $f_{0} : \Ff_{0} \to \mathfrak{g}_{0}$.
\end{enumerate}

Il en résulte déjà que $t_{\ast}$ est pleinement fidèle; il reste à prouver
que $\Ff_{0}$ est un faisceau. Si $f_j : U_j \to U$ est
un recouvrement, ce recouvrement est image d'un recouvrement dans
$\Ss^{\ast}$, de sorte que $\Ff_{0}(U)$ s'injecte dans
$\prod_j \Ff_{0}(U_j)$. Si des $x_j
\in \Ff_{0}(U_j)$ se recollent, alors, à fortiori, ils se
recollent dans $\Ss^{\ast}$ donc proviennent d'un élément de
$\Ff_{0}(U)$. Ceci achève la démonstration de (\hyperref[V.prop8.1.12]{8.1.12}).

\setcounter{subsubsection}{13}
\setcounter{paragraph}{-1}
\paragraph{}\etiquette[8.1.13.0]{V.par8.1.13.0}

Soient $q : \Tt \to \Ss$, $p : \Ll \to \Ss$ et
$\Ll^{\ast}$ comme en (\hyperref[V.subsubsec8.1.10]{8.1.10}) et $\mathsf{e}$ un champ sur
$\Tt$. Le champ $q_{\ast}\mathsf{e}$ sur $\Ss$ est le champ «
défini » par $(q_{\ast} \mathsf{e})_U = \mathsf{e}_q x_{U}$.

Si $(X_{\lambda})_{\lambda \in \Ll}$ est un système projectif
local indexé par $\Ll$ et à valeur dans $\mathsf{e} : (X_L)_{L
\in \mathsf{L^{x}}}$ par la formule
$$
X_L = \phi^{\ast} X_{\lambda} \text{ pour } L = (V, U, \phi, \lambda){\quoi}.
$$

On laisse au lecteur le soin de vérifier que:

\begin{enonce*}{Proposition 8.1.13.1}\etiquette[8.1.13.1]{V.prop8.1.13.1}
Avec les notations précédentes, on a
$$
\varprojlim_{\lambda \in \Ll} X_{\lambda} \xrightarrow{\sim}
\varprojlim_{L \in \Ll^{\ast}} X_L{\quoi},
$$
les deux membres étant simultanément définis ou non.
\end{enonce*}

Soient $\Ss$ un site et $p^{\ast} : \Ss^{\ast} \to \Ss$ une catégorie localement
filtrante sur $\Ss$. Si $\mathsf{e}$ est un champ sur $\Ss$, alors $p_{\ast}
\mathsf{e}$ est un champ sur $\Ss^{\ast}$ (\hyperref[V.par8.1.13.0]{8.1.13.0}); de plus, le foncteur
identique : $\Ss^{\ast} \to \Ss^{\ast}$ est une catégorie localement filtrante sur
$\Ss^{\ast}$, muni de la topologie induite, et tout système inductif local
$(X_{\lambda})$ indexé par $\lambda^{\ast}$ définit un système inductif
local, indexé par $\Ss^{\ast}$, sur $\Ss^{\ast}$.

On vérifie aisément:

\begin{enonce*}{Proposition 8.1.14}\etiquette[8.1.14]{V.prop8.1.14}
Avec\pageoriginale les notations précédentes, on a
$$
\varprojlim_{\lambda \in \Ss^{\ast} /\Ss} X_{\lambda} =
\varprojlim_{\lambda \in \Ss^{\ast} /\Ss^{\ast}} X_{\lambda} {\quoi},
$$
les deux membres étant simultanément définis ou non..
\end{enonce*}

\subsection{Limites inductives locales dans les catégories de
  faisceaux}\etiquette[8.2]{V.subsec8.2}

\setcounter{subsubsection}{-1}
\subsubsection{}\etiquette[8.2.0]{V.subsubsec8.2.0}

Soit $p : \Tt \to \Ss$ un morphisme de sites. Soit $\mathsf{e}$
le champ sur $\Ss$ suivant
\begin{enumerate}
\renewcommand{\theenumi}{\roman {enumi}}
\renewcommand{\labelenumi}{(\theenumi)}
\item Un objet de $\mathsf{e}$ est un couple $(U, \Ff)$ d'un
  objet de $\Ss$ et d'un faisceau $\Ff$ sur $p^{\ast}
  U$.

\item Une flèche $f : (U, \Ff) \to (V, \Gg)$ est un couple formé
  d'une flèche $f_{0} : U \to V$ et d'un morphisme de
  $\Tt/U$-faisceaux:
$$
f_1 : p^{\ast} (f_0)^{\ast} \Gg \to \Ff
$$

\item Le foncteur de $\mathsf{e}$ dans $\Ss$ est donné par $f \mapsto
  f_0$.
\end{enumerate}

On prendra garde que $\mathsf{e}_U$ est la \emph{catégorie opposée} de
  la catégorie des faisceaux sur $p^{\ast} U$.

Soit $\Ll$ une catégorie localement cofiltrante (à gauche) sur $\Ss$.

\begin{enonce*}{Définition 8.2.1}\etiquette[8.2.1]{V.def8.2.1}
Avec les notations précédentes, la catégorie des systèmes inductifs
locaux, indexés par $\Ll$, de faisceaux sur $\Tt$, est la catégorie
$\Hom_{\Ss}(\Ll, \mathsf{e})^{\circ}$.
\end{enonce*}

Le foncteur « limite projective »  de \hyperref[V.def8.1.3]{8.1.3} s'appelle ici \emph{foncteur
« limite inductive locale »}. C'est un foncteur de
$\Hom_{\Ss}(\Ll, \mathsf{e})^{\circ}$ dans $(\Gamma
\mathsf{e})^{\circ} = (\Tt)^{\sim}$.

\begin{enonce*}{Théorème 8.2.2}\etiquette[8.2.2]{V.thm8.2.2}
Le foncteur « limite inductive locale »  de la catégorie des
systèmes inductifs de faisceaux sur $\Tt$, indexés par
$\Ll$, dans la catégorie des faisceaux sur $\Tt$, est
partout défini, commute aux limites projectives finies et commute aux
limites inductives quelconques.
\end{enonce*}

Les\pageoriginale propositions 8.1.13,
\hyperref[V.prop8.1.14]{8.1.14} et \hyperref[V.prop8.1.12]{8.1.12} permettent de se ramener au
cas où $\Ss = \Tt = \Ll$. Un système inductif local est
alors la donnée, pour chaque $U \in \Ob \Ss$, d'un faisceau
$\Ff_U$ sur $\Ss/U$ et, pour chaque flèche $f: V \to U$ dans $\Ss$
d'une flèche, fonctorielle en $f$, de $f^{\ast} \Ff_U$ dans
$\Ff_V$.

Pour de tels systèmes inductifs, les $\varinjlim$ se calculent
comme suit:

\begin{enonce*}{Lemme 8.2.3}\etiquette[8.2.3]{V.lem8.2.3}
Soit $U \to \Ff_U$ une section de la catégorie des
préfaisceaux sur les objets de $\Ss$. On a
$$
a(U \mapsto \Gamma(U, \Ff_U)) \simeq \varinjlim_{U} a \Ff_U{\quoi}.
$$
\end{enonce*}

Par définition, le second membre représente le foncteur qui, à
chaque faisceau $\Ff$ sur $\Ss$, associe l'ensemble des systèmes
cohérents de flèches $\phi_U : \Ff_U \to \Ff|U$. À son tour, une
flèche $\phi_U$ est un système cohérent de flèches $\phi_g$, une
pour chaque $g : V \to U$, $\phi_g : \Ff_U (V) \to \Ff(V)$.
Appliquons \hyperref[V.lem8.1.11]{8.1.11} au foncteur identique de $\Ss$,
obtenant ainsi $\Ss^{\ast}$ localement filtrante sur $\Ss$. On a vu
que
$$
\varinjlim_{g \in \Ss^{\ast}} \Ff_U(V_g)^{\sim} \simeq
\varinjlim_{U} a \Ff_U{\quoi},
$$
où $V_g$ est la source de $g$ et ${}^{\sim}$ le foncteur « faisceau
constant engendré ». Le foncteur $U \mapsto 1_U$ de $\Ss$ dans $\Ss^{\ast}$
est cofinal, d'où encore par \hyperref[V.prop8.1.5]{8.1.5}
$$
\varinjlim_{U} \Ff_U(U)^{\sim} \simeq
\varinjlim_{U} a \Ff_U{\quoi}.
$$
Revenant aux définitions, on trouve enfin
$$
\varinjlim_{U} \Ff_U(U)^{\sim} = a (U \mapsto \Gamma(U,
\Ff_U)){\quoi}.
$$

Ceci montre que dans le cas auquel on s'est réduit, le foncteur
$\varinjlim$ est partout défini, et qu'il se calcule comme
composé du foncteur « faisceau engendré par un préfaisceau » et du
foncteur $(\Ff_U) \mapsto (U \mapsto \Gamma (U,
\Ff_U))$. Ces deux foncteurs commutent aux limites projectives
finies, et le foncteur $\varinjlim$ commute de plus aux
limites inductives quelconques de par sa définition comme foncteur
adjoint.

Le \pageoriginale lemme \hyperref[V.lem8.2.3]{8.2.3} fournit un procédé systématique pour
démontrer des propriétés du foncteur limite inductive locale à partir
des propriétés analogues du foncteur « faisceau associé à un
préfaisceau ». On trouve ainsi:

\begin{enonce*}{Proposition 8.2.4}\etiquette[8.2.4]{V.prop8.2.4}
Soit $\Tt$, $\Ss$ et $\Ll$ comme précédemment,
$\Aa_{\lambda}$ un système inductif local de faisceaux
d'anneaux sur $\Tt$, $\Ll_{\lambda}$ (\resp
$\Mm_{\lambda}$) un système inductif local de faisceaux de
modules à droite (\resp à gauche) sur $\Aa_{\lambda}$, tous
trois indexés par $\Ll$.

On a
$$
\varinjlim (\Ll_{\lambda}
\otimes_{\Aa_{\lambda}} \Mm_{\lambda}) =
\varinjlim \Ll_{\lambda} \mathop{\otimes}_{\varinjlim
    \Aa_{\lambda}} \varinjlim \Mm_{\lambda}{\quoi}.
$$
\end{enonce*}

(Se ramener au cas $\Tt = \Ss = \Ll$, et noter que les
deux membres coïncident avec le faisceau engendré par le préfaisceau
des
$$
\Ll_U (U) \mathop{\otimes}_{\Aa_{U} (U)} \Mm_U (U) ){\quoi}.
$$

\begin{enonce*}{Proposition 8.2.5}\etiquette[8.2.5]{V.prop8.2.5}
Soit $\Tt_1 \xrightarrow{q} \Tt_2 \to \Ss$ deux morphismes de sites,
et $(F_{\lambda})$ un système inductif local de faisceau sur $\Tt_2$
indexé par une catégorie localement cofiltrante $\Ll$ sur $\Ss$. On
a
$$
q^{\ast} \varinjlim_{\lambda \in \Ll} F_{\lambda} \simeq
\varinjlim_{\lambda \in \Ll} q^{\ast} F_{\lambda} {\quoi}.
$$
\end{enonce*}

Ceci résulte aussitôt de ce que $q^{\ast}$ admet un adjoint à droite $q_{\ast}$.

Soient $\Ss$ et $\Tt$ deux sites dans lesquels les
$\varprojlim$ finies sont représentables, soit
$p:\Tt\to\Ss$ un morphisme de sites tel que $p^{\ast}:\Ss\to\Tt$
commute aux $\varprojlim$ finies, et soit $\Ll$
la catégorie localement cofiltrante sur $\Tt$ ayant pour
objets les triples $(V, U, \phi)$ où $V \in \Ob \Tt$, $U \in
\Ob \Ss$ et $\phi \in \Hom (V, p^{\ast} U)$. On l'obtient en appliquant \hyperref[V.lem8.1.11]{8.1.11}
à $1_{\Ss} : \Ss \to \Ss$.

\begin{enonce*}{Proposition 8.2.6}\etiquette[8.2.6]{V.prop8.2.6}
Pour tout faisceau $\Ff$ sur $\Tt$, on a
$$
\Ff = \varinjlim_{\lambda \in \Ll} \phi^{\ast}
\phi_{\ast} \Ff
$$
où, par abus de notation, $\phi$ désigne le morphisme induit par $p$
de $\Tt/V$ dans $\Ss/V$.
\end{enonce*}

Soit\pageoriginale $\phi'$ le foncteur image réciproque au sens des préfaisceaux;
de \hyperref[V.lem8.2.3]{8.2.3}, on tire
$$
\varinjlim_{\lambda \in \Ll} \phi^{\ast} \phi_{\ast}
\Ff = \varinjlim_{\lambda \in \Ll} a {\phi}'
\phi_{\ast} \Ff = \varinjlim_{\lambda \in \Ll}
\Ff (V)^{\sim} = \varinjlim_{V} \Ff (V)^{\sim}
= \Ff{\quoi}.
$$

\begin{enonce*}{Proposition 8.2.7}\etiquette[8.2.7]{V.prop8.2.7}
Soient $p : T \to S$ un morphisme de topos, $\Oo_S$ un
faisceau d'anneaux sur $S$ et $\Oo_T$ son image réciproque sur
$T$. Pour qu'un $\Oo_T$-Module à droite $\Mm$ soit
plat sur $\Oo_T$, il faut et il suffit que le foncteur
\begin{equation*}
\Nn \to \Mm \mathop{\otimes}_{\Oo_{T}} p^{\ast} \Nn\tag{8.2.7.1}
\end{equation*}\etiquette[8.2.7.1]{V.8.2.7.1}
de la catégorie des $\Oo_S$-Modules à gauche dans celle des
faisceaux abéliens sur $T$ soit exact.
\end{enonce*}

L'assertion « il faut » est triviale; supposons donc le foncteur
\eqhyperref[V.8.2.7.1]{8.2.7.1} exact, et prouvons que $\Mm$ est plat.

On peut supposer $p$ défini par un morphisme de sites $p :
\Tt \to \Ss$ vérifiant les hypothèses faites en 2.6. Soit $V
\in \Ob \Tt$, $U \in \Ob \Ss$ et $\phi : V \to p^{\ast} U$. On désigne
encore par $\phi$ le morphisme de sites induit: $\phi :
\Tt/V \to \Ss/U$.

\begin{enonce*}{Lemme 8.2.8}\etiquette[8.2.8]{V.lem8.2.8}
Le foncteur $\Nn \mapsto \Mm | V
\displaystyle\mathop{\otimes}_{\Oo_{V}} \phi^{\ast} \Nn$,
de la catégorie des $\Oo_S | U$-Modules sur $\Ss/U$ dans celle
des $\Oo_T | V$-Modules sur $\Tt/V$, est exact.
\end{enonce*}

On se ramène à prendre $V = p^{\ast} U$. Soient $j$ et $j'$ les «
morphismes d'inclusion »: $j : \Ss / U \to \Ss$ et $j': \Tt /p^{\ast}
U \to \Tt$. On a:
$$
j'_{!} (\Mm | V \mathop{\otimes}_{\Oo_{V}} \phi^{\ast}
\Nn) = \Mm \mathop{\otimes}_{\Oo_T} p^{\ast} j_{!} \Nn{\quoi},
$$
d'où 8.2.8 puisque $j_{!}$ et $j'_{!}$ sont exacts et fidèles.

Soit alors $Q$ un $\Oo_T$-Module. On a (8.2.6)
$$
Q = \varinjlim_{\phi : V \to p^{\ast} U} \phi^{\ast} \phi_{\ast} Q | V
$$
de sorte que (\hyperref[V.prop8.2.4]{8.2.4})
$$
\Mm \mathop{\otimes}_{\Oo_{T}} Q = \varinjlim_{\phi
  : V \to p^{\ast} U} \Mm | V
\mathop{\otimes}_{\Oo_{V}} \phi^{\ast} \phi_{\ast} Q | V {\quoi}.
$$
D'après\pageoriginale \hyperref[V.lem8.2.8]{8.2.8}, les foncteurs $Q \mapsto \Mm | V
\displaystyle\mathop{\otimes}_{\Oo_{V}} \phi^{\ast}
\phi_{\ast} Q | V$ sont exacts à gauche, de sorte
que $\Mm \displaystyle\mathop{\otimes}_{\Oo_T} Q$ est exact à
gauche en $Q$, donc exact.

\begin{enonce*}{Corollaire 8.2.9}\etiquette[8.2.9]{V.cor8.2.9}
Soit $f : (T, \Oo_T) \to (S, \Oo_S)$ un morphisme de
topos annelés. Alors, l'image réciproque par $f$ d'un faisceau de
modules (à droite) plats est un faisceau de modules plats.
\end{enonce*}

Factorisant $f$ par $(T, f^{\ast} \Oo_S)$, on se ramène au cas
$\Oo_T = f^{\ast} \Oo_S$. Soit alors $\Mm$ un
faisceau de modules plats sur $S$. D'après \hyperref[V.prop8.2.7]{8.2.7}, pour vérifier que
$f^{\ast} \Mm$ est plat, il suffit de vérifier l'exactitude du
foncteur
$$
\Nn \mapsto f^{\ast} \Mm \mathop{\otimes}_{\Oo_T}
f^{\ast} \Nn = f^{\ast} (\Mm \mathop{\otimes}_{\Oo_S} \Nn)
$$
d'où l'assertion, puisque $f^{\ast}$ est exact.

\begin{enonce*}{Lemme 8.2.10}\etiquette[8.2.10]{V.lem8.2.10}
Soit $(S, \Oo_S)$ un topos annelé, $\Ff$ un faisceau de modules à
gauche localement de présentation finie, $\Gg$ un faisceau de
bimodules et $\Hh$ un faisceau plat de modules à gauche. Alors, la
flèche canonique
$$
\Hhom (\Ff, \Gg) \otimes \Hh \to \Hhom
(\Ff, \Gg \otimes \Hh)
$$
est un isomorphisme.
\end{enonce*}

La question est locale sur $S$, ce qui permet de supposer que
$\Ff$ admet une présentation finie
$$
\Ff_1 \to \Ff_0 \to \Ff \to 0{\quoi}.
$$
La première ligne  du diagramme suivant est exacte, car $\Hh$
est plat:
\[
\xymatrix{
0 \to \Hhom(\Ff,\Gg) \otimes \Hh \ar[r] \ar[d]&
\Hhom(\Ff_1, \Gg) \otimes \Hh \ar[r]\ar[d]^{S} &
\Hhom(\Ff_2,\Gg) \otimes \Hh \ar[d]^{S}\\
0 \to \Hhom(\Ff,\Gg \otimes \Hh) \ar[r]&
\Hhom(\Ff_1, \Gg \otimes \Hh) \ar[r] &
\Hhom(\Ff_2,\Gg \otimes \Hh)
}
\]
d'où l'assertion. Faisant $\Gg = \Aa$, on déduit de \hyperref[V.lem8.2.10]{8.2.10}:

\begin{enonce*}{Lemme 8.2.11}\etiquette[8.2.11]{V.lem8.2.11}
Tout\pageoriginale morphisme d'un faisceau de modules localement de
présentation finie dans un faisceau de modules plat se factorise,
localement sur $S$, par un faisceau libre $\Oo_S^n (n \in \NN)$.
\end{enonce*}

\begin{enonce*}{Théorème 8.2.12}\etiquette[8.2.12]{V.thm8.2.12}
\emph{(D. Lazard)}. Pour qu'un faisceau de modules $\Mm$ sur un site
  annelé $(S, \Oo_S)$ soit plat, il faut et il suffit qu'il
  soit limite inductive locale de modules libres de type fini.
\end{enonce*}


Le « il suffit » résulte de \hyperref[V.thm8.2.2]{8.2.2}, \hyperref[V.prop8.2.4]{8.2.4}.

Quel que soit $U \in \Ob \Ss$, soit $\Ll_0(U)$ (\resp
$\Ll_1(U)$) la catégorie des faisceaux de modules libres de
type fini (\resp de présentation finie) sur $U$, munis d'une
application linéaire dans $\Mm | U$. Pour toute flèche $f : V \to U$
dans $\Ss$, soit $f^{\ast}$ le foncteur de restriction de $\Ll_i(U)$
dans $\Ll_i(V)$. Les foncteurs définissent une catégorie
$\Ll_i$, fibrée sur $\Ss$, dont les fibres sont les catégories
opposées des catégories $\Ll_i(U)$.

Le lecteur vérifiera que la catégorie $\Ll_1$ est localement
filtrante sur $\Ss$ et que
$$
\Mm = \varinjlim_{(M , f) \in \Ll_1} M
$$

Si $\Mm$ est plat, l'inclusion de $\Ll_0$ dans
$\Ll_1$ est pleinement fidèle et, d'après \hyperref[V.lem8.2.11]{8.2.11}, vérifie $(C
\mathfrak{F} 1)$ de \hyperref[V.def8.1.4]{8.1.4}. Elle vérifie dès lors automatiquement
$(C \mathfrak{F} 2)$, et $\Ll_0$ est localement
filtrante. D'après \hyperref[V.prop8.1.5]{8.1.5}, on a
$$
\Mm = \varinjlim_{(M, f)\in \Ll_2} M{\quoi},
$$
d'où l'assertion \hyperref[V.thm8.2.12]{8.2.12}.

\begin{thebibliography}{99}\pageoriginale
\bibitem{SGA4-05:1} M. Artin et B. \sisi{Mazure}{Mazur} : Etale Homotopy \sisi{Theory}{}, Lecture Notes nᵒ\sisi{}{100}, Springer Verlag.
\bibitem{SGA4-05:2} Blum et Herrera : Article à paraître aux Inventiones.
\bibitem{SGA4-05:3} H. Cartan et S. Eilenberg : Homological Algebra.
\bibitem{SGA4-05:4} P. Deligne : Théorie de Hodge (Publication de l'I.H.E.S).
\bibitem{SGA4-05:5} P. Gabriel et N. Popescu : CRAS.
\bibitem{SGA4-05:6} P. Gabriel et M. Zisman : Homotopie Theory and Calculus of
  Fraction, Ergebnisse der Mathematik, Bd 35.
\bibitem{SGA4-05:7} R. Godement : Théorie des faisceaux. Herman.
\bibitem{SGA4-05:8} J. Giraud : Méthode de la \sisi{Descente}{descente}. Mémoire de la S.M.F.
\bibitem{SGA4-05:9} J. Giraud : Algèbre homologique non commutative (à paraître).
\bibitem{SGA4-05:10} A. Grothendieck : On the De Rham Cohomology of Algebraic
  Varieties, I. H. E. S nᵒ 29.
\bibitem{SGA4-05:11} A. Grothendieck : Sur quelques points d'Algèbre Homologique,
  Tohoku, Math. Journal.
\bibitem{SGA4-05:12} R. Hartshorne : Residues and Duality. Lecture notes
  nᵒ 20, Springer Verlag.
\bibitem{SGA4-05:13} L. Illusie : \SGA 6 I.
\bibitem{SGA4-05:14} S. Lubkin : On a Conjecture of A. Weil. Am. J. of
  Math. p.456, 1967.
\bibitem{SGA4-05:15} G. Segal : Classifying spaces and Spectral Sequences,
  I.H.E.S. nᵒ 34.
\bibitem{SGA4-05:16} Séminaire Cartan 1957-1957.
\bibitem{SGA4-05:17} D.Sullivan : Geometric Topology, part I, Notes
  miméographiées. M.I.T. 1970.
\end{thebibliography}

\ifx\danslelivre\undefined
\end{document}
\fi
